\PassOptionsToPackage{unicode=true}{hyperref} % options for packages loaded elsewhere
\PassOptionsToPackage{hyphens}{url}
%
\documentclass[]{article}
\usepackage{lmodern}
\usepackage{amssymb,amsmath}
\usepackage{ifxetex,ifluatex}
\usepackage{fixltx2e} % provides \textsubscript
\ifnum 0\ifxetex 1\fi\ifluatex 1\fi=0 % if pdftex
  \usepackage[T1]{fontenc}
  \usepackage[utf8]{inputenc}
  \usepackage{textcomp} % provides euro and other symbols
\else % if luatex or xelatex
  \usepackage{unicode-math}
  \defaultfontfeatures{Ligatures=TeX,Scale=MatchLowercase}
\fi
% use upquote if available, for straight quotes in verbatim environments
\IfFileExists{upquote.sty}{\usepackage{upquote}}{}
% use microtype if available
\IfFileExists{microtype.sty}{%
\usepackage[]{microtype}
\UseMicrotypeSet[protrusion]{basicmath} % disable protrusion for tt fonts
}{}
\IfFileExists{parskip.sty}{%
\usepackage{parskip}
}{% else
\setlength{\parindent}{0pt}
\setlength{\parskip}{6pt plus 2pt minus 1pt}
}
\usepackage{hyperref}
\hypersetup{
            pdfborder={0 0 0},
            breaklinks=true}
\urlstyle{same}  % don't use monospace font for urls
\usepackage{color}
\usepackage{fancyvrb}
\newcommand{\VerbBar}{|}
\newcommand{\VERB}{\Verb[commandchars=\\\{\}]}
\DefineVerbatimEnvironment{Highlighting}{Verbatim}{commandchars=\\\{\}}
% Add ',fontsize=\small' for more characters per line
\newenvironment{Shaded}{}{}
\newcommand{\AlertTok}[1]{\textcolor[rgb]{1.00,0.00,0.00}{\textbf{#1}}}
\newcommand{\AnnotationTok}[1]{\textcolor[rgb]{0.38,0.63,0.69}{\textbf{\textit{#1}}}}
\newcommand{\AttributeTok}[1]{\textcolor[rgb]{0.49,0.56,0.16}{#1}}
\newcommand{\BaseNTok}[1]{\textcolor[rgb]{0.25,0.63,0.44}{#1}}
\newcommand{\BuiltInTok}[1]{#1}
\newcommand{\CharTok}[1]{\textcolor[rgb]{0.25,0.44,0.63}{#1}}
\newcommand{\CommentTok}[1]{\textcolor[rgb]{0.38,0.63,0.69}{\textit{#1}}}
\newcommand{\CommentVarTok}[1]{\textcolor[rgb]{0.38,0.63,0.69}{\textbf{\textit{#1}}}}
\newcommand{\ConstantTok}[1]{\textcolor[rgb]{0.53,0.00,0.00}{#1}}
\newcommand{\ControlFlowTok}[1]{\textcolor[rgb]{0.00,0.44,0.13}{\textbf{#1}}}
\newcommand{\DataTypeTok}[1]{\textcolor[rgb]{0.56,0.13,0.00}{#1}}
\newcommand{\DecValTok}[1]{\textcolor[rgb]{0.25,0.63,0.44}{#1}}
\newcommand{\DocumentationTok}[1]{\textcolor[rgb]{0.73,0.13,0.13}{\textit{#1}}}
\newcommand{\ErrorTok}[1]{\textcolor[rgb]{1.00,0.00,0.00}{\textbf{#1}}}
\newcommand{\ExtensionTok}[1]{#1}
\newcommand{\FloatTok}[1]{\textcolor[rgb]{0.25,0.63,0.44}{#1}}
\newcommand{\FunctionTok}[1]{\textcolor[rgb]{0.02,0.16,0.49}{#1}}
\newcommand{\ImportTok}[1]{#1}
\newcommand{\InformationTok}[1]{\textcolor[rgb]{0.38,0.63,0.69}{\textbf{\textit{#1}}}}
\newcommand{\KeywordTok}[1]{\textcolor[rgb]{0.00,0.44,0.13}{\textbf{#1}}}
\newcommand{\NormalTok}[1]{#1}
\newcommand{\OperatorTok}[1]{\textcolor[rgb]{0.40,0.40,0.40}{#1}}
\newcommand{\OtherTok}[1]{\textcolor[rgb]{0.00,0.44,0.13}{#1}}
\newcommand{\PreprocessorTok}[1]{\textcolor[rgb]{0.74,0.48,0.00}{#1}}
\newcommand{\RegionMarkerTok}[1]{#1}
\newcommand{\SpecialCharTok}[1]{\textcolor[rgb]{0.25,0.44,0.63}{#1}}
\newcommand{\SpecialStringTok}[1]{\textcolor[rgb]{0.73,0.40,0.53}{#1}}
\newcommand{\StringTok}[1]{\textcolor[rgb]{0.25,0.44,0.63}{#1}}
\newcommand{\VariableTok}[1]{\textcolor[rgb]{0.10,0.09,0.49}{#1}}
\newcommand{\VerbatimStringTok}[1]{\textcolor[rgb]{0.25,0.44,0.63}{#1}}
\newcommand{\WarningTok}[1]{\textcolor[rgb]{0.38,0.63,0.69}{\textbf{\textit{#1}}}}
\usepackage{longtable,booktabs}
% Fix footnotes in tables (requires footnote package)
\IfFileExists{footnote.sty}{\usepackage{footnote}\makesavenoteenv{longtable}}{}
\usepackage{graphicx,grffile}
\makeatletter
\def\maxwidth{\ifdim\Gin@nat@width>\linewidth\linewidth\else\Gin@nat@width\fi}
\def\maxheight{\ifdim\Gin@nat@height>\textheight\textheight\else\Gin@nat@height\fi}
\makeatother
% Scale images if necessary, so that they will not overflow the page
% margins by default, and it is still possible to overwrite the defaults
% using explicit options in \includegraphics[width, height, ...]{}
\setkeys{Gin}{width=\maxwidth,height=\maxheight,keepaspectratio}
\setlength{\emergencystretch}{3em}  % prevent overfull lines
\providecommand{\tightlist}{%
  \setlength{\itemsep}{0pt}\setlength{\parskip}{0pt}}
\setcounter{secnumdepth}{0}
% Redefines (sub)paragraphs to behave more like sections
\ifx\paragraph\undefined\else
\let\oldparagraph\paragraph
\renewcommand{\paragraph}[1]{\oldparagraph{#1}\mbox{}}
\fi
\ifx\subparagraph\undefined\else
\let\oldsubparagraph\subparagraph
\renewcommand{\subparagraph}[1]{\oldsubparagraph{#1}\mbox{}}
\fi

% set default figure placement to htbp
\makeatletter
\def\fps@figure{htbp}
\makeatother


\date{}

\begin{document}

% pandoc-tablenos: caption name
\renewcommand{\tablename}{Tabela}

% pandoc-fignos: caption name
\renewcommand{\figurename}{Figura}

% pandoc-xnos: cleveref fakery
\newcommand{\plusnamesingular}{}
\newcommand{\starnamesingular}{}
\newcommand{\xrefname}[1]{\protect\renewcommand{\plusnamesingular}{#1}}
\newcommand{\Xrefname}[1]{\protect\renewcommand{\starnamesingular}{#1}}
\providecommand{\cref}{\plusnamesingular~\ref}
\providecommand{\Cref}{\starnamesingular~\ref}
\providecommand{\crefformat}[2]{}
\providecommand{\Crefformat}[2]{}

% pandoc-xnos: cleveref formatting
\crefformat{figure}{figura~#2#1#3}
\Crefformat{figure}{Figure~#2#1#3}
\crefformat{table}{tabela~#2#1#3}
\Crefformat{table}{Table~#2#1#3}
\crefformat{equation}{equação~#2#1#3}
\Crefformat{equation}{Equation~#2#1#3}

\hypertarget{introduuxe7uxe3o}{%
\section{Introdução}\label{introduuxe7uxe3o}}

A área denominada Inteligência Artificial é um dos campos da ciência
mais explorados na atualidade, estando presente em uma gama crescente de
aplicações que vão desde o entretenimento (Risi e Togelius, 2017) à
saúde (Esteva \emph{et al.}, 2017), passando por campos como segurança
digital (Aljumah e Ahamad, 2016; Omar, Johari e Smith, 2017), entre
inúmeros outros.

Sundar Pichai, CEO do Google desde 2015, afirmou recentemente em um
discurso no Fórum Econômico Mundial de 2018 que a inteligência
artificial é provavelmente a área mais importante na qual a humanidade
já trabalhou, equiparando sua importância à eletricidade e ao fogo e
declarando ainda que esta tecnologia é o coração do Google (Parker,
2018). Dave Coplin, \emph{Chief Envisioning Officer} da Microsoft,
afirmou no \emph{AI Summit} de 2016, que a inteligência artificial é a
tecnologia mais importante em que todos estão trabalhando no momento
(Shead, 2016). Um das páginas (``What is Artificial Intelligence? -
Amazon Web Services'', {[}s.d.{]}) do \emph{website} da Amazon é
dedicada à explicar o que é inteligência artificial, e deixa clara a
importância do aprendizado de máquinas (\emph{machine learning}), uma de
suas subáreas, em algumas das competências essenciais da companhia.

Isso ajuda a entender a relevância do campo no modelo de negócios de
algumas das maiores empresas do mundo, já que esse interesse é
acompanhado por outras gigantes da tecnologia, como Apple e Facebook. Na
contramão desse movimento, no entanto, observa-se uma escassa penetração
dessas tecnologias na indústria de instrumentos virtuais, acompanhado de
um baixo volume de pesquisas ne área; Um levantamento bibliográfico na
base Scopus revela que trabalhos relacionados à utilização de redes
neurais para a simulação de instrumentos musicais ainda são bastante
escassos, a despeito do sucesso desta ferramenta em iniciativas afins.

O potencial de inovação nesse campo é sugerido pelo papel central
(muitas vezes representando o estado da arte) que as redes neurais
artificiais vêm desempenhando em áreas diretamente correlatas, com o
campo de síntese de voz, ou ainda em campos menos obviamente
relacionados, à exemplo dos grandes avanços na área de visão de
computadores (\emph{computer vision}), como a geração de imagens e
vídeos.

Esses resultados estimulam a transposição de algumas dessas inovações
para o contexto da síntese sonora, além do desenvolvimento de novas
abordagens. Nesse contexto, o presente trabalho investiga, a partir de
diferentes perspectivas, o potencial de aplicação dos recentes
desenvolvimentos na área de redes neurais à modelagem de instrumentos
acústicos, com vistas à síntese sonora em tempo real.

Ênfase é dada a instrumentos de caráter percussivo, no sentido de
sistemas em que o som é gerado por uma excitação inicial,
aproximadamente impulsiva do ponto de vista físico, e a subsequente
vibração livre do(s) componentes pertinentes do instrumento. Esse é o
caso, por exemplo, do conjunto de peças de uma bateria ao serem
golpeados, das cordas ou conjuntos de cordas de um piano acionadas pelo
martelo ao ser alavancado pelo pressionar das teclas e, em uma
aproximação mais livre, a excitação causada pelos dedos ou palheta em
instrumentos de corda.

Do ponto de vista da indústria, embora o interesse por instrumentos
musicais digitais tenha crescido bastante na última década (Staudt,
2016), os instrumentos virtuais de ponta, utilizados em estúdios para
gravações profissionais ainda baseiam-se prioritariamente em coleções de
amostras sonoras pré-gravadas (Smith, 2008), demandando uma alta
quantidade de memória em disco e razoável poder de processamento do
\emph{hardware} (em geral computadores) utilizado.

Em plataformas onde o poder de processamento e armazenamento são
limitados, como é o caso de teclados digitais e baterias eletrônicas,
por exemplo, é comum que esses instrumentos sejam guarnecidos com
bibliotecas de menor tamanho e qualidade, suficientes apenas para a
prática do instrumento; a qualidade de produção é alcançada apenas
conectando-se o instrumento à um computador, e usando-o como controlador
de um VST de forma a acessar técnicas de manipulação e coleções de
amostras sonoras de melhor qualidade.

O foco em algoritmos neurais voltados à síntese em tempo real tem origem
tanto na escassez de pesquisas sobre o tema quanto nas possibilidades
mercadológicas: Aumentos de eficiência e redução do tamanho em memória
ocupado pelo modelo podem, por exemplo, basear implementações mais
verossímeis e baratas nessas plataformas dedicadas.

\hypertarget{objetivo}{%
\subsection{Objetivo}\label{objetivo}}

O objetivo geral deste trabalho é desenvolver um modelo de simulação em
tempo real de instrumentos acústicos que tire proveito do estado da arte
relacionado à redes neurais artificias para exibir maior eficiência,
quando comparado aos algoritmos tradicionais, dando origem a simulações
mais realistas, principalmente do ponto de vista da percepção humana, e
menos intensivas computacionalmente.

Esse objetivo pode ser desmembrado em objetivos específicos:

Identificar formas compactas de representação das ondas sonoras mais
adequadas à predição/generalização via redes neurais;

Identificar arquiteturas neurais que possam ser utilizadas na modelagem
acústica;

Identificar na literatura sobre modelos acústicos de instrumentos
musicais ferramentas que possam ser utilizadas para simplificação e
aumento de eficiência, tanto do modelo quanto das representações;

Tendo em vista o foco na síntese em tempo real, identificar as
arquiteturas e hiperparametros neurais e técnicas em geral mais
eficientes,

Delinear a capacidade de generalização do modelo criado.

\hypertarget{escopo}{%
\subsection{Escopo}\label{escopo}}

O trabalho limita-se à emulação de instrumentos musicais acústicos
convencionais e não aborda, senão tangencialmente, a área mais subjetiva
de desenho de novos instrumentos, através de um exemplo em que um
instrumento híbrido é apresentado; um tratamento razoavelmente completo
do tema, pelo grande número de possibilidades que oferece (Dalgleish,
Spencer e Foster, 2014), está além do objetivo deste trabalho.

Adicionalmente, limitar o escopo do trabalho à investigação de
instrumentos que possam ser aproximados por um modelo de excitação
impulsivo, como aqui é feito, possibilita excluir da observação uma
subárea significativa do processamento de sinais que lida com a evolução
das frequências de uma onda sonora no domínio do tempo, trazendo o
volume de pesquisa a níveis práticos, sem uma perda substancial em
termos conceituais.

Uma outra vantagem desse recorte é de cunho técnico, já que essa escolha
reduz a duração média das ondas investigadas; investigações do
comportamento de arquiteturas neurais densas e profundas, por exemplo,
tornam-se computacionalmente custosas quando uma dimensionalidade muito
grande é envolvida. Esse também é o comportamento de alguns dos
algoritmos utilizados durante o trabalho, como a Transformada Discreta
de Fourier.

Em relação ao tratamento das redes neurais, prioridade foi dada às
arquiteturas já consolidadas na literatura, o que excluiu, por exemplo,
uma investigação da promissora arquitetura Capsnet (Hinton, Sabour e
Frosst, 2018; Sabour, Frosst e Hinton, 2017). A escolha é necessária
devido à limitações de tempo frente ao grande volume de pesquisas
envolvendo novas arquiteturas.

\hypertarget{estrutura-do-trabalho-e-suxedntese-das-etapas-de-pesquisa}{%
\subsection{Estrutura do Trabalho e Síntese das Etapas de
Pesquisa}\label{estrutura-do-trabalho-e-suxedntese-das-etapas-de-pesquisa}}

Durante a revisão bibliográfica, após um breve resumo da evolução das
redes neurais do ponto de vista histórico, o trabalho apresenta as
arquiteturas mais relevantes encontradas na literatura. O objetivo é
introduzir o tema, inicialmente de maneira ampla, lançando as bases para
as seções seguintes, que exploram o estado da arte das aplicações dessas
arquiteturas à 3 áreas de interesse: Música, discurso e imagens.
Enquanto o interesse nas duas primeiras áreas é mais evidente, uma
investigação dos desenvolvimentos em visão de computadores oferece
insights importantes; além de ser uma das áreas do aprendizado de
máquinas mais ativamente pesquisadas, alguns esforços de transposição
dos desenvolvimentos nessa área para o campo da modelagem acústica tem
alcançado resultados interessantes.

A seguir, algumas plataformas disponíveis são discutidas e comparadas;
novamente, por limitações de tempo, ênfase foi dada aos com maior
penetração tanto na indústria quanto na academia, e a escolha da
plataforma utilizada nesse trabalho obedeceu mais à critérios de
documentação, adoção na comunidade e flexibilidade, do que critérios
relacionados diretamente à performance.

A seção seguinte investiga a modelagem acústica convencional, com ênfase
no estado da arte da síntese sonora em tempo real. Por tratar-se de uma
área mais hermética, sobretudo quando comparada ao campo de redes
neurais e sua filosofia código aberto, alguns temas relevantes são
apresentados com uma maior profundidade conceitual. Por esse mesmo
motivo, implementações didáticas dos dois algoritmos mais utilizados são
apresentadas.

Em seguida é apresentado um referencial teórico, onde são formalizadas
algumas das bases conceituais do trabalho, como a equação da onda e o
funcionamento, do ponto de vista matemático, de uma rede neural;
ademais, uma interpretação geométrica para a simetria da transformada
discreta de Fourier, quando aplicada a sinais no domínio dos números
reais, é introduzida.

A seguir, no capítulo dedicado à metodologia, são apresentadas as etapas
executadas, desde o processo de obtenção das amostras de som utilizadas
na etapa de treinamento das redes, passando pelas investigações no
domínio do tempo e da frequência, até a apresentação em mais detalhes do
modelo final proposto, que tem seus resultados comentados na seção
seguinte. Na conclusão algumas limitações e possibilidades de
encaminhamentos futuros são apresentados.

\hypertarget{relevuxe2ncia-socio-econuxf4mica}{%
\subsection{Relevância
socio-econômica}\label{relevuxe2ncia-socio-econuxf4mica}}

Os computadores introduziram a produção musical em uma nova era,
permitindo, por um lado, experimentações com novos timbres e formas de
interação homem máquina para a criação musical e, por outro, a emulação
de instrumentos tradicionais, através de técnicas promissoras, como a
modelagem acústica (Bovermann \emph{et al.}, 2016).

Avançar essas pesquisas significa, na medida em que, por exemplo,
diminui custos de \emph{software} e \emph{hardware} dedicados, oferecer
a um maior número de pessoas acesso à educação musical e, por
conseguinte, uma formação intelectual mais rica.

Algumas vantagens do acesso precoce à educação musical são delineadas no
estudo de Forgeard \emph{et al.} (2008), que sugere uma maior habilidade
verbal e uma maior capacidade de raciocínio não verbal em crianças que
praticaram instrumentos em sua infância. Na mesma linha, o trabalho de
Vaughn (2000) aponta para uma relação entre o estudo voluntário de
música e uma melhoria do desempenho matemático, a partir de uma meta
análise de 20 estudos.

Uma outra motivação é estimular a descentralização do capital
intelectual relacionado ao campo de instrumentos virtuais, desde a
década de 70 fortemente concentrado em algumas empresas japonesas,
situadas em sua maioria na cidade de Hamamatsu, como Yamaha e Roland
(Reiffenstein, 2006). Essas empresas detém, exclusivamente ou
parcialmente, as patentes para alguns dos mais importantes algoritmos de
síntese sonora, como os \emph{digital waveguides}.

\hypertarget{revisuxe3o-bibliogruxe1fica}{%
\section{Revisão Bibliográfica}\label{revisuxe3o-bibliogruxe1fica}}

Após uma breve exposição sobre o nascimento e evolução das redes
neurais, com o objetivo principal de conferir familiaridade à área onde
o trabalho é desenvolvido, a revisão bibliográfica aqui apresentada
assume um caráter pragmático, investigando na literatura, em especial no
que tange as pesquisas relacionadas à área de redes neurais, técnicas
pertinentes ao desenvolvimento do modelo proposto.

As última seção elenca de maneira bastante concisa o conteúdo
relacionado à modelagem acústica ``tradicional'', limitando-se aos
modelos fisicamente informados mais importantes para o trabalho em tela.
A área de síntese digital, que tem aproximadamente a mesma idade da área
de redes neurais é, talvez, mais extensa e errática do que esta em sua
evolução, e o leitor interessado em uma revisão histórica mais
abrangente é referido ao trabalho de Bilbao (2009), que delineia o
desenvolvimento do campo desde sua origem em algoritmos abstratos.

\hypertarget{redes-neurais-artificiais}{%
\subsection{Redes Neurais Artificiais}\label{redes-neurais-artificiais}}

Multidisciplinar desde o nascimento, o desenvolvimento das redes neurais
artificiais pode ser remontado aos primeiros esforços para
sistematização teórica da forma como o cérebro humano funciona, a partir
dos trabalhos de Hermann von Helmholtz, Ernst Mach e Ivan Pavlov, na
virada do século 19 (Hagan \emph{et al.}, 1996).

Em 1943, o neurofisiologista Warren McCulloch e o matemático Walter
Pitts foram responsáveis por formular o primeiro modelo matemático
conhecido do cérebro humano (McCulloch e Pitts, 1943), mostrando que
topologias simples podem, em princípio, encarregar-se de operações
aritméticas e lógicas (Yadav, Yadav e Kumar, 2015) complexas. O trabalho
desses autores é considerado o início da pesquisa na área de redes
neurais artificiais (Russell e Norvig, 2016).

O termo Inteligência Artificial foi cunhado em 1956 por John McCarthy em
uma conferência no Dartmouth College (Coppin, 2004), da qual
participaram muitos dos primeiros pesquisadores influentes na área
(Brunette, Flemmer e Flemmer, 2009).

Esse termo, assim como o próprio termo inteligência, não se presta a uma
definição fácil; para o propósito deste trabalho, no entanto, pode ser
entendido como um termo guarda-chuva sob o qual algoritmos que buscam
emular aspectos da inteligência humana estão abrigados.

O termo aprendizado de máquinas, por sua vez, é utilizado para denominar
um subgrupo da área de inteligência artificial, do qual fazem parte as
redes neurais, que ocupa-se de técnicas que permitam que máquinas
resolvam problemas sem que sejam explicitamente programadas para tal
(Raschka, 2015).

Na área mais ampla denominada de inteligência artificial, redes neurais
artificiais possuem um papel primordial pois tem sido, em suas diversas
encarnações e arquiteturas, o principal representante do aprendizado de
máquinas em diversas aplicações comerciais (Goodfellow \emph{et al.},
2016).

O primeiro passo para a utilização prática das redes neurais foi dado no
final da década de 1950 por Frank Rosenblatt (Hagan \emph{et al.},
1996), em sua proposta ao laboratório aeronáutico de Cornell de um
autômato baseado em seu modelo simplificado de neurônio, o Perceptron
(Rosenblatt, 1957).

Tomando como base também as teorias de McCulloch, Widrow desenvolveu a
ADALINE (Adaptive Linear Neuron) (Widrow e Hoff, 1960) bastante similar
em estrutura ao modelo de Rosenblatt, e partilhando de suas limitações.
O modelo de treinamento proposto, no entanto, era consideravelmente mais
robusto.

A motivação final de Rosenblatt, inspirada nas teorias sobre o
funcionamento do cérebro humano apresentadas por autores como
Culbertson, Von Neumann e Ashby (Rosenblatt, 1958) era a construção de
uma máquina capaz de aprender à responder diretamente à estímulos
físicos externos, como sinais luminosos.

Para tanto, tal máquina utilizaria como unidade fundamental o
Perceptron, que funciona matematicamente como uma função aplicada à soma
ponderada das entradas e dos vieses dando origem a um classificador
linear, capaz de atualizar seus pesos para aprender, através de exemplos
de entradas e as correspondentes saídas desejáveis (alvos), a separar
linearmente classes diferentes.

A despeito das limitações práticas do modelo proposto, como expostas de
maneira um tanto pessimista por Minsky (Minski e Papert, 1969), é
interessante notar que muitas das características contemporâneas já
estavam presentes nos trabalho seminais de Rosenblatt e Widrow,
notadamente o caráter estocástico das redes neurais, a necessidade de
treiná-las com base em grandes conjuntos de dados e a característica de
caixa preta do modelo treinado.

É interessante observar que Widrow introduziu, como forma de
treinamento, uma caso restrito do algoritmo de descida em gradiente, que
veio a ser amplamente responsável pelo ressurgimento do interesse em
redes neurais décadas mais tarde, em conjunto com uma técnica mais
eficiente de atualização dos pesos e vieses de uma rede.

O modelo de Rosenblat, tendo tido sua primeira implementação na forma de
uma simulação em um computador IBM 704 (Bishop, 2006), ao contrário da
intenção inicial de utilizar hardware específico, inaugura também a
prática da simulação neural.

Enquanto uma arquitetura baseada em uma única camada de Perceptrons
apresenta severas limitações, o agrupamento sucessivo dessas camadas dá
origem a uma topologia, conhecida como Multilayer Perceptron, capaz de
atuar como um aproximador universal (Hornik, 1991). Desde que utilize
uma função de ativação não-linear (Leshno \emph{et al.}, 1993), e dado
um número suficiente de neuronios na camada oculta, essa topologia é
capaz de mapear qualquer conjunto de números finitos a qualquer outro
com precisão arbitrária (Hornik, Stinchcombe e White, 1989).

Um outro impedimento práticos do trabalhos de Rosenblatt e Widrow foi a
ausência de uma metodologia eficiente para a atualização dos pesos da
rede, sobretudo envolvendo múltiplas camadas de neurônios - caso não
coberto pelos algoritmos iniciais tanto de Rosenblat quanto de Widrow;
tal metodologia veio a ser proposta originalmente por Werbos em 1974
(Werbos, 1974). Contudo, essa técnica permaneceu pouco conhecida na
comunidade até ser redescoberta por Parker (Parker, 1985) e, pouco tempo
depois, Rumelhart (Rumelhart, Hinton e Williams, 1985) na segunda metade
da década de 1980 (Mizutani, Dreyfus e Nishio, 2000; Widrow e Lehr,
1990).

O algoritmo, que é conhecido como \emph{backpropagation}, foi um dos
responsáveis por reaquecer o interesse no campo (Hagan \emph{et al.},
1996) na época, inaugurando sua fase atual, com o surgimento dos
principais congressos sobre o assunto, como o \emph{IEEE International
Conference on Neural Networks} e periódicos, a exemplo do \emph{INNS
Neural Networks}, ao fim da década de 1980 (Yadav, Yadav e Kumar, 2015).

Nos anos seguintes, observou-se uma profusão de novas arquiteturas, que
foram aprofundando-se na media em que a utilização de um maior número de
camadas ocultas foi possibilitada pelos avanços no \emph{hardware}
computacional. Outras formas de organizar camadas sucessivas foram
também introduzidas, além de vários avanços incrementais nos algoritmos
de treinamento.

\hypertarget{redes-neurais---principais-arquiteturas}{%
\subsection{Redes Neurais - Principais
Arquiteturas}\label{redes-neurais---principais-arquiteturas}}

Esta seção apresenta as principais arquiteturas consolidadas na
literatura, dando ênfase às arquiteturas relevantes ao trabalho aqui
desenvolvido. Para uma apresentação mais abrangente, ainda que não
exaustiva, o leitor é referido à página preparada pelo The Asimov
Institute (Veen,
2018)(http://www.asimovinstitute.org/neural-network-zoo/), onde uma gama
maior de arquiteturas é elencada, junto com referência aos artigos onde
foram introduzidas.

\hypertarget{redes-densas}{%
\subsubsection{Redes Densas}\label{redes-densas}}

A arquitetura básica no campo das redes neurais artificiais é a chamada
densa (\emph{Feed-Forward}), que consiste de várias camadas sucessivas
totalmente ligadas entre si por meio de pesos. Os impulsos recebidos
pelas camadas mais baixas fluem sucessivamente para as camadas
posteriores, como sugere a nomenclatura em inglês desta arquitetura. São
uma generalização do Multilayer Perceptron na medida em que utilizam em
geral uma gama mais vasta de funções de ativação, muitas delas de forma
sigmoidal, como a função logística \(y = \frac{1}{1 + e ^ {-x}}\) e a
tangente hiperbólica \(\tanh(x) = \frac{e^x − e^{-x}}{e^x + e^{-x}}\)
(Goldberg, 2016)

\hypertarget{redes-neurais-profundas}{%
\subsubsection{Redes Neurais Profundas}\label{redes-neurais-profundas}}

Com a melhoria do hardware, a topologia densa foi recebendo um acréscimo
de camadas, gerando as redes profundas (Deep Neural Networks). O aspecto
mais importante dessa movimentação foi que as redes ganharam a
habilidade de gerar representações sucessivas, abstraindo diferentes
aspectos dos dados em cada uma de suas camadas.

A grande vantagem desta topologia foi sua capacidade automática de
extração dos componentes descritivos dos dados, um trabalho que ficava a
cargo dos pesquisadores anteriormente (Socher, 2014). Em contrapartida,
a forma como os dados são interpretados pela rede pode não ser
facilmente inferida pelo pesquisador, e as redes podem assumir uma forte
característica de caixa-preta.

\begin{figure}
\centering
\includegraphics{im/arqdense.png}
\caption{Representação de uma Rede Densa e uma Rede Densa
Profunda\label{fig:dense}}
\end{figure}

Fonte: Elaboração própria

A \xrefname{figura}\cref{fig:dense} observa-se uma representação
esquemática de duas arquiteturas densas: uma rede ``rasa'', à esquerda,
e uma arquitetura profunda, à direita. Os círculos brancos representam
as entradas da rede, enquanto os pretos representam suas saídas. Em
cinza, temos a representação dos neurônios que compõe as camadas ocultas
da rede.

\hypertarget{redes-recorrentes}{%
\subsubsection{Redes Recorrentes}\label{redes-recorrentes}}

As redes recorrentes são representações compactas, do ponto de vista
técnico, das arquiteturas profundas (Schmidhuber, 2015), devido ao seu
caráter recursivo. Conceitualmente, contudo, a abordagem é distinta:
Nestas redes os neurons são parcialmente alimentados com seus próprios
estados anteriores, emulando um efeito similar à utilização de ligações
entre neurons de uma camada anterior com neurons de uma camada posterior
não adjacente (Veit, Wilber e Belongie, 2016).

Essa arquitetura foi proposta por Elman (Elman, 1990) com o propósito de
capturar informações codificadas no encadeamento temporal de séries de
dados, e é bastante poderosa em várias aplicações, como modelos de
previsão e classificação de informações (Xu, Auli e Clark, 2015), por
exemplo.

\begin{figure}
\centering
\includegraphics{im/arqrecurrent.png}
\caption{Representação da evolução temporal de uma Rede
Recorrente\label{fig:recur}}
\end{figure}

Fonte: Elaboração própria

A \xrefname{figura}\cref{fig:recur} representa 6 passos iniciais de uma
rede recorrente. Observa-se que em cada passo, a partir do primeiro, a
rede recebe, além do vetor de entrada externo, as próprias saídas que
gerou no passo anterior. É oportuno notar que, em algumas aplicações, a
rede recebe apenas um vetor de entradas, no primeiro passo, e opera em
todos os passos seguintes processando suas entradas anteriores, em uma
configuração que recebe o nome de um para muitos (\emph{one-to-many}) na
literatura.

\hypertarget{redes-convolucionais}{%
\subsubsection{Redes Convolucionais}\label{redes-convolucionais}}

Trata-se de um tipo de arquitetura, geralmente profunda, que é
amplamente utilizado em problemas relacionados a imagens, atingindo
resultados de ponta em várias áreas relacionadas à visão de
computadores, como reconhecimento de objetos e rostos em imagens (Pang
\emph{et al.}, 2017).

o problema da alta dimensionalidade dos vetores de entrada, como é o
caso da representação de uma imagem por exemplo, é amenizado através da
substituição de camadas totalmente conectadas por camadas
convolucionais, que varrem a imagem, movimentando-se em cada uma de suas
dimensões um passo por vez cada vez (LeCun \emph{et al.}, 1998), e
atualizando todos os pesos da camada convolucional de acordo.

Esse procedimento permite a geração, nessas camadas convolucionais da
rede, de uma só representação para padrões que aparecem em diferentes
pontos do vetor de entrada; tais representações são geralmente
interpretadas em camadas posteriores, totalmente conectadas, de forma a
gerar o resultado final.

\begin{figure}
\centering
\includegraphics{im/arqconv.png}
\caption{Representação da evolução temporal de uma Rede
Convolucional\label{fig:arqconv}}
\end{figure}

Fonte: Elaboração própria

Na \xrefname{figura}\cref{fig:arqconv} é possível acompanhar as 3
primeiras entradas de uma rede convolucional unidimensional: o vetor de
entrada é ``varrido'' pela rede, geralmente andando um passo à cada
leitura, em um processo análogo à definição matemática de convolução. O
trabalho de Dumoulin e Visin (2016) apresenta uma exposição de vários
tipos de convolução, dentre as quais a convolução básica em duas
dimensões, apresentada na \xrefname{figura}\cref{fig:conv}

\begin{figure}
\centering
\includegraphics{im/conv.gif}
\caption{Convolução Bidimensional\label{fig:conv}}
\end{figure}

Fonte: Dumoulin e Visin (2016)

\hypertarget{redes-neurais-aplicadas-uxe0-musica}{%
\subsection{Redes Neurais Aplicadas à
Musica}\label{redes-neurais-aplicadas-uxe0-musica}}

A maioria dos trabalhos que investigam a aplicação de redes neurais à
música, sobretudo envolvendo sua síntese, ocorrem em um nível de
abstração mais alto do que a geração direta dos sons: Geralmente tomam
como base a manipulação de representações musicais como partituras, por
exemplo, ou representações sonoras compactas, como espectrogramas. Os
motivos passam pela alta dimensionalidade dos outputs: no caso de um
\emph{framerate} de 44100 samples por segundo, a qualidade encontrada em
cds, por exemplo, a síntese de 10 segundos de áudio envolve a criação de
mais de 4 milhões de samples.

O trabalho desenvolvido em conjunto pelas equipes do Google Brain team e
do DeepMind é um exemplo de esforço nesse sentido (Engel \emph{et al.},
2017). Nele, uma arquitetura desenvolvida a partir da Wavenet
(aprofundada no tópico sobre discurso falado) é utilizada para a geração
de ondas sonoras a partir do treinamento direto com amostras de áudio de
vários instrumentos musicais. Os resultados mostram que a arquitetura
baseada em várias camadas convolucionais utilizada é capaz de aprender
representações no domínio do tempo para vários tipos de instrumentos
diferentes.

Uma extensão experimental desse trabalho, denominada de Magenta
(``Magenta'', {[}s.d.{]}), investiga representações latentes para
sequências musicais no domínio do tempo, a partir de uma abordagem
probabilística (Roberts, Engel e Oore \emph{et al.}, 2018; Roberts,
Engel e Raffel \emph{et al.}, 2018).

Uma plataforma física para a manipulação e reprodução dos sons gerados
pela rede, na forma de um \emph{hardware} código aberto denominado
\emph{Nsynth Super} (Eck, 2018; ``NSynthSuper'', {[}s.d.{]}) é
introduzida pela equipe do projeto. Embora não fique claro se a etapa de
síntese ocorre diretamente na plataforma, os esforços corroboram o
interesse aplicação de redes neurais para a síntese sonora em tempo
real.

Enquanto as camadas das arquiteturas anteriores são predominantemente
convolucionais, Pfalz (2018) apresenta um trabalho bastante abrangente
sobre o uso de redes neurais recorrentes puras para essa finalidade. Um
dos primeiros trabalhos na área é, talvez, o desenvolvido por Stanley
(2007), consistindo na investigação da capacidade de uma rede em prever
sons, em analogia direta com séries temporais. Os resultados são
bastante tímidos, e uma rede capaz de prever apenas uma onda sonora,
oriunda de uma nota de saxofone, é apresentada. De maneira semelhante,
Sarroff e Casey (2014) investiga a aplicação direta de redes neurais, de
arquitetura densa, na geração direta de múltiplas amostras de áudio

A geração de música em níveis mais altos de abstração é, para o trabalho
aqui desenvolvido, apenas marginalmente interessante, na medida em que
pode oferecer inspiração sobre representações compactas dos sons. O
trabalho de (Hutchings, 2017), por exemplo, utiliza tablaturas para a
previsão de partes completas de bateria tomando como base o ritmo de uma
de suas peças.

Tomando emprestados os desenvolvimentos obtidos na área da visão de
computadores, os trabalhos de Grinstein \emph{et al.} (2017) e Mital
(2017) buscam transpor os avanços alcançados no campo da transferência
de estilo entre imagens para o domínio do som. O primeiro, de cunho
exploratório, investiga as arquiteturas mais promissoras para a tarefa,
chamando a atenção para o desempenho de arquiteturas rasas, enquanto o
segundo investiga, a partir da Wavenet, representações apropriadas.
Ainda inspirado em técnicas relacionadas à geração de imagens, Donahue,
McAuley e Puckette (2018) faz uso de GANs para a síntese direta de
áudio, no contexto da emulação do discurso falado, gerando resultados
inteligíveis.

Os esforços de classificação costumam introduzir inovações interessantes
nas representações utilizadas:(Costa, Oliveira e Silla, 2017) e (Choi,
Fazekas e Sandler, 2016) abordam a tarefa a partir de redes
convolucionais; no primeiro trabalho as redes são alimentadas com música
representada por um tipo especial de espectrograma enquanto no segundo
arquiteturas convolucionais profundas são investigadas, assim como em
(Choi \emph{et al.}, 2017) que acrescenta recorrência às arquiteturas
convolucionais, no intuito de explorar as inter-relações temporais das
amostras sonoras utilizadas.

Representações não canônicas também possuem importante papel nas tarefas
relacionadas à transcrição musical a partir do uso de redes neurais,
onde busca-se traduzir amostras sonoras em representações de nível mais
alto, simbólicas, como tablaturas e partituras, por exemplo. Uma das
primeiras contribuições nesse sentido pode ser vista em (Tuohy, Potter e
Center, 2006) por meio de um modelo que combina redes neurais e uma
heurística \emph{hill-climber}.

Com o uso de uma rede recorrente, (Boulanger-Lewandowski, Bengio e
Vincent, 2013) transcreve, a partir de espectrogramas, diferentes partes
musicais em comandos MIDI. Na mesma linha (Böck e Schedl, 2012) e
(Sigtia, Benetos e Dixon, 2016) introduzem trabalhos focados na
transcrição de sons polifônicos produzidos por pianos. Em relação á
transcrição de partes de um \emph{kit} de bateria, (Southall, Stables e
Hockman, 2016) aborda a tarefa a partir de uma arquitetura recorrente
alimentada com representações espectrais.

\hypertarget{redes-neurais-aplicadas-ao-discurso-falado}{%
\subsection{Redes Neurais Aplicadas ao Discurso
Falado}\label{redes-neurais-aplicadas-ao-discurso-falado}}

Um dos trabalho mais notáveis nesse sentido, desenvolvido pela equipe de
inteligência artificial do Google, é a chamada Wavenet (Van Den Oord
\emph{et al.}, 2016; ``WaveNet: A Generative Model for Raw Audio'',
{[}s.d.{]}), uma arquitetura baseada em camadas convolucionais, para a
geração ponta a ponta de voz. Um sample é gerado por vez, levando em
conta um grande número de entradas passadas; os resultados introduzem um
novo estado da arte para a tarefa. Os desdobramentos dessa arquitetura
na área musical foram apresentados na seção anterior.

(Hinton \emph{et al.}, 2012) apresenta uma revisão da literatura sobre o
uso de modelagem acústica em abordagens baseadas em redes neurais na
área de reconhecimento de voz. No mesmo ano (Graves, Mohamed e Hinton,
2013) alcança resultados no \emph{TIMIT phoneme recognition benchmark},
com a aplicação de redes recorrentes profundas, que estabelecem o novo
estado da arte, enquanto (Maas, Hannun e Ng, 2013) aponta a
superioridade da função de ativação relu (\emph{rectifier
nonlinearities}) sobre ativações de caráter sigmoidal em tarefas
relacionadas ao reconhecimento contínuo da fala.

A partir de uma arquitetura hibrida, combinando um modelo fonético
recorrente com um classificador acústico baseado em uma rede neural
profunda, (Boulanger-Lewandowski \emph{et al.}, 2014) aplica estratégias
de sequenciamento telefônico para estabelecer um novo benchmark no
\emph{TIMIT dataset}, uma técnica que foi comprovada por (Sak \emph{et
al.}, 2015) como superior ao uso de arquiteturas como \emph{deep long
short-Term memory recurrent neural networks} e as abordagens mais
utilizadas à época baseadas no modelo oculto de Markov.

Com relação à utilização de redes convolucionais, Sainath \emph{et al.}
(2015) investiga a otimização de seus hiperparametros, além de
estratégias de \emph{pooling} e treinamento para aplicações em
reconhecimento de fala, enquanto Zweig \emph{et al.} (2017) e Zhang
\emph{et al.} (2017) exploram o desempenho de sistemas baseados somente
em redes neurais (\emph{end-to-end}), com o último combinando redes
convolucionais hierárquicas com classificação temporal coneccionista.
Nessa mesma linha, o trabalho de (Zhang, Chan e Jaitly, 2017)também
investiga sistemas \emph{end-to-end} a partir de uma arquitetura
combinando redes neurais convolucionais profundas com recorrência e
princípios baseados no aninhamento de redes neurais (NIN, \emph{network
in a network}). Alguns desses princípios, notadamente a justaposição de
convolução e recorrência, inspiraram a implementação apresentada na
seção sobre convolução e recorrência.

\hypertarget{redes-neurais-aplicadas-uxe0-imagem}{%
\subsection{Redes Neurais Aplicadas à
Imagem}\label{redes-neurais-aplicadas-uxe0-imagem}}

Recentemente tem-se visto esforços para adaptar métodos utilizados com
sucesso na área da visão de computadores ao campo da modelagem acústica.
O trabalho de (Engel \emph{et al.}, 2017), batizado de Nsynth, busca
inspiração na área de geração de imagens para elaborar uma arquitetura
apropriada a síntese sonora. Ainda em paralelo com a área de relacionada
à imagens, o autor propõe um base de dados sonora à luz de bases de
dados clássicas no campo de visão de computadores, como o MNIST, em um
esforço para alavancar pesquisas nessa direção.

Da mesma forma a Wavenet, antecessora do Nsynth, é uma adaptação das
arquiteturas PixelRNN (Oord, Kalchbrenner e Kavukcuoglu, 2016) e
PixelCNN(Oord \emph{et al.}, 2016), desenvolvidas para a geração de
imagens.

Os trabalhos que tem como foco a síntese de imagens geralmente envolvem
as chamadas GANs: trata-se de uma abordagem onde uma rede é treinada
para gerar imagens, enquanto outra as classifica em artificiais ou não.
O objetivo da primeira rede é enganar a segunda, dando origem a imagens
mais realistas, do ponto de vista da percepção humana.

Isola \emph{et al.} (2016), por exemplo, utiliza essa técnica para gerar
imagens a partir de representações de seus contornos, enquanto Hwang e
Zhou (2016), Zhang, Isola e Efros (2016), Larsson, Maire e Shakhnarovich
(2016) e Iizuka, Simo-Serra e Ishikawa (2016) abordam a tarefa de
colorização automática de imagens.

Ainda usando redes adversariais (Frans, 2017) introduz controle ao
processo de colorização, com a utilização de mapas de cores, que são
usados como entradas para as redes, um conceito que é explorado também
por Sangkloy \emph{et al.} (2016), onde a interação do usuário e uma
rede densa permitem a colorização de imagens em tempo real, através de
linhas desenhadas em áreas da imagem que possuam as cores pretendidas.

Gatys, Ecker e Bethge (2016) propõe uma técnica de transferência de
estilo, utilizando uma rede convolucional capaz de aprender
características do estilo de uma imagem e transferir para outra, sem
alterar o seu conteúdo semântico. Nessa linha,{[}Kulkarni \emph{et al.}
(2015) utiliza uma arquitetura de convolução-desconvolução para separar
características de posicionamento de uma imagem, permitindo sua
reconstrução posterior em outras posições e situações de iluminação, a
partir de mudanças manuais nas variáveis de entrada da rede. Ainda
investigando a manipulação de imagens a partir de redes adversariais o
trabalho de Zhu \emph{et al.} (2016) tem como objetivo aprender
características das imagens de maneira direta.

O trabalho de Oord, Kalchbrenner e Kavukcuoglu (2016), que veio a
inspirar a Wavenet, propõe um topologia recorrente profunda com conexões
residuais melhoradas capaz de reconstruir imagens parcialmente
obstruídas, enquanto (Theis e Bethge, 2015) investiga arquiteturas LSTM
multi dimensionais no contexto da modelagem da distribuição de imagens.

Na área de vídeo, que pode ser entendida como um caso geral da
manipulação de imagens, com complexidades adicionais relacionadas à alta
dimensionalidade dos dados e ao encadeamento temporal Karpathy \emph{et
al.} (2014) investiga abordagens capazes de estender as redes
convolucionais de forma a permitir que tirem proveito das
características temporais dos dados de entrada. De forma semelhante,
portando desenvolvimentos na área de reconhecimento de expressões
faciais obtidos em imagens, Khorrami \emph{et al.} (2016) combina redes
recorrentes e convolucionais, medindo a importância relativa de cada uma
nos resultados finais.

Buscando combater a alta dimensionalidade, Yang, Krompass e Tresp (2017)
apresenta o conceito de \emph{tensor-train}, permitindo a transferência
de conhecimento de outras arquiteturas para a utilização em dados
sequenciais com alta dimensionalidade. O trabalho de He \emph{et al.}
(2015) introduz novos benchmarks à tarefa de reconhecimento de emoções
em vídeos, utilizando uma arquitetura profunda baseada em camadas LSTM
bidirecionais no processamento de vídeos, incluindo seu áudio.

\hypertarget{redes-neurais---ferramentas-e-frameworks}{%
\subsection{Redes Neurais - Ferramentas e
Frameworks}\label{redes-neurais---ferramentas-e-frameworks}}

A área de aprendizado de máquinas encontra-se muito ativa atualmente,
tanto no âmbito da academia quanto no da indústria, com várias grandes
empresas de tecnologia incorporando essa tecnologia em suas competências
essenciais. Dessa forma, assistimos à uma proliferação de ferramentas e
plataformas focadas em diferentes aspectos da área, muitas delas
desenvolvidos ou endossados por essas empresas, como Microsoft, Google e
Amazon. Somando-se a isso o fato de que, em maior ou menor grau, todos
eles apresentam uma curva de aprendizado significativa, fica evidente
que uma comparação direta de todos, ou mesmo da maioria das ferramentas
disponíveis, encontra-se além do escopo deste trabalho.

Algum direcionamento, contudo, pode ser obtido a partir da literatura
disponível. Bahrampour \emph{et al.} (2016), por exemplo, compara quatro
ferramentas: Caffe (``Caffe'', {[}s.d.{]}), desenvolvida pelo grupo de
pesquisas em inteligência artificial de Berkeley (``Berkeley Artificial
Intelligence Research Lab'', {[}s.d.{]}); Neon (``Neon'', {[}s.d.{]}),
desenvolvido pelo setor de inteligência artificial da Intel; Torch
(``Scientific computing for LuaJIT.'', {[}s.d.{]}) e Theano (``Theano'',
{[}s.d.{]}) os dois últimos sem afiliações diretas com grandes empresas.
O autor conclui pela superioridade da plataforma Theano, em termos de
velocidade em geral e flexibilidade, chamando atenção para sua
capacidade de auto-diferenciação.

Embora o time principal tenha abandonado o desenvolvimento do Theano no
final de 2017 a plataforma Tensorflow (``TensorFlow'', {[}s.d.{]}),
introduzida por pesquisadores do \emph{Google Brain Team} em 2015 é, em
muitos aspectos, sucessora do Theano, conservando a maioria de suas
características.

Se um estudo realizado pelo mesmo autor em 2015 (Bahrampour \emph{et
al.}, 2015) indicava alguma ineficiência da plataforma, à época ainda em
sua infância, o trabalho de Shi \emph{et al.} (2016), ao comparar 5
plataformas, conclui que a situação 1 ano depois é diferente, e que não
há uma plataforma sensivelmente superior às outras em termos gerais.

O trabalho de Erickson \emph{et al.} (2017) é ainda mais abrangente,
comparando 12 \emph{frameworks}, e chamando atenção para a importância
das habilidades prévias e características da pesquisa na escolha da
ferramenta mais apropriada. Na mesma linha, Parvat \emph{et al.} (2017)
compara 5 plataformas, chamando atenção para as potenciais vantagens, em
termos do tamanho do código e rapidez na criação de protótipos, da
utilização de bibliotecas de nível de abstração mais alto, como a
biblioteca Keras (``Keras: The Python Deep Learning library'',
{[}s.d.{]}), no topo das plataformas convencionais.

\hypertarget{suxedntese-sonora-em-tempo-real}{%
\subsection{Síntese Sonora em Tempo
Real}\label{suxedntese-sonora-em-tempo-real}}

Excluídas as técnicas baseadas em manipulação de amostras pré-gravadas
de sons, talvez a abordagem mais utilizada por músicos atualmente
(Bilbao, 2009), podemos dividir área de síntese sonora propriamente em
duas escolas, na medida em que ocupa-se em modelar os processos físicos
que dão origem ao som (Modelagem Física - \emph{Physical Modelling}) ou
diretamente as características das ondas sonoras (Modelagem Espectral -
\emph{Spectral Modelling}) (Serra, 2007).

Dessas escolas, a modelagem física é a mais proeminente tanto do ponto
de vista da pesquisa quanto de aplicações; a área de modelagem espectral
parece ter despertado pouco interesse dos pesquisadores nas últimas duas
décadas, encontrando a maioria de suas aplicações recentes em áreas fora
da síntese musical. A modelagem física apresenta como principais
vantagens um menor número de parâmetros, que costumam ser mais
intuitivos, na medida em que são diretamente relacionados à grandezas
existentes no mundo real.

Abordagens baseadas na modelagem física ocupam-se majoritariamente em
formular modelos discretos baseados em descrições matemáticas dos
instrumentos, como por exemplo a solução de D'alembert para a equação da
onda (Karjalainen e Erkut, 2004), normalmente incorporando ao modelo
termos relacionados por exemplo à rigidez exibida por cordas não ideais,
que modelem características relevantes ao timbre do instrumento
representado (Bensa \emph{et al.}, 2003). Destacam-se, nessa escola,
duas abordagens: o método das diferenças finitas, acurado e
computacionalmente intensivo, e os \emph{Digital Waveguides}, um modelo
bastante eficiente (Bensa \emph{et al.}, 2005) e elegante que é, talvez,
o mais utilizado para a emulação em tempo real de instrumentos
acústicos.

\hypertarget{muxe9todo-das-diferenuxe7as-finitas}{%
\subsubsection{Método das Diferenças
Finitas}\label{muxe9todo-das-diferenuxe7as-finitas}}

A formulação mais básica desse método consiste na resolução da equação
da onda propagando-se em uma corda ideal, introduzida abaixo e derivada
com mais detalhes na seção seguinte.

\begin{equation}\frac{\partial^2 y(x,t)}{\partial t^2} =\frac{T}{\rho} \frac{\partial^2 y(x,t)}{\partial x^2}\label{eq:onda}\end{equation}

A solução é buscada para pontos em uma malha obtida discretizando-se
tanto o tempo \(t\) quanto a posição horizontal \(x\). Isso permite a
aproximação das derivadas parciais de segunda ordem pelo método das
diferenças centrais. Assim procedendo, podemos reescrever a equação da
onda como abaixo.

\[\frac{y[i,j+1]-2y[i,j]+y[i,j-1]}{\Delta t^2} = c^2 \frac{y[i+1,j]-2y[i,j]+y[i-1,j]}{\Delta x^2}\]

onde

\(c=\sqrt{\frac{T}{\rho}}\)

\(x_i = i \Delta x, i \in \{0,1,...,M\}\)

\(t_j = j \Delta t, j \in \{0,1,...,N\}\)

e, colocando \(y[i,j+1]\) em evidência:

\[ y[i,j+1] = 2 y[i, j] - y[i,j-1] + C^2 (y[i+1, j] - 2 y[i, j] + y[i-1, j])\]

Em \(C=c \frac{\Delta t}{\Delta x}\), denominado número de Courant,
foram aglutinados todos os parâmetros que governam a qualidade da
simulação. Para garantir sua estabilidade, os parâmetros devem ser
fixados de modo que \(C\) seja menor que 1.

Repare que \$ y{[}i,j+1{]}\$ é calculado com base no ultimo e penúltimo
momentos discretos, o que demanda, para o primeiro passo da simulação em
\(j = 1\), a definição do estado inicial do sistema em \(j=0\), e
informações sobre o sistema no momento discreto indefinido \(j=-1\).
Essa última dificuldade será atacada a partir de considerações
adicionais sobre o sistema.

Como condições de contorno para o caso contínuo temos que, no momento
inicial, a corda encontra-se deslocada do seu equilíbrio, em repouso.
Portanto, a velocidade em \(t=j=0\), para qualquer ponto \(x\) da corda,
tem valor zero. Ademais, define-se a corda como fixa em suas
extremidades nos pontos \(x=0\) e \(x=L\). Para primeira condição
podemos escrever:

\[ \frac{\partial y(x,t)}{\partial t} \approx \frac{y[i,1]-y[i,-1]}{2 \Delta t} = 0 ~ \forall ~ i \in \{0,1,...,M\} \]

o que implica:

\[ y[i,1] = y[i,-1] ~ \forall ~ i \in \{0,1,...,M\} \]

Assim, para o primeiro momento da simulação em \(j=1\) podemos
substituir \(y[i,-1]\) por \(y[i,1]\), obtendo a equação para quando
\(j=0\):

\[ y[i, 1] = y[i, 0] + \frac{1}{2} C^2 (y[i+1, 0] - 2 y[i, 0] + y[i-1, 0])\]

Em relação à segunda condição, devemos reforças o fato de que
\(y[0, j] = y[L, j] = 0 ~ \forall ~ j \in \{0,1,...,N\}\) em cada passo
da simulação.

É conveniente colocar a formulação em termos físicos: se considerarmos,
por exemplo, um \emph{framerate} típico de 44100 amostras por segundo
podemos evitar desperdício computacional fazendo com que os passos da
simulação coincidam com o intervalo de 1/44100 segundo entre os pontos
amostrados. De modo geral, levando em conta que \(C \le 1\), podemos
parametrizar a simulação em termos dos componentes físicos e o
\emph{framerate} desejado da seguinte forma:

\(dt = \frac{1}{FPS} = \frac{D}{N}\)

\(N = D ~ FPS\)

\(dx = \frac{L}{M}\)

\(M \le \frac{FPS}{2f}\)

\(c = 2fL\)

onde \(D\) é a duração desejada, em segundos; \(L\) é o comprimento da
corda em metros e FPS é o \emph{framerate} pretendido.

O fragmento de implementação no quadro 1, na linguagem Python, faz as
vezes de pseudocódigo e ajuda a ilustrar o funcionamento do algoritmo. A
versão completa do código, com a possibilidade da geração do gráfico da
onda, pode ser encontrada em Tesserato (2018), no caminho
\href{https://github.com/tesserato/Dissertacao/blob/master/Finite_Difference.py}{``resources/Finite\_Difference.py''}.

\begin{Shaded}
\begin{Highlighting}[numbers=left,,]
\NormalTok{amplitude }\OperatorTok{=} \FloatTok{0.005} \CommentTok{# meters}
\NormalTok{pluck_position }\OperatorTok{=} \FloatTok{.5} \CommentTok{# fraction of L}
\NormalTok{pickup_position }\OperatorTok{=} \FloatTok{.5} \CommentTok{# fraction of L}
\NormalTok{fps }\OperatorTok{=} \DecValTok{44100} \CommentTok{# samples / second}
\NormalTok{frequency }\OperatorTok{=} \DecValTok{440} \CommentTok{#hz}
\NormalTok{duration }\OperatorTok{=} \DecValTok{1} \CommentTok{# seconds}
\NormalTok{L }\OperatorTok{=} \FloatTok{0.6} \CommentTok{# meters}
\NormalTok{sustain }\OperatorTok{=} \FloatTok{.9998}

\NormalTok{N }\OperatorTok{=} \BuiltInTok{int}\NormalTok{(duration }\OperatorTok{*}\NormalTok{ fps) }\CommentTok{# number of time points}
\NormalTok{dt }\OperatorTok{=} \DecValTok{1} \OperatorTok{/}\NormalTok{ fps }\CommentTok{# spacing beetwen discrete time points}
\NormalTok{M }\OperatorTok{=} \BuiltInTok{int}\NormalTok{(fps }\OperatorTok{/}\NormalTok{ (}\DecValTok{2} \OperatorTok{*}\NormalTok{ frequency)) }\CommentTok{# number of position points}
\NormalTok{dx }\OperatorTok{=}\NormalTok{ L }\OperatorTok{/}\NormalTok{ M }\CommentTok{# spacing beetwen discrete position points}
\NormalTok{c }\OperatorTok{=}\NormalTok{ frequency }\OperatorTok{*} \DecValTok{2} \OperatorTok{*}\NormalTok{ L }\CommentTok{# meters / second}
\NormalTok{C }\OperatorTok{=}\NormalTok{ c }\OperatorTok{*}\NormalTok{ dt }\OperatorTok{/}\NormalTok{ dx }\CommentTok{# Courant number}
\NormalTok{x }\OperatorTok{=}\NormalTok{ np.arange(}\DecValTok{0}\NormalTok{, M }\OperatorTok{+} \DecValTok{1}\NormalTok{, }\DecValTok{1}\NormalTok{) }\OperatorTok{*}\NormalTok{ dx}
\NormalTok{t }\OperatorTok{=}\NormalTok{ np.arange(}\DecValTok{0}\NormalTok{, N }\OperatorTok{+} \DecValTok{1}\NormalTok{, }\DecValTok{1}\NormalTok{) }\OperatorTok{*}\NormalTok{ dt}
\NormalTok{pickup }\OperatorTok{=} \BuiltInTok{int}\NormalTok{(}\BuiltInTok{round}\NormalTok{(M }\OperatorTok{*}\NormalTok{ pickup_position))}
\NormalTok{asc }\OperatorTok{=} \BuiltInTok{int}\NormalTok{(}\BuiltInTok{round}\NormalTok{(M }\OperatorTok{*}\NormalTok{ pluck_position))}
\NormalTok{dsc }\OperatorTok{=}\NormalTok{ M }\OperatorTok{+} \DecValTok{2} \OperatorTok{-}\NormalTok{ asc}
\NormalTok{X_asc }\OperatorTok{=}\NormalTok{ np.linspace(}\DecValTok{0}\NormalTok{, }\DecValTok{1}\NormalTok{, asc)}
\NormalTok{X_dsc }\OperatorTok{=}\NormalTok{ np.linspace(}\DecValTok{0}\NormalTok{, }\DecValTok{1}\NormalTok{, dsc)[::}\OperatorTok{-}\DecValTok{1}\NormalTok{][}\DecValTok{1}\NormalTok{: ]}
\NormalTok{y }\OperatorTok{=}\NormalTok{ np.zeros(M }\OperatorTok{+} \DecValTok{1}\NormalTok{)}
\NormalTok{y_1 }\OperatorTok{=}\NormalTok{ np.hstack([X_asc,X_dsc]) }\OperatorTok{*}\NormalTok{ amplitude }\OperatorTok{*}\NormalTok{ np.random.normal(}\DecValTok{1}\NormalTok{, }\FloatTok{.01}\NormalTok{, M }\OperatorTok{+} \DecValTok{1}\NormalTok{) }\CommentTok{# initial displacement shape}
\NormalTok{y_2 }\OperatorTok{=}\NormalTok{ np.zeros(M }\OperatorTok{+} \DecValTok{1}\NormalTok{)}

\NormalTok{ctr}\OperatorTok{=}\DecValTok{0}
\NormalTok{w }\OperatorTok{=}\NormalTok{ []}
\ControlFlowTok{for}\NormalTok{ i }\KeywordTok{in} \BuiltInTok{range}\NormalTok{(}\DecValTok{1}\NormalTok{, M):}
\NormalTok{  y[i] }\OperatorTok{=}\NormalTok{ y_1[i] }\OperatorTok{+} \FloatTok{0.5} \OperatorTok{*}\NormalTok{ C}\OperatorTok{**}\DecValTok{2} \OperatorTok{*}\NormalTok{(y_1[i}\OperatorTok{+}\DecValTok{1}\NormalTok{] }\OperatorTok{-} \DecValTok{2}\OperatorTok{*}\NormalTok{y_1[i] }\OperatorTok{+}\NormalTok{ y_1[i}\DecValTok{-1}\NormalTok{])}
\NormalTok{y[}\DecValTok{0}\NormalTok{] }\OperatorTok{=} \DecValTok{0}
\NormalTok{y[M] }\OperatorTok{=} \DecValTok{0}
\NormalTok{w.append(y[pickup])}
\NormalTok{y_2[:] }\OperatorTok{=}\NormalTok{ y_1}
\NormalTok{y_1[:] }\OperatorTok{=}\NormalTok{ y}
\NormalTok{ctr }\OperatorTok{+=} \DecValTok{1}
\ControlFlowTok{for}\NormalTok{ j }\KeywordTok{in} \BuiltInTok{range}\NormalTok{(}\DecValTok{1}\NormalTok{, N):}
  \BuiltInTok{print}\NormalTok{(}\StringTok{'step '}\NormalTok{, ctr,}\StringTok{' of '}\NormalTok{, N)}
  \ControlFlowTok{for}\NormalTok{ i }\KeywordTok{in} \BuiltInTok{range}\NormalTok{(}\DecValTok{1}\NormalTok{, M):}
\NormalTok{    y[i] }\OperatorTok{=} \DecValTok{2} \OperatorTok{*}\NormalTok{ y_1[i] }\OperatorTok{-}\NormalTok{ y_2[i] }\OperatorTok{+}\NormalTok{ C}\OperatorTok{**}\DecValTok{2} \OperatorTok{*}\NormalTok{ (y_1[i}\OperatorTok{+}\DecValTok{1}\NormalTok{] }\OperatorTok{-} \DecValTok{2}\OperatorTok{*}\NormalTok{y_1[i] }\OperatorTok{+}\NormalTok{ y_1[i}\DecValTok{-1}\NormalTok{])}
\NormalTok{  y[}\DecValTok{0}\NormalTok{] }\OperatorTok{=}\NormalTok{ y[}\DecValTok{0}\NormalTok{] }\OperatorTok{*} \FloatTok{0.5}
\NormalTok{  y[M] }\OperatorTok{=}\NormalTok{ y[M] }\OperatorTok{*} \FloatTok{0.5}
\NormalTok{  y }\OperatorTok{=}\NormalTok{ y }\OperatorTok{*}\NormalTok{ sustain}
\NormalTok{  w.append(y[pickup])}
\NormalTok{  y_2[:] }\OperatorTok{=}\NormalTok{ y_1 }
\NormalTok{  y_1[:] }\OperatorTok{=}\NormalTok{ y}
\NormalTok{  ctr }\OperatorTok{+=} \DecValTok{1}
\NormalTok{w }\OperatorTok{=}\NormalTok{ np.array(w) }\OperatorTok{*} \DecValTok{4000000}
\end{Highlighting}
\end{Shaded}

Implementação na linguagem Python do Método das Diferenças Finitas

Fonte: Elaboração própria

Uma animação relativa a uma onda de frequência real igual a 440 Hz,
criada pela implementação acima, pode ser encontrado no repositório da
dissertação (Tesserato, 2018), no caminho
\href{https://github.com/tesserato/tesserato.github.io/tree/master/resources/Demo\%20Finite\%20Difference}{``resources/Demo
Finite Difference/''}, junto com áudios de diversas ondas, com locais de
excitação e captação variados. É interessante observar o efeito da
variação desses parâmetros no timbre da onda sonora gerada. Cabe notar
que a velocidade e a escala da animação estão bastante distorcidas, para
facilitar a visualização.

\begin{figure}
\centering
\includegraphics{im/FD_440_Pluck_0-1_Pick_0-1.png}
\caption{Excitação: 0.1 - Captação: 0.1\label{fig:id}}
\end{figure}

Fonte: Elaboração própria

\begin{figure}
\centering
\includegraphics{im/FD_440_Pluck_0-1_Pick_0-5.png}
\caption{Excitação: 0.1 - Captação: 0.5\label{fig:id}}
\end{figure}

Fonte: Elaboração própria

\begin{figure}
\centering
\includegraphics{im/FD_440_Pluck_0-3_Pick_0-7.png}
\caption{Excitação: 0.3 - Captação: 0.7\label{fig:id}}
\end{figure}

Fonte: Elaboração própria

\begin{figure}
\centering
\includegraphics{im/FD_440_Pluck_0-5_Pick_0-1.png}
\caption{Excitação: 0.5 - Captação: 0.1\label{fig:id}}
\end{figure}

Fonte: Elaboração própria

\begin{figure}
\centering
\includegraphics{im/FD_440_Pluck_0-5_Pick_0-5.png}
\caption{Excitação: 0.5 - Captação: 0.5\label{fig:id}}
\end{figure}

Fonte: Elaboração própria

\hypertarget{digital-waveguides}{%
\subsubsection{Digital Waveguides}\label{digital-waveguides}}

Esse algoritmo consiste em uma abordagem simplificada da modelagem
física, na medida em que concentra em alguns pontos discretos os
cálculos necessários à simulação, aumentando bastante a eficiência
computacional do modelo (Van Duyne e Smith, 1995).

Conceitualmente, podem ser vistos como uma caso especial do método das
diferenças finitas (Van Duyne e Smith, 1993), e a maioria de suas
implementações consistem na utilização de \emph{delay lines} e filtros
digitais para a modelagem da propagação da onda (Smith, 2006). O modelo
também parte da equação da onda \xrefname{equação}\cref{eq:onda},
discretizada da forma apresentada no contexto do método das diferenças
finitas. A formulação analítica, no entanto, é abandonada, em favor de
uma interpretação mais diretamente baseada na interpretação das duas
ondas viajando em sentidos opostos na corda. A imagem abaixo apresenta a
intuição do método, que representa ainda hoje o estado da arte da
modelagem física (Bilbao, 2009).

\begin{figure}
\centering
\includegraphics{im/DigitalWaveguides.png}
\caption{Digital Waveguides\label{fig:waveguides}}
\end{figure}

Fonte: Elaboração própria

As duas linhas quadriculadas apresentadas na
\xrefname{figura}\cref{fig:waveguides} são as chamadas \emph{delay
lines}, por onde a onda circula inalterada, em uma direção: no exemplo,
na linha superior a perturbação viaja da esquerda para a direita, e da
direita para a esquerda na linha de baixo. Essa movimentação geralmente
ocorre em uma estrutura de dados denominada \emph{circular buffer}. O
movimento de reflexão da onda quando alcança as extremidades fixas da
corda é calculado nos pontos representados por cículos na
\xrefname{figura}\cref{fig:waveguides}, que são responsáveis por enviar
os \emph{samples} de um \emph{buffer} para o outro, invertendo seu
sinal; nesses pontos, geralmente a partir da utilização de filtros
digitais, outras manipulações são introduzidas, no sentido de emular
perdas e fatores relacionados à rigidez de uma corda não ideal. O
quadrado, na \xrefname{figura}\cref{fig:waveguides} representa a posição
onde a onda é amostrada, a partir da soma do conteúdo das duas
\emph{delay lines}. Esse sistema, como um todo, recebe o nome de
\emph{digital waveguide}.

Do ponto de vista físico os \emph{digital waveguides}, representam uma
corda fixa em suas terminações, por onde um pulso viaja, sendo refletida
nas extremidades. Ao passar de uma linha pra outra, a onda é espelhada e
invertida e todos os cálculos são concatenados nesses pontos ou,
dependendo da simulação, em apenas um desses pontos, como é o caso da
implementação apresentada no quadro 2, aumentando consideravelmente a
eficiência da simulação.

No modelo mais básico, um termo de perda é calculado durante a transição
de uma linha para outra, enquanto implementações mais sofisticadas fazem
uso dos filtros para emular além da rigidez das cordas, alguns outros
fenomenos mais sofisticados, como dispersão, por exemplo. O modelo pode
ser extendido também para emulações bidimensionais (Mullen, 2006) e
tridimensionais {[}Fontana, Rocchesso e Apollonio (2000); Laird (2001);
speed2012voice{]} , e presta-se à implementações bastante eficientes,
principalmente em uma dimensão.

No quadro 2, é apresentado um exemplo em Python, à guisa de
pseudocódigo. Observa-se que a implementação é mais simples do que a do
método das diferenças finitas, e mais eficiente, como a
\ref{tbl:wavexdiff} demonstra. Como antes, a versão funcional do código
pode ser acessada em Tesserato (2018), no caminho
\href{https://github.com/tesserato/tesserato.github.io/blob/master/resources/Digital_waveguide.py}{``resources/Digital\_waveguide.py''},
e uma animação ilustrando uma onda gerada pelo método, além dos com
diferentes valores para o ponto de excitação e captação áudios, podem
ser acessados no mesmo repositório, em
\href{https://github.com/tesserato/tesserato.github.io/tree/master/resources/Demo\%20Digital\%20Waveguide}{``resources/Demo
Digital Waveguide/''}

\begin{Shaded}
\begin{Highlighting}[numbers=left,,]
\NormalTok{path }\OperatorTok{=} \StringTok{'Demo Digital Waveguide/'}
\NormalTok{n }\OperatorTok{=} \DecValTok{44100}
\NormalTok{fps }\OperatorTok{=} \DecValTok{44100}
\NormalTok{frequency }\OperatorTok{=} \DecValTok{440}
\NormalTok{amplitude }\OperatorTok{=} \DecValTok{20000}
\NormalTok{pluck_position }\OperatorTok{=} \FloatTok{.1}
\NormalTok{pickup_position }\OperatorTok{=} \FloatTok{.1}
\NormalTok{sustain }\OperatorTok{=} \FloatTok{.99}
\NormalTok{smoothing }\OperatorTok{=} \DecValTok{3}
\NormalTok{plot }\OperatorTok{=} \VariableTok{True}

\NormalTok{L }\OperatorTok{=} \BuiltInTok{int}\NormalTok{(}\BuiltInTok{round}\NormalTok{(fps }\OperatorTok{/}\NormalTok{ (}\DecValTok{2} \OperatorTok{*}\NormalTok{ frequency)))}
\NormalTok{pickup }\OperatorTok{=} \BuiltInTok{int}\NormalTok{(}\BuiltInTok{round}\NormalTok{(L }\OperatorTok{*}\NormalTok{ pickup_position))}
\NormalTok{asc }\OperatorTok{=} \BuiltInTok{int}\NormalTok{(}\BuiltInTok{round}\NormalTok{(L }\OperatorTok{*}\NormalTok{ pluck_position))}
\NormalTok{dsc }\OperatorTok{=}\NormalTok{ L }\OperatorTok{-}\NormalTok{ asc}
\NormalTok{X_asc }\OperatorTok{=}\NormalTok{ np.linspace(}\DecValTok{0}\NormalTok{, }\DecValTok{1}\NormalTok{, asc)}
\NormalTok{X_dsc }\OperatorTok{=}\NormalTok{ np.linspace(}\DecValTok{0}\NormalTok{, }\DecValTok{1}\NormalTok{, dsc }\OperatorTok{+} \DecValTok{1}\NormalTok{)[::}\OperatorTok{-}\DecValTok{1}\NormalTok{][}\DecValTok{1}\NormalTok{: ]}
\NormalTok{delay_r }\OperatorTok{=}\NormalTok{ np.hstack([X_asc,X_dsc])}
\NormalTok{delay_l }\OperatorTok{=}\NormalTok{ np.zeros(L)}

\NormalTok{w }\OperatorTok{=}\NormalTok{ np.zeros(n)}
\NormalTok{zeros }\OperatorTok{=}\NormalTok{ np.zeros(L)}
\ControlFlowTok{for}\NormalTok{ i }\KeywordTok{in} \BuiltInTok{range}\NormalTok{(n):}
  \BuiltInTok{print}\NormalTok{(}\StringTok{'step '}\NormalTok{, i }\OperatorTok{+} \DecValTok{1}\NormalTok{,}\StringTok{' of '}\NormalTok{, n)}
\NormalTok{  w[i] }\OperatorTok{=}\NormalTok{ delay_r[pickup] }\OperatorTok{+}\NormalTok{ delay_l[pickup]}
\NormalTok{  to_l }\OperatorTok{=} \DecValTok{-1} \OperatorTok{*}\NormalTok{ np.average(delay_r[}\OperatorTok{-}\NormalTok{smoothing:]) }\OperatorTok{*}\NormalTok{ sustain }\CommentTok{# to add in delay_l[L-1] AFTER rolling}
\NormalTok{  delay_r }\OperatorTok{=}\NormalTok{ np.roll(delay_r, }\DecValTok{1}\NormalTok{)}
\NormalTok{  delay_r[}\DecValTok{0}\NormalTok{] }\OperatorTok{=} \DecValTok{-1} \OperatorTok{*}\NormalTok{ delay_l[}\DecValTok{0}\NormalTok{]}
\NormalTok{  delay_l }\OperatorTok{=}\NormalTok{ np.roll(delay_l, }\DecValTok{-1}\NormalTok{)}
\NormalTok{  delay_l[L}\DecValTok{-1}\NormalTok{] }\OperatorTok{=}\NormalTok{ to_l}
\end{Highlighting}
\end{Shaded}

Implementação na linguagem Python do algoritmo Digital Waveguides

Fonte: Elaboração própria

\begin{figure}
\centering
\includegraphics{im/DW_440_Pluck_0-1_Pick_0-1.png}
\caption{Excitação: 0.1 - Captação: 0.1\label{fig:id}}
\end{figure}

Fonte: Elaboração própria

\begin{figure}
\centering
\includegraphics{im/DW_440_Pluck_0-1_Pick_0-5.png}
\caption{Excitação: 0.1 - Captação: 0.5\label{fig:id}}
\end{figure}

Fonte: Elaboração própria

\begin{figure}
\centering
\includegraphics{im/DW_440_Pluck_0-3_Pick_0-7.png}
\caption{Excitação: 0.3 - Captação: 0.7\label{fig:id}}
\end{figure}

Fonte: Elaboração própria

\begin{figure}
\centering
\includegraphics{im/DW_440_Pluck_0-5_Pick_0-1.png}
\caption{Excitação: 0.5 - Captação: 0.1\label{fig:id}}
\end{figure}

Fonte: Elaboração própria

\begin{figure}
\centering
\includegraphics{im/DW_440_Pluck_0-5_Pick_0-5.png}
\caption{Excitação: 0.5 - Captação: 0.5\label{fig:id}}
\end{figure}

Fonte: Elaboração própria

\hypertarget{digital-waveguides-x-diferenuxe7as-finitas}{%
\subsubsection{Digital Waveguides X Diferenças
Finitas}\label{digital-waveguides-x-diferenuxe7as-finitas}}

Para a comparação dos dois algoritmos, alguns exemplos de ondas
equivalentes, geradas por ambos, foram preparados, a fim de basear um
julgamento subjetivo. A \xrefname{figura}\cref{fig:wavexdiff} compara o
primeiro frame e último das simulações da evolução do pulso em uma corda
para os dois métodos.

\begin{figure}
\centering
\includegraphics{im/DigitalWaveguidesXFiniteDifferences.png}
\caption{Animação da Corda - Diferenças Finital x Digital
Waveguide\label{fig:wavexdiff}}
\end{figure}

Fonte: Elaboração própria

A movimentação da onda produzida a partir do método das diferenças
finitas tem um caráter mais organico, enquanto é, também, mais propensa
a desvios da condição de estabilidade, enquanto a simulação a partir do
algoritmo de \emph{digital waveguides} é mais robusta, porém mais
mecânica Em relação às eficiências, abaixo são elencados os tempos que
foram necessários para a geração dos sons apresentados anteriormente.
Todas as amostras tem \emph{framerate} igual a 44100 FPS e duração de 1
segundo. Os tempos são apresentados em segundos.

\begin{longtable}[]{@{}cccc@{}}
\caption{Latência em segundos para uma onda com 44100 \emph{samples}
\label{tbl:wavexdiff}}\tabularnewline
\toprule
Excitação & Captação & Digital Waveguides & Diferenças
Finitas\tabularnewline
\midrule
\endfirsthead
\toprule
Excitação & Captação & Digital Waveguides & Diferenças
Finitas\tabularnewline
\midrule
\endhead
0.1 & 0.1 & 13.811148881912231 & 16.9521381855011\tabularnewline
0.1 & 0.5 & 13.020658254623413 & 19.620429277420044\tabularnewline
0.3 & 0.7 & 11.332738161087036 & 24.256462335586548\tabularnewline
0.5 & 0.1 & 12.350087404251099 & 17.89053726196289\tabularnewline
0.5 & 0.5 & 12.122233390808105 & 18.002467393875122\tabularnewline
\bottomrule
\end{longtable}

Fonte: Elaboração própria

\hypertarget{referencial-teuxf3rico}{%
\section{Referencial Teórico}\label{referencial-teuxf3rico}}

O objetivo deste capítulo é oferecer uma definição mais formal das
principais ferramentas utilizadas neste trabalho.

\hypertarget{transformada-discreta-de-fourier}{%
\subsection{Transformada Discreta de
Fourier}\label{transformada-discreta-de-fourier}}

A transformada de Fourier, em sua forma contínua, é um tipo de
transformada integral, como abaixo:

\begin{equation}
X(f) = \mathcal{F}\{x(t)\}=\int_{-\infty}^\infty x(t) e^{- 2 \pi f t \ i} dt
\label{eq:id}\end{equation}

\begin{equation}
x(t) = \mathcal{F}^{-1}\{X(f)\}=\int_{-\infty}^\infty X(f) e^{2 \pi f t \ i} df
\label{eq:id}\end{equation}

Tendo sido formulada por Fourier, enquanto investigava o fenômeno da
transferência de calor (Stein e Shakarchi, 2003), trata-se de uma das
ferramentas mais utilizadas na investigação de sistemas físicos (Lyons,
2011; Olson e Levitt, 2017).

Esta ferramenta parte da premissa de que qualquer sinal contínuo e
periódico pode ser representado a partir de uma soma de senoides (Smith
e others, 1997) permitindo que a descrição de um sinal \(x(t)\), em
relação ao tempo, seja transformada em uma representação deste mesmo
sinal em termos das frequências que o compõe, na forma
\(X(f)=\mathcal{F}\{x(t)\}\); \(x(t)\) e \(X(f)\) são denominados um par
de Fourier.

A versão discreta da transformada de Fourier, apresentada abaixo, é uma
das ferramentas mais importantes e utilizadas no campo de processamento
de sinais (Gazi, 2018); contribui para essa penetração a descoberta de
um algoritmo eficiente , denominado de \emph{Fast Fourier Transform},
que diminuiu consideravelmente o número de operações necessárias ao
cômputo da transformada discreta de Fourier (Pereyra e Ward, 2012).

\begin{equation}
X[m] =\sum_{n=0}^{N-1} x[n] e^{- 2 \pi m n / N ~ i}
\label{eq:id}\end{equation}

\begin{equation}
x[n] = \frac{1}{N}\sum_{m=0}^{N-1} X[m] e^{2 \pi m n / N ~ i}
\label{eq:id}\end{equation}

onde \(x\) e \(X\) são sequências de \(N\) números complexos. Utilizando
a fórmula de Euler, \(e^{x ~ i} = \cos(x) + \sin(x)~i\), podemos alterar
a definição acima, da forma exponencial, para a trigonométrica,
reorganizando o expoente de forma a evidenciar algumas características
importantes, na forma abaixo.

\[
X[m] =\sum_{n=0}^{N-1}
x[n] \cos \left(m ~ \frac{2 \pi}{N} ~ n \right)-
x[n] \sin\left(m ~ \frac{2 \pi}{N} ~ n \right) ~ i
\]

A forma acima torna mais fácil uma interpretação geométrica do
algoritmo. Para cada valor de \(m ~ \in ~ {0,1,2,...,N-1}\) é feita a
multiplicação, elemento a elemento, das \(N\) posições do impulso \(x\)
com \(N\) posições de um onda em forma de cosseno (\(C(m)\)), pura, de
``frequência'' \(m\), na parte real da equação. Analogamente, em sua
parte imaginária, o mesmo processo ocorre: cada posição medida do
impulso \(x\) também é multiplicada pela posição equivalente, desta vez,
de uma onda em formato de seno(\(S(m)\)), com ``frequência'' \(m\).

Note-se que esta ``frequência'', e por isso a insistência nas aspas,
refere-se a quantos ciclos das ondas (\(C(m)\)) e (\(S(m)\)) acima
definidas estão contidos no espaço dos \(N\) pontos utilizados na
transformada e não corresponde diretamente à frequência física da onda;
essa relação é regida por \(f_l = \frac{f_r}{FPS . N}\), onde \(f_l\) é
a frequência local, e \(f_r\) sua frequência real, física.

Para o caso de sinais que, no domínio do tempo, consistem de uma
sequência finita de números reais, como é o caso dos sinais considerados
neste trabalho, pode-se, à luz da definição acima, derivar uma
interessante e útil propriedade do algoritmo: A descrição do sinal no
domínio da frequência é conjugado-simétrica em torno de sua mediana, de
forma que \(X[m] = X^* [N-m]\), com \(X^*\) denotando o conjugado
complexo de \(X\). Dessa forma, sendo \(N\) par, precisamos de apenas de
\(N/2+1\) termos de \(X[m]\) para descrever completamente o pulso no
domínio da frequência, enquanto \((N+1)/2\) são suficientes caso \(N\)
seja ímpar. A \xrefname{figura}\cref{fig:rfft} ilustra geometricamente
essa propriedade.

\begin{figure}
\centering
\includegraphics{im/fourier.png}
\caption{Transformada Discreta Real de Fourier - Interpretação
Geométrica da Simetria\label{fig:rfft}}
\end{figure}

Fonte: Elaboração própria

Na \xrefname{figura}\cref{fig:rfft} é a presentado o resultado da
transformada discreta de uma onda composta por 8 amostras, que pode ser
vista no retangulo inferior da imagem. Os outros retangulos ilustram os
respectivos passos e a simetria entre passos equidistantes de \(n/2\). O
motivo da simetria, causado pela perda de informação ao tomar apenas
alguns pontos, torna-se explícito quando observado do ponto de vista
geométrico: os 8 pontos equidistantes tomados como amostra, por exemplo,
de uma senoide com frequência igual a 2 coincidem com os retirados de
uma senoide com frequência igual a 6, já que partes do ciclo são
ignorados.

O código em Python apresentado no quadro 3, utilizado para gerar os
quadros da \xrefname{figura}\cref{fig:rfft}, serve para elucidar
numericamente o exposto. É interessante notar que a interpretação
geométrica da simetria da transformada discreta de Fourier aplicada em
ondas puramente reais, até onde alcança o conhecimento do autor, é
apresentada pela primeira vez neste trabalho. Cabe observar que
implementações mais eficientes, como o referido algoritmo FFT,
utilizadas na prática, tiram proveito da simetria da transformada, entre
outras alterações mais sofisticadas, para economizar operações (Sorensen
\emph{et al.}, 1987).

\begin{Shaded}
\begin{Highlighting}[numbers=left,,]
\ImportTok{import}\NormalTok{ numpy }\ImportTok{as}\NormalTok{ np}
\ImportTok{import}\NormalTok{ matplotlib.pyplot }\ImportTok{as}\NormalTok{ plt}
\ImportTok{from}\NormalTok{ scipy.fftpack }\ImportTok{import}\NormalTok{ fft}
\CommentTok{# Discrete Fourier Transform | Transformada Discreta de Fourier}
\KeywordTok{def}\NormalTok{ DFT(x):}
\NormalTok{  N }\OperatorTok{=}\NormalTok{ x.shape[}\DecValTok{0}\NormalTok{]}
  \CommentTok{# cria um vetor de N entradas complexas, preenchido com zeros}
\NormalTok{  X }\OperatorTok{=}\NormalTok{ np.zeros(N, dtype}\OperatorTok{=}\StringTok{'complex64'}\NormalTok{)}
\NormalTok{  n }\OperatorTok{=}\NormalTok{ np.array([p }\ControlFlowTok{for}\NormalTok{ p }\KeywordTok{in} \BuiltInTok{range}\NormalTok{(N)]) }\CommentTok{# n = \{0,1,2,...,N-1\}}
  \CommentTok{# Versão de "alta resolução" de n, usada para o gráfico apenas}
\NormalTok{  Hn }\OperatorTok{=}\NormalTok{ np.linspace(}\DecValTok{0}\NormalTok{, N}\DecValTok{-1}\NormalTok{, }\DecValTok{100}\OperatorTok{*}\NormalTok{N) }\CommentTok{# Hn =\{0,...,1,...,2,...,N-1\}}
  \ControlFlowTok{for}\NormalTok{ m }\KeywordTok{in} \BuiltInTok{range}\NormalTok{(N): }\CommentTok{# m E \{0,1,2,...,N-1\}}
\NormalTok{    C_Si }\OperatorTok{=}\NormalTok{ np.exp(}\OperatorTok{-}\DecValTok{2} \OperatorTok{*}\NormalTok{ np.pi }\OperatorTok{*}\NormalTok{ m }\OperatorTok{*}\NormalTok{ n }\OperatorTok{/}\NormalTok{ N }\OperatorTok{*}\NormalTok{ 1j) }\CommentTok{# C(m) - S(m) i}
    \CommentTok{# Versão de "alta resolução" de C_Si, usada para o gráfico apenas}
\NormalTok{    HC_Si }\OperatorTok{=}\NormalTok{ np.exp(}\OperatorTok{-}\DecValTok{2} \OperatorTok{*}\NormalTok{ np.pi }\OperatorTok{*}\NormalTok{ m }\OperatorTok{*}\NormalTok{ Hn }\OperatorTok{/}\NormalTok{ N }\OperatorTok{*}\NormalTok{ 1j)}
\NormalTok{    X[m] }\OperatorTok{=} \BuiltInTok{sum}\NormalTok{(np.multiply(x, C_Si)) }\CommentTok{# multiplicação termo a termo}
    \CommentTok{# plotando:}
\NormalTok{    plt.figure(}\DecValTok{1}\NormalTok{)}
\NormalTok{    plt.suptitle(}\StringTok{'M ='} \OperatorTok{+} \BuiltInTok{str}\NormalTok{(m), fontsize}\OperatorTok{=}\DecValTok{16}\NormalTok{)}
\NormalTok{    plt.subplot(}\DecValTok{211}\NormalTok{)}
\NormalTok{    plt.plot(n,x[n].real, }\StringTok{'k:.'}\NormalTok{, label}\OperatorTok{=}\StringTok{'x[n]'}\NormalTok{)}
\NormalTok{    plt.plot(n,C_Si.real, }\StringTok{'k--o'}\NormalTok{, label}\OperatorTok{=}\StringTok{'C[n]'}\NormalTok{)}
\NormalTok{    plt.plot(Hn,HC_Si.real, }\StringTok{'k-'}\NormalTok{,label}\OperatorTok{=}\StringTok{'C[n] ideal'}\NormalTok{)}
\NormalTok{    plt.ylabel(}\StringTok{'Real'}\NormalTok{, fontsize}\OperatorTok{=}\DecValTok{14}\NormalTok{, color}\OperatorTok{=}\StringTok{'k'}\NormalTok{)    }
\NormalTok{    plt.subplot(}\DecValTok{212}\NormalTok{)}
\NormalTok{    plt.plot(n,x[n].imag, }\StringTok{'k:.'}\NormalTok{, label}\OperatorTok{=}\StringTok{'x[n]'}\NormalTok{)}
\NormalTok{    plt.plot(n,C_Si.imag, }\StringTok{'k--o'}\NormalTok{, label}\OperatorTok{=}\StringTok{'C[n]'}\NormalTok{)}
\NormalTok{    plt.plot(Hn,HC_Si.imag, }\StringTok{'k-'}\NormalTok{,label}\OperatorTok{=}\StringTok{'C[n] ideal'}\NormalTok{)}
\NormalTok{    plt.ylabel(}\StringTok{'Imaginário'}\NormalTok{, fontsize}\OperatorTok{=}\DecValTok{14}\NormalTok{, color}\OperatorTok{=}\StringTok{'k'}\NormalTok{)}
\NormalTok{    plt.legend()}
\NormalTok{    plt.savefig(}\StringTok{'0'} \OperatorTok{+} \BuiltInTok{str}\NormalTok{(m) }\OperatorTok{+} \StringTok{'.png'}\NormalTok{, dpi}\OperatorTok{=}\DecValTok{150}\NormalTok{)}
\NormalTok{    plt.close(}\StringTok{'all'}\NormalTok{)}
  \CommentTok{#calcula o erro em relação ao algoritmo da biblioteca Scipy}
\NormalTok{  error }\OperatorTok{=} \BuiltInTok{sum}\NormalTok{((X}\OperatorTok{-}\NormalTok{fft(x))}\OperatorTok{**}\DecValTok{2}\NormalTok{)}
  \BuiltInTok{print}\NormalTok{(}\BuiltInTok{round}\NormalTok{(error,}\DecValTok{5}\NormalTok{))}
  \CommentTok{# retorna a parte não redundante da transformada}
  \ControlFlowTok{return}\NormalTok{ X[}\DecValTok{0}\NormalTok{ : }\BuiltInTok{int}\NormalTok{(N}\OperatorTok{/}\DecValTok{2}\OperatorTok{+}\DecValTok{1}\NormalTok{)] }\ControlFlowTok{if}\NormalTok{ (N }\OperatorTok{%} \DecValTok{2} \OperatorTok{==} \DecValTok{0}\NormalTok{) }\ControlFlowTok{else}\NormalTok{ X[}\DecValTok{0}\NormalTok{ : }\BuiltInTok{int}\NormalTok{((N}\OperatorTok{+}\DecValTok{1}\NormalTok{)}\OperatorTok{/}\DecValTok{2}\NormalTok{)]}

\KeywordTok{def}\NormalTok{ s(t):}
  \ControlFlowTok{return}\NormalTok{ np.cos(}\DecValTok{2} \OperatorTok{*}\NormalTok{ np.pi }\OperatorTok{*}\NormalTok{ t)}

\NormalTok{t }\OperatorTok{=}\NormalTok{ np.linspace(}\DecValTok{0}\NormalTok{, }\DecValTok{1}\NormalTok{, }\DecValTok{8}\NormalTok{) }\CommentTok{# 8 pontos entre 0 e 1}
\NormalTok{x }\OperatorTok{=}\NormalTok{ s(t) }\CommentTok{# aplica a função ponto a ponto}
\NormalTok{X }\OperatorTok{=}\NormalTok{ DFT(x)}
\CommentTok{# plotando:}
\NormalTok{plt.figure(}\DecValTok{1}\NormalTok{)}
\NormalTok{plt.subplot(}\DecValTok{211}\NormalTok{)}
\NormalTok{plt.ylabel(}\StringTok{'Real'}\NormalTok{, fontsize}\OperatorTok{=}\DecValTok{14}\NormalTok{, color}\OperatorTok{=}\StringTok{'k'}\NormalTok{)}
\NormalTok{plt.stem(X.real,linefmt}\OperatorTok{=}\StringTok{'k--'}\NormalTok{,markerfmt}\OperatorTok{=}\StringTok{'ko'}\NormalTok{, basefmt}\OperatorTok{=}\StringTok{'k--'}\NormalTok{)}
\NormalTok{plt.subplot(}\DecValTok{212}\NormalTok{)}
\NormalTok{plt.stem(X.imag,linefmt}\OperatorTok{=}\StringTok{'k--'}\NormalTok{,markerfmt}\OperatorTok{=}\StringTok{'ks'}\NormalTok{, basefmt}\OperatorTok{=}\StringTok{'k--'}\NormalTok{)}
\NormalTok{plt.ylabel(}\StringTok{'Imaginário'}\NormalTok{, fontsize}\OperatorTok{=}\DecValTok{14}\NormalTok{, color}\OperatorTok{=}\StringTok{'k'}\NormalTok{)}
\NormalTok{plt.savefig(}\StringTok{'transform.png'}\NormalTok{, dpi}\OperatorTok{=}\DecValTok{150}\NormalTok{)}
\NormalTok{plt.close(}\StringTok{'all'}\NormalTok{)}
\end{Highlighting}
\end{Shaded}

Implementação didática na linguagem Python do algoritmo DFT

Fonte: Elaboração própria

\hypertarget{a-equauxe7uxe3o-da-onda}{%
\subsection{A Equação da Onda}\label{a-equauxe7uxe3o-da-onda}}

A equação diferencial que descreve o movimento de uma onda em duas
dimensões pode ser derivada a partir das leis de Newton, assumindo
algumas simplificações {[}@ garrett2017understanding{]}. Considerando a
\xrefname{figura}\cref{fig:corda}:

\begin{figure}
\centering
\includegraphics{im/waveequation.png}
\caption{Representação de um segmento de corda\label{fig:corda}}
\end{figure}

Fonte: Elaboração própria

e as seguintes definições:

\(y(x,t)\) a distância entre um ponto qualquer de uma corda e o eixo
horizontal no tempo \(t\) e na posição \(x\);

\(\theta(x,t)\) o angulo entre a tangente da corda na posição \(x\) e no
momento \(t\) e a direção horizontal;

\(\vec{T}(x,t)\) o vetor representando a tensão na corda, no ponto \(x\)
e no momento \(t\);

Assumindo ainda que qualquer ponto da corda move-se somente na posição
vertical temos que a resultante das forças é também vertical, e pode ser
descrita como:

\[F(x,t) = T(x + dx,t)\sin(\theta + d\theta) - T(x,t)\sin(\theta) = 
T(x + dx,t)\cos(\theta + d\theta) \tan(\theta + d\theta) - T(x,t)\cos(\theta)\tan(\theta)\]

se considerarmos somente pequenos deslocamentos verticais da corda em
relação à sua posição de equilíbrio, temos que \(\theta \ll 1\) e,
portanto \(\cos(\theta + d\theta) \approx \cos(\theta) \approx 1\).

Além disso, lembrando que as tangentes podem ser escritas como a
derivada parcial da função de deslocamento em relação à posição, temos:

\[F(x,t) = T(x+d x, t) \left(\frac{\partial y(x+d x)}{\partial x}\right) - T(x, t) \left(\frac{\partial y(x)}{\partial x}\right)\]

Assumindo tensão uniforme e invariável em toda a extensão da corda,
condições bastante razoáveis no contexto de instrumentos musicais,
podemos escrever \(T(x+d x, t) = T(x, t) = T\) e \(F(x,t)\) toma a forma
\(T \left(\frac{\partial y(x+d x)}{\partial x} - \frac{\partial y(x)}{\partial x}\right)\),
quando colocamos \(T\) em evidência. Se assumirmos que a massa do
segmento infinitesimal da corda pode ser escrita na forma \(\rho dx\),
onde \(\rho\) é a massa da corda por unidade de comprimento e
\(dx \approx \sqrt{dx^2+dy^2}\) para oscilações pequenas, aplicando a
segunda lei de Newton para força resultante na corda temos \$F(x,t) =
\rho dx \frac{\partial^2 y}{\partial t^2} \$.

Assim, temos
\(T \left(\frac{\partial y(x+d x)}{\partial x} - \frac{\partial y(x)}{\partial x}\right) = \rho dx \frac{\partial^2 y}{\partial t^2}\)
que podemos reorganizar como
\(\frac{\partial^2 y}{\partial t^2} = \frac{T}{\rho} \frac{\left(\frac{\partial y(x+d x)}{\partial x} - \frac{\partial y(x)}{\partial x}\right)}{dx}\).
Notando que o termo à direita é a segunda derivada de \(y\) em relação a
\(x\), chegamos à equação da onda:

\[\frac{\partial^2 y(x,t)}{\partial t^2} =\frac{T}{\rho} \frac{\partial^2 y(x,t)}{\partial x^2}\]

onde \(v = \sqrt{\frac{T}{\rho}}\) é a velocidade de propagação da onda
na corda.

Uma solução para essa equação, quando não há uma fonte de excitação, foi
proposta por d'Alembert (Chaigne e Kergomard, 2016), na forma:

\[y(x,t)=F(x+vt)-G(x-vt)\]

Essa equação pode ser intepretada como dois pulsos viajando em sentidos
opostos em uma corda ideal com velocidade \(v\).

Em \(t=0\) podemos escrever:

\(y_0(x)=y(x,0)=F(x)-G(x)\)

\(v_0(x)=y'(x,0)=vF'(x)-vG'(x)\)

\(vF(x)-vG(x)= \int_{-\infty}^x v_0(\epsilon) d\epsilon\)

O que nos dá um sistema de equações que, resolvido, nos permite
reescrever a equação de d'Alembert em função das condições iniciais de
deslocamento e velocidade na corda, como abaixo:

\[y(x,t)= \frac{1}{2} \Big(y_0(x+vt)-y_0(x-vt) \Big) + \frac{1}{2v}\Big(\int_{x-vt}^{x+vt} v_0(\epsilon) d\epsilon \Big)\]

Essa formulação é importante já que em muitos casos pode-se trabalhar
apenas com o deslocamento inicial da corda, o que torna o termo à
direita zero.

Alternativamente, pelo método da separação de variáveis, pode-se obter a
solução em termos de um soma infinita de senoides estacionárias(Gracia e
Sanz-Perela, 2016):

\[y(x,t)= \sum_{n=1}^{\infty} \big(a_n \cos(2 \pi f_n t) +  b_n \sin(2 \pi f_n t\big) \sin(\frac{n \pi x}{L})\]

onde considera-se uma corda de comprimento \(L\) fixa em suas
extremidades. \(f_n\) são as frequências parciais da corda e são dados
pela equação \(\frac{nv}{2l}\). Além disso, os coeficientes \(a_n\) e
\(b_n\) são dados pelos coeficiente da trasformada de Fourier das
condições iniciais em \(t=0\) (\(y_0(x)\) e \(v_0(x)\)) (Salsa, 2016)

essas duas interpretações dão origem às duas escolas distintas de
modelagem acústica, como referido anteriormente, uma focada no domínio
do tempo e outra no domínio da frequência. Repare, no entanto, que ambas
as soluções modelam uma corda ideal, não levando em conta a rigidez
encontrada em cordas reais.

No domínio do tempo, (Smith, 1992) aponta que a rigidez implica no fato
de que ondas com diferentes frequências viajam através da corda em
velocidades diferentes, relação regida por
\(c(w) = c_0(\frac{1+kw^2}{2Kc_0^2})\) com \(k = E \pi r^4 \pi / 4\)
onde \(E\) é o modulo de Young e \(r\) é o raio da corda.

No domínio da frequência (Rigaud, David e Daudet, 2013) demonstra que a
rigidez pode ser incorporada ao modelo através da substituição do termo
que descreve as frequências de cada parcial por
\(fn = f_0 \sqrt{1 + Bn^2}\) onde
\(f_0 = \frac{1}{2l} \sqrt{\frac{T}{\rho}}\) e
\(B = \frac{\pi^3Ed^4}{64Tl^2}\) onde \(E\) é o modulo de Young e \(d\)
é o diâmetro da corda.

A maneira mais simples de acomodar o decaimento devido à forças
dissipativas, completando a teoria necessária ao presente trabalho, é
modular as amplitudes em ambas as soluções a partir de um termo da forma
\(e^{-\alpha t}\)

\hypertarget{formalizauxe7uxe3o-de-uma-rede-neural---forward-pass}{%
\subsection{Formalização de uma Rede Neural - Forward
Pass}\label{formalizauxe7uxe3o-de-uma-rede-neural---forward-pass}}

Definindo-se de maneira esquemática uma camada de uma rede neural
artificial conforme a \xrefname{figura}\cref{fig:camada}:

\begin{figure}
\centering
\includegraphics[width=2.08333in,height=\textheight]{im/layer.png}
\caption{representação de uma camada de uma rede\label{fig:camada}}
\end{figure}

Fonte: Elaboração própria

Onde:

\(I_{1 \times m} := [i_1 \dots i_m]\) é um vetor de inputs à rede, e
\(T_{1 \times n} := [t_1 \dots t_n]\)um vetor de alvos, de forma que
para cada um dos ventores de input há um vetor de alvo correspondente,
com a mesma dimensão da saída da camada;

\[ W_{m \times n} :=
\left[ \begin{matrix}
w_{11} & \dots & w_{1n} \\
\vdots & \ddots & \vdots \\
w_{m1} & \dots & w_{mn}
\end{matrix} \right] \]

é uma matriz representando os pesos das ligações entre os neurônios de
entrada e saída, de forma que \(w_{ij}\) seria o peso da ligação entre o
iésimo neurônio de entrada e o jésimo neurônio de saída.

\(A_{1 \times n} := [a_1 \dots a_n] \ |\
a_y := w_{1y}i_1 + \dots + w_{ny}i_n\)

representa cada elemento da multiplicação entre o vetor de entrada e os
pesos da camada, antes da aplicação, elemento à elemento, da função de
ativação;

\(O_{1 \times n} := [o_1 \dots o_n] \ |\
o_y := f(a_y) = f(w_{1y}i_1 + \dots + w_{ny}i_n)\)

é um vetor representando as saídas finais da camada;

\(E_{1 \times n} := [e_1 \dots e_n] \ |\
e_y := \varepsilon (o_y, t_y) = \varepsilon (f(w_{1y}i_1 + \dots + w_{ny}i_n) , t_y)\)

representa os erros individuais da rede para cada uma das saída
motivadas por cada um dos inputs.

Cabe notar que \(\varepsilon (o_y, t_y) = (o_y - t_y)^2\) é a forma mais
usual de computo dos erros individuais; define-se ainda o escalar
\(q := e_1 + \dots + e_n\) como a soma dos erros individuais para cada
um dos vetores de erros.

De posse das definições acima, somos capazes de representar de forma
compacta, em notação matricial, o movimento de predição (\emph{forward
pass}) de uma camada, como será visto. E, por extensão, de uma rede, na
medida em que estas podem ser obtidas pela justaposição de um número
arbitrários de camadas.

O mesmo pensamento, de forma simétrica, presta-se à definição do
movimento que atualizará os pesos (\emph{backward pass}), e que é
formalizado na próxima seção. Assim, derivaremos a seguir a forma básica
do algoritmo de \emph{backpropagation}.

Além de introduzir um referencial matemático para o algoritmo, isso
permitirá que suas extensões, como Adam e Adagrad, utilizadas adiante,
possam ser melhor compreendidas.

\hypertarget{formalizauxe7uxe3o-de-uma-rede-neural---backpropagation}{%
\subsection{Formalização de uma Rede Neural -
Backpropagation}\label{formalizauxe7uxe3o-de-uma-rede-neural---backpropagation}}

Tem-se que o gradiente do tensor de erros em relação a um peso
arbitrário de uma camada pode ser escrito como na
\xrefname{equação}\cref{eq:grad}, com aplicação direta da regra da
cadeia, e observando a seguinte notação para uma função qualquer
\(h(x)\) : \(\frac{\partial}{\partial x}h(x) = h'(x)\).

\begin{equation}
\frac{\partial q }{\partial w_{xy}} =
\varepsilon' (o_y, t_y) f'(a_y)i_x
\quad \text{Eq.:(I)}
\label{eq:grad}\end{equation}

A matriz de incrementos para cada um dos pesos, a cada iteração, pode
ser definida como abaixo, com a adição de uma taxa de aprendizado
\(\alpha\) com sinal negativo. Assim é, uma vez que o gradiente acima
aponta para sentido de maior crescimento dos erros no espaço dos pesos:
minimizar os erros implica, portanto, em mover os pesos em sentido
oposto.

\[
\Delta W_{m \times n} :=
\left[ \begin{matrix}
\Delta w_{11} & \dots & \Delta w_{1n} \\
\vdots & \ddots & \vdots \\
\Delta w_{m1} & \dots & \Delta w_{mn}
\end{matrix} \right]  \ |\
\Delta w_{xy} := - \alpha \frac{\partial q}{\partial w_{xy}} = - \alpha \varepsilon' (o_y, t_y) f'(a_y)i_x
\]

Resta definir uma forma de propagar os erros às entradas da rede,
permitindo assim que um número arbitrário de camadas sejam conectadas e
treinadas. Para tanto, convém considerar uma camada anterior à rede em
tela, com a forma e notação apresentadas na
\xrefname{figura}\cref{fig:camadas}:

\begin{figure}
\centering
\includegraphics[width=4.16667in,height=\textheight]{im/back02.png}
\caption{Duas camadas sobrepostas\label{fig:camadas}}
\end{figure}

Fonte: Elaboração própria

E, de forma análoga:

\[
\begin{align}
Z_{1 \times m} :=& [z_1 \dots z_m] \ |\
z_x := u_{1x}p_1 + \dots + u_{kx}p_k
\\
I_{1 \times m} =& [i_1 \dots i_m] \ |\
i_x := g(z_x) = g(u_{1x}p_1 + \dots + u_{kx}p_k)
\end{align}
\]

\[
\begin{align}
\frac{\partial q}{\partial u_{rx}} =&
[
\varepsilon' (o_1, t_1) f'(a_1) w_{x1} g'(z_x)p_r
+ \dots +
\varepsilon' (o_n, t_n) f'(a_n) w_{xn} g'(z_x)p_r
]
\\ =&
[
\varepsilon' (o_1, t_1) f'(a_1) w_{x1}
+ \dots +
\varepsilon' (o_n, t_n) f'(a_n) w_{xn}
] g'(z_x)p_r
\end{align}
\]

Comparando com \(\text{I}\) é fácil observar que o somatório na primeira
parte da última expressão corresponde a derivada do erro na primeira
equação. Assim, Podemos definir o tensor do erro propagado da forma
apresentada abaixo:

\[
\begin{align}
E_{1 \times m}^b :=& [e_1^b \dots e_m^b] \ |\
e_x^b := \varepsilon' (o_1, t_1)f'(a_1)w_{x1}
+ \dots +
\varepsilon' (o_n, t_n) f'(a_n)w_{xn}
\end{align}
\]

Fixamos, dessa forma, todos os termos necessários à uma arquitetura
modular: para camadas ocultas, para as quais não há como aferirmos
diretamente o erro, o mesmo é propagado a partir das camadas
posteriores. Em forma vetorial, pode-se sintetizar o algoritmo como na
\xrefname{equação}\cref{eq:net}, observando a conveniência de
introduzirmos o tensor \(H_{1 \times n}\) para evitar redundância nos
cálculos, e a notação \(\odot\) denotando o produto de Hadamard
(elemento a elemento) entre dois tensores:

\begin{equation}
\begin{align}
O =& f(IW) \\
H_{1 \times n} :=& E \odot f'(A)\\
\Delta W =& -\alpha I^t H= - \alpha I^t (E \odot f'(A)) \\
E^b =& HW^t= (E \odot f'(A)) W^t
\end{align}
\label{eq:net}\end{equation}

\hypertarget{metodologia}{%
\section{Metodologia}\label{metodologia}}

\hypertarget{amostras-sonoras-utilizadas}{%
\subsection{Amostras Sonoras
Utilizadas}\label{amostras-sonoras-utilizadas}}

Para a busca da coleção de amostras sonoras a serem utilizados para o
treinamento da rede, algumas características são essenciais. Em primeiro
lugar, uma licença permissiva, que não limite a utilização, modificação
e posterior divulgação dos resultados obtidos deve ser explicitamente
atribuída pelo autor.

Além disso, para o caso da bateria, todas as peças de um \emph{kit}
ordinário devem estar presentes, e idealmente, organizadas de maneira
inteligível, enquanto que para instrumentos temperados é bastante útil
algum tipo de indicação da frequência fundamental de cada uma das
amostras.

É desejável também que, para cada peça, amostras representando várias
dinâmicas tenham sido gravados, com especial atenção à intensidade com
que cada peça é golpeada. Interessante ainda é a presença dos chamados
\emph{round-robins}: gravações redundantes de cada uma das amostras, que
apresentam uma ideia de como fatores aleatórios influenciam no som
produzido pelo instrumento.

Com isso em mente, procedeu-se à busca, em sites e blogs especializados,
de indicações sobre trabalhos disponíveis que pudessem enquadrar-se nas
condições citadas. No caso do \emph{kit} de bateria, dois trabalhos em
especial foram considerados:

O primeiro é um esforço de disponibilizar, através da plataforma
\emph{Github}, uma coleção livre de sons de bateria. O trabalho leva o
nome de The Open Source Drumkit (Crabacus, 2015) em consiste em uma
coleção de amostras em formato .wav, separadas em pastas e nomeadas de
acordo com a intensidade do golpe, e as peças específicas de um kit de
bateria convencional, apresentando uma média de 10 articulações por
peça.

O segundo é um esforço para a elaboração de um instrumento virtual de
bateria denominado DrumGizmo (``DrumGizmo Wiki'', {[}s.d.{]}), e
apresenta samples gravadas a partir de 5 kits físicos de bateria. O
número de articulações é bem maior, girando em torno de 20 por kit, e as
gravações apresentam múltiplos microfones, com pelo menos um microfone
por peça e são, em geral, de melhor qualidade. Além disso, a maioria dos
\emph{kits} adere à licença Creative Commons Attribution 4.0
International que permite o livre uso e adaptação do material
disponibilizado em trabalhos derivados.

No caso das amostras de piano, em (``Ivy Audio'', {[}s.d.{]}) pode ser
encontrado, entre outros, um DMI baseado em samples de um Steinway Model
B grand piano, com 5 dinâmicas e 2 \emph{round-robins} para cada tecla,
além de um total de 4 condições de microfonação. Detalhes dos termos de
uso são discutidos na página de FAQ do site.

No site do University of Iowa Electronic Music Studios (``Musical
Instrument Samples'', {[}s.d.{]}) um conjunto bastante semelhante de
amostras do mesmo piano pode ser encontrado, com a vantagem de uma
atribuição mais livre de condições de uso; as amostras, no entanto,
diretamente hospedadas no site, devem ser baixadas uma por vez.

Uma vasta gama de amostras sonoras de instrumentos clássicos, com
variadas dinâmicas, é oferecida no site da Philharmonia Orchestra
(``Sound Samples'', {[}s.d.{]}).

\hypertarget{redes-densas-1}{%
\subsection{Redes Densas}\label{redes-densas-1}}

É natural que a investigação tenha início com uma análise da arquitetura
Feed-Forward, com redes compostas de um número arbitrário de camadas
totalmente conectadas. Isso possibilita formar uma ideia de como essa
arquitetura básica se comporta em um ambiente de alta dimensionalidade,
pavimentando o caminho para a subsequente incorporação de novas
arquiteturas e formulações na medida em que limites de aplicabilidade
forem sendo encontrados.

Convém lembrar que o trabalho tem como foco a síntese sonora em tempo
real, de forma que o tamanho, no sentido do número de neurons de uma
rede neural, e sua complexidade, aspecto relacionado à arquitetura,
tornam-se fatores potencialmente limitantes, e passam a requerer
especial atenção. Redes convolucionais, por exemplo, devem ser tratadas
com extrema cautela, por não possuírem um histórico de implementações
eficientes.

Em primeiro lugar, procedeu-se ao tratamento dos samples: foram
escolhidos, para esta investigação inicial, 5 tambores (os 4 tons
disponíveis, e o bumbo esquerdo) do kit denominado Aasimonster
disponibilizado pela iniciativa DrumGizmo, de cada um dos quais 20
dinâmicas foram escolhidas, gerando um total de 100 amostras, cada um
com 16 canais (1 por microfone utilizado na gravação).

foram extraídos os canais pertinentes à cada peça, resultando em 100
samples monofonicos em formato .wav, com \emph{framerate} de 44100
frames por segundo. As dinâmicas para cada peça foram normalizadas em
grupos de 4, de forma a emular 4 samples redundantes
(\emph{round-robins}) para 5 níveis diferentes de intensidade por peça.
Cada uma das amostras foi truncada, para possuir exatamente um segundo
de duração, e um efeito \emph{fade-off} foi aplicado às amostras que
estendiam-se além desse intervalo.

Esse conjunto de samples, de forma bastante direta, representou os alvos
da rede neural, na forma de um vetor de dimensão 100 x 44100, como cada
uma de suas linhas consistindo em uma representação digital da onda
sonora de cada amostra. A conversão foi efetuada utilizando o módulo
Pywave da linguagem Python, e cada vetor foi normalizado no espaço
{[}-1,1{]}. Os samples foram renomeados de forma a serem mais
descritivos.

Exemplificando, temos que a amostra de rótulo 154, corresponde à quarta
amostra redundante referente à primeira peça (o tom mais agudo),
golpeada com intensidade 5 (maior intensidade). Assim, as amostras 321,
322, 323 e 324, por exemplo, correspondem à quatro variações do mesmo
evento: cada um deles representa a terceira peça golpeada com
intensidade 2. Os dados de entrada da simplesmente refletem essa
nomenclatura em um intervalo {[}0,1{]} e são apresentados na
\xrefname{tabela}\cref{tbl:nom}.

\begin{longtable}[]{@{}llllll@{}}
\caption{Nomenclatura utilizada \label{tbl:nom}}\tabularnewline
\toprule
Peça & Peça Norm. & Intensidade & Intensidade Norm. & Sample \# & Sample
\# Norm.\tabularnewline
\midrule
\endfirsthead
\toprule
Peça & Peça Norm. & Intensidade & Intensidade Norm. & Sample \# & Sample
\# Norm.\tabularnewline
\midrule
\endhead
1 & 0.00 & 1 & 0.00 & 1 & 0.00\tabularnewline
2 & 0.25 & 2 & 0.25 & 2 & 0.33\tabularnewline
3 & 0.50 & 3 & 0.50 & 3 & 0.66\tabularnewline
4 & 0.75 & 4 & 0.75 & 4 & 1.00\tabularnewline
5 & 1.00 & 5 & 1.00 & - & -\tabularnewline
\bottomrule
\end{longtable}

Fonte: Elaboração própria

O que pretende-se é forçar uma associação de cada uma das
características da onda à cada uma das 3 dimensões do vetor de entrada.
Esse rationale tem algumas implicações importantes: critérios de
validação, embora possíveis de construir, perdem bastante do seu
significado e foram deixados de lado em favor de um julgamento subjetivo
da qualidade das amostras geradas.

\hypertarget{algoritmos-de-otimizauxe7uxe3o}{%
\subsubsection{Algoritmos de
Otimização}\label{algoritmos-de-otimizauxe7uxe3o}}

Utilizando as 100 amostras, processados como descrito, como alvos, e
seus rótulos normalizados como entradas, buscou-se investigar os
hiperparâmetros de uma rede neural densa mais adequados ao apredizado e
reprodução das ondas; essa é uma etapa importante, pois vários
parametros aqui definidos serão assumidos como ótimos e utilizados em
etapas posteriores, em arquiteturas diferentes.

Essa necessidade tem origem nas restrições de tempo e poder
computacional, que impedem a realização de novos testes dessa magnitude
para cada uma das arquiteturas futuramente investigadas. A plataforma
escolhida para a implementação das redes foi o Keras, rodando com o
Tensorflow como \emph{backend}. Essa escolha permite aproveitar a
velocidade de experimentação conferida pela arquitetura em nível mais
alto da biblioteca Keras sem comprometer muito a flexibilidade.

Em primeiro lugar, é conveniente investigar o melhor algoritmo de
atualização dos pesos: embora, conceitualmente, mesmo a versão pura da
descida em gradiente, utilizando uma taxa de aprendizado adequada, leve
à eventual convergência da rede (Ruder, 2016), na prática a velocidade
dessa convergência varia consideravelmente dependendo do método
utilizado, e um algoritmo eficiente permite a realização de mais testes
no mesmo intervalo de tempo.

A literatura disponível, embora esclarecedora, não é unanime ao apontar
uma técnica universalmente superior às outras, de forma que uma
investigação empírica faz-se necessária, e é descrita a seguir.

testou-se, em uma rede com uma camada oculta composta por 100 neurons,
em lotes (\emph{batches}) de 100 vetores de entrada e durante 500 épocas
os seguintes algoritmos de otimização: SGD com uma \emph{learning rate}
de 0,01, o algoritmo experimental RMSProp (Tieleman e Hinton, 2012), os
otimizadores Adam (Kingma e Ba, 2014; Reddi, Kale e Kumar, 2018),
Adagrad (Duchi, Hazan e Singer, 2011), Adadelta (Zeiler, 2012),
Adamax(Kingma e Ba, 2014) e Nadam (Dozat, 2016; Sutskever \emph{et al.},
2013). A função de ativação utilizada em todas as camadas foi a tangente
hiperbólica, e a função de perda foi a média dos quadrados dos erros.
Foi tomado cuidado para que todos os testes iniciassem em condições
pseudo-randomicas idênticas.

SGD e Adadelta ofereceram resultados bastante semelhantes, e aquém dos
outros. Todos os outros convergiram de forma semelhante, sendo o Adagrad
o mais rápido de todos, e Nadam o mais eficaz na diminuição da função de
perda.

A comparação por tempo relativo, disponível a partir do Tensorboard,
permite visualizar a superioridade do algoritmo Nadam: com 2,57 minutos
de treino, o tempo em que o treinamento foi concluído com o algoritmo
Adagrad, ele proporcionou o menor valor da função objetivo.

Convém investigar, dessa forma, o comportamento da rede treinada com
ambos os algoritmos, já que o parâmetro mais importante é a rapidez com
que a rede alcança um nível de erro capaz de gerar samples com qualidade
satisfatória. Após essa comparação exploratória, mais profunda,
poderemos identificar tanto o valor limite da função de perda, quanto o
tempo em que a rede mais eficiente leva para alcançá-lo, utilizando
essas informações como parâmetros para investigações posteriores,
inclusive de outras arquiteturas.

\hypertarget{grid-search-para-topologia-e-paruxe2metros}{%
\subsubsection{\texorpdfstring{\emph{Grid search} para topologia e
parâmetros}{Grid search para topologia e parâmetros}}\label{grid-search-para-topologia-e-paruxe2metros}}

Fixado os potenciais algoritmos de otimização, procedeu-se a um
\emph{grid search}, para investigar o número de camadas ocultas,
respectivos neurons e o efeitos dos tamanhos dos \emph{batchs} , tanto
em relação ao impacto no tempo total de treinamento quanto à rapidez (e
limite) de convergência.

\hypertarget{adagrad}{%
\paragraph{Adagrad}\label{adagrad}}

Testou-se \emph{batches} de 1, 10, 50, 100, 150 e 200 samples para redes
com 25, 50, 75 e 100 neurônios em cada camada oculta, em um intervalo de
0 a 3 camadas ocultas, gerando um total de 82 redes distintas, que foram
treinadas durantes 500 épocas.

Para uma rede sem camadas ocultas, um \emph{batch} igual a 1 foi o mais
efetivo, em 500 épocas, porém o mais lento. \emph{batches} de 100 e 50
foram os mais rápidos. O impacto no valor final da perda, no entanto,
foi desprezível para todos os tamanhos de \emph{batches}. para um camada
oculta com 25 neurons, \emph{batches} de 100 e 150 neurons ofereceram
resultados idênticos, sendo os mais rápidos e mais eficientes.

A partir desse ponto, \emph{batches} de 100, 150 e 200 começaram a
apresentar comportamentos semelhantes. Cabe notar que \emph{batches}
unitários parecem piorar muito o tempo de execução na medida em que o
número de neurônios aumenta. temos abaixo os melhores resultados por
\emph{batch} e seus respectivos tempos. Todos eles encontram-se nas
redes com 3 camadas ocultas.

\begin{figure}
\centering
\includegraphics{im/adagrad loss x step.png}
\caption{Adagrad: Loss x Step - 3 camadas ocultas\label{fig:id}}
\end{figure}

Fonte: Elaboração própria

\begin{figure}
\centering
\includegraphics{im/adagrad loss x time.png}
\caption{Adagrad: Loss x Time - 3 camadas ocultas\label{fig:id}}
\end{figure}

Fonte: Elaboração própria

\begin{longtable}[]{@{}llll@{}}
\caption{Menores valores da perda após 500 iterações
\label{tbl:id}}\tabularnewline
\toprule
Batches & neurônios por camada oculta & Loss (e-3) &
tempo\tabularnewline
\midrule
\endfirsthead
\toprule
Batches & neurônios por camada oculta & Loss (e-3) &
tempo\tabularnewline
\midrule
\endhead
200 & 75 & 4,0934 & 2m 48s\tabularnewline
100 & 100 & 4,1426 & 3m 41s\tabularnewline
150 & 100 & 4,1644 & 3m 17s\tabularnewline
1 & 50 & 4,4069 & 2m 28s\tabularnewline
10 & 75 & 5,6130 & 8m 15s\tabularnewline
50 & 25 & 5,7133 & 2m 24s\tabularnewline
\bottomrule
\end{longtable}

Fonte: Elaboração própria

Nesse ponto fica evidente que uma arquitetura formada por 3 camadas
ocultas de 75 neurons é suficiente para aprender a uma representação
direta das ondas, e a mais eficiente sob o otimizador Adagrad. O gráfico
da função de perda em função do tempo de treinamento torna essa
visualização mais fácil.

\hypertarget{nadam}{%
\paragraph{Nadam}\label{nadam}}

Para essa etapa \emph{batches} de 10, 50 e 100 amostras sonoras foram
utilizados; dispensou-se \emph{batches} unitários, pelo comportamento
ineficiente identificado, e os \emph{batches} acima de 100, já que não
introduziram melhorias significativas. Foram testadas, para cada um
desses \emph{batches}, redes com todas as combinações entre 25, 50, 75 e
100 neurons, até o limite de 3 camadas ocultas. Os resultados são
apresentados abaixo, para as arquiteturas com 3 camadas.

\begin{figure}
\centering
\includegraphics[width=8.33333in,height=\textheight]{im/nadam loss x step.png}
\caption{Nadam: Loss x Step - 3 camadas ocultas\label{fig:id}}
\end{figure}

Fonte: Elaboração própria

\begin{figure}
\centering
\includegraphics[width=8.33333in,height=\textheight]{im/nadam loss x time.png}
\caption{Nadam: Loss x Time - 3 camadas ocultas\label{fig:id}}
\end{figure}

Fonte: Elaboração própria

\begin{longtable}[]{@{}llll@{}}
\caption{Menores valores da perda após 500 iterações
\label{tbl:id}}\tabularnewline
\toprule
Batches & neurônios por camada oculta & Loss (e-3) &
tempo\tabularnewline
\midrule
\endfirsthead
\toprule
Batches & neurônios por camada oculta & Loss (e-3) &
tempo\tabularnewline
\midrule
\endhead
10 & 75 100 100 & 0.44632 & 16m 59 s\tabularnewline
10 & 50 50 100 & 1.5559 & 15m\tabularnewline
100 & 25 75 75 & 2,4384 & 3m 6 s\tabularnewline
50 & 75 75 75 & 2.2692 & 3m 36 s\tabularnewline
100 & 75 75 75 & 2.2692 & 3m 38 s\tabularnewline
50 & 75 100 & 4.9025 & 5m 4s\tabularnewline
\bottomrule
\end{longtable}

Fonte: Elaboração própria

Uma rede feedforward com 3 camadas de 75 neurons parece oferecer a
arquitetura ideal, com ambos os otimizadores. Esses hisperparametros
servirão da base para arquiteturas futuras desenvolvidas neste trabalho.

\hypertarget{ativauxe7uxf5es-e-otimizadores}{%
\paragraph{Ativações e
Otimizadores}\label{ativauxe7uxf5es-e-otimizadores}}

A performance dos dois, para 2000 \emph{steps}, é comparada abaixo;
também é comparado o impacto da função de ativação softsign, em relação
à tanh utilizada até agora.

\begin{figure}
\centering
\includegraphics{im/softsign tanh adagrad nadam.png}
\caption{Comparação: Ativações e Otimizadores\label{fig:id}}
\end{figure}

Fonte: Elaboração própria

Podemos ver que a combinação do otimizador Nadam com a função de
ativação tanh é a mais eficiente.

\hypertarget{anuxe1lise-dos-resultados}{%
\subsubsection{Análise dos resultados}\label{anuxe1lise-dos-resultados}}

Os testes acima proporcionam uma visão objetiva sobre o comportamento
geral da arquitetura densa na resolução do problema em tela. Faz-se
necessário, contudo, uma análise das amostras de som geradas pelas redes
investigadas, de forma a entender, do ponto de vista sensorial, sua
verossimilhança, além da capacidade de generalização da rede.

Tomando como métrica a o valor da função de perda ao final do
treinamento, podemos selecionar algumas redes para investigar seu
comportamento.

Analisando a rede com maior valor de função de perda após o treinamento
(a rede sem camadas ocultas treinada com Adagrad em \emph{batches} de
100), a rede escolhida como baseline (75 neurons em cada uma de 3
camadas ocultas, otimizador Nadam, batches de 100) e a rede com a menor
perda após o treinamento (a rede com 75, 100 e 100 neurons em suas
camadas ocultas, treinada com Nadam em \emph{batches} de 10),
desenvolveremos alguma intuição acerca da relação entre a função de
perda e a verossimilhança das amostras produzidas. Os dados são
sintetizados na \xrefname{tabela}\cref{tbl:comp}:

\begin{longtable}[]{@{}llllll@{}}
\caption{Redes para análise das amostras sonoras
\label{tbl:comp}}\tabularnewline
\toprule
N° & Camadas & Batch & Otimizador & Valor final da função de perda &
Parametros\tabularnewline
\midrule
\endfirsthead
\toprule
N° & Camadas & Batch & Otimizador & Valor final da função de perda &
Parametros\tabularnewline
\midrule
\endhead
01 & 75 100 100 & 10 & Nadam & 4.4632 E-4 & 4,472,100\tabularnewline
02 & 75 75 75 & 100 & Nadam & 2.2692 E-3 & 3,363,300\tabularnewline
03 & 0 & 100 & Adagrad & 7.6421 E-3 & 176,400\tabularnewline
\bottomrule
\end{longtable}

Fonte: Elaboração própria

O formato das ondas das amostras da primeira peça da bateria, segundos
\emph{round-robins}, são mostrados na coluna à esquerda da
\xrefname{figura}\cref{fig:compensa}, enquanto que as outras colunas
mostram os resultados das redes. As amostras podem ser ouvidas no
repositório preparado para o trabalho (Tesserato, 2018), no caminho
\href{https://github.com/tesserato/tesserato.github.io/tree/master/resources/01_wavs_FF}{``resources/01\_wavs\_FF/''}.

Amostra

Rede 1

Rede 2

Rede 3

Alvo

0.0

1.0

Nas linhas da \xrefname{figura}\cref{fig:compensa} em que não há um
alvo, os resultados foram extrapolados pelas redes. A primeira rede
apresenta um balanço razoável entre acurácia e poder de generalização,
enquanto a segunda, com um número um pouco menor de parametros, introduz
uma quantidade considerável de ruído ao fim das amostras. A terceira,
consideravelmente menor, parece ter aprendido uma representação
intermediária entre os exemplos, que varia apenas ligeiramente de peça
para peça.

\begin{figure}
\centering
\includegraphics{im/compFF.png}
\caption{Resultados: Rede Densa\label{fig:compensa}}
\end{figure}

Fonte: Elaboração própria

Cabe observar que a primeira rede, com pouco mais de 4 milhões de
parâmetros, ocupa-se de cobrir uma família de peças da bateria, e ocupa
em disco aproximadamente 50 MB, implicando que um conjunto completo de
redes girará em torno de 200MB, na medida em que uma rede ocupe-se dos
tambores (incluindo o bumbo), como foi o caso, uma outra da caixa, um
terceira do contratempo, e uma quarta para os pratos. Embora o tamanho
seja razoável, sobretudo à luz dos DMIs comerciais de ponta, como o
Superior Drummer 3 da empresa Toontrack, com 230GB, estratégias podem
ser utilizadas para aumentar a eficiência da rede.

Isso torna-se especialmente importante se lembrarmos que limitamos a
duração das amostras ao tempo de 1 segundo. Para os tambores, caixa e a
maior parte das articulações do contratempo, essa duração é razoável. Os
pratos, contudo, principalmente quando alvos de um ataque mais forte,
costumam apresentar um tempo de decaimento de aproximadamente 10
segundos, o que demandaria uma rede no mínimo 10 vezes maior.

Em outras palavras, essa arquitetura não se presta à generalização da
duração de uma onda sonora, impondo um limite à duração máxima que
poderá ser gerada, e é inaplicável, ao menos de uma modo direto, à
emulação de instrumentos que possuem o som sustentado pela injeção
constante de energia da parte do instrumentista, como é o caso dos
metais, como trombone, saxofone, etc, e instrumentos da família do
violino, onde, respectivamente, o sopro e o arco injetam constantemente
energia no sistema, e o som pode ser mantido por estendidos períodos de
tempo.

\hypertarget{redes-recorrentes-1}{%
\subsection{Redes Recorrentes}\label{redes-recorrentes-1}}

Os resultados obtidos com a utilização de redes densas encorajam a
investigação do comportamento de redes recorrentes, já que possuem o
potencial de diminuir o número de parâmetros, apresentando ainda a
possibilidade de generalização do tempo de duração dos sons gerados. Um
escolha sensível é dividir os alvos em intervalos de 4410 samples, o que
equivale a intervalos de 0.1 segundo para o \emph{sampling rate} de
44100 FPS dos arquivos utilizados, o que equivale aproximadamente à 5
vezes a resolução temporal mínima do ouvido humano (Zwicker e Fastl,
2013). Após alguns testes, a arquitetura recorrente apresentada na
figura abaixo foi escolhida.

\begin{figure}
\centering
\includegraphics[width=4.16667in,height=\textheight]{im/RNNnoOverlap.png}
\caption{Topologia RNN\label{fig:id}}
\end{figure}

Fonte: Elaboração própria

Com 516 210 parametros treináveis (contra 4 472 100 da rede densa), essa
topologia gera resultados semelhantes aos da rede totalmente conectada,
tanto em relação à função de perda quanto à capacidade de generalização,
além de possuir um tamanho em disco aproximadamente 10 vezes menor. A
\xrefname{figura}\cref{fig:comprnn} ilustra as ondas geradas
correspondendo às amostras 454 e 554, e a generalização da rede para uma
amostra entre essas duas, correspondente à um tom entre o tom mais grave
e o bumbo. Esse intervalo foi escolhido justamente por apresentar a
maior disparidade, tanto sonora quanto visual, entre as ondas.

\begin{figure}
\centering
\includegraphics{im/compRNN.png}
\caption{Resultados Rede Recorrente\label{fig:comprnn}}
\end{figure}

Fonte: Elaboração própria

Os testes psicoacústicos revelam, porém, que as pequenas
descontinuidades entre duas previsões subsequentes (não visíveis na
figura acima) traduzem-se em um som de batidas, o que compromete a
qualidade do som gerado pela rede.

\hypertarget{redes-convolucionais-e-convolucionais-recorrentes}{%
\subsection{Redes Convolucionais e
Convolucionais-Recorrentes}\label{redes-convolucionais-e-convolucionais-recorrentes}}

Uma estratégia adotada para contornar o problema, inspirada no
funcionamento das redes convolucionais, foi adotar uma interseção entre
previsões subsequentes da rede; essa região seria ao fim multiplicada
por uma função \emph{smoothstep} decrescente, da forma
\(1 - (3x^2 - 2x^3)\), definida no intervalo entre 0 e 1, a fim de
equacionar de forma suave as constribuições dos dois passos da rede, sem
introduzir novas descontinuidades. A arquitetura proposta é ilustrada
abaixo, e de fato é bastante eficaz em eliminar as descontinuidades
diretas entre as previsões.

\begin{figure}
\centering
\includegraphics{im/smoothstep.png}
\caption{Arquitetura - Convolução Alisada\label{fig:id}}
\end{figure}

Fonte: Elaboração própria

Os quadrados na região superior esquerda da figura representam os
vetores de entrada iniciais fornecidos no primeiro passo. Nos passos
subsequentes, a rede recebe seus próprios vetores de saída, e a partir
deles gera a fração seguinte da onda, de forma recursiva, em uma
configuração denominada \emph{one-to-many} (na prática, como pode ser
observado no código do programa disponibilizado no repositório
{[}tesserato\_2018{]}, encontrando no caminho
\href{https://github.com/tesserato/tesserato.github.io/blob/master/resources/02_RNN.py}{resources/02\_RNN.py},
as entradas recebem o vetor de entrada em cada passo, ou um vetor nulo a
partir do segundo passo, por limitações técnicas. Após algumas épocas,
no entanto, a rede aprende a ignorar essas entradas adicionais).

Todas as frações de onda geradas, com exceção da última, possuem uma
interseção com o início da fração seguinte, e a influência de cada uma
das previsões no resultado final dá-se a partir dos pesos retirados da
função \emph{smoothstep}: O primeiro ponto da região de interseção é
determinado, portanto, inteiramente pela previsão mais antiga enquanto o
último o é totalmente pela previsão mais recente. A influência
majoritária nos pontos internos passa gradativamente, de maneira suave,
da previsão mais antiga para a mais nova.

Contudo, nos resultados alisados persiste ainda um problema relacionado,
principalmente, ao transiente no meio onde propaga-se a perturbação, uma
etapa bastante mal comportada que ocorre logo após o impulso inicial ser
aplicada à corda, membrana ou prato e durante a qual o regime (quasi)
periódico não foi ainda alcançado. Isso resulta, do ponto de vista da
onda, em uma etapa curta, que pode ser vista no início da
\xrefname{figura}\cref{fig:descs}, e que apresenta vibrações aberrantes
em frequências diferentes e maiores do que frequência fundamental do
sistema.

Essas frequências são capturadas pela rede recorrente, e equivocadamente
aplicadas ao início de cada um dos seus passos (e não só ao primeiro). A
\xrefname{figura}\cref{fig:descs}, um detalhe aumentado da onda
correspondente à amostra 111, onde é capturada a fronteira entre duas
previsões subsequentes, ilustra esse efeito na coluna da direita,
enquanto as descontinuidades não tratadas podem ser vistas na coluna da
esquerda.

As consequências sonoras podem ser ouvido no caminho
\href{https://github.com/tesserato/tesserato.github.io/tree/master/resources/02_wavs_RNN}{``resources/02\_wavs\_RNN/''}
do repositório preparado para o trabalho (Tesserato, 2018), nos dois
arquivos .wav que possuem o prefixo ``beating'' e refletem-se em
batidas, mais acentuadas no caso das ondas sem o tratamento.

\begin{figure}
\centering
\includegraphics{im/descontinuidades.png}
\caption{Descontinuidades e Transientes\label{fig:descs}}
\end{figure}

Fonte: Elaboração própria

O problema é tratável, e um dos encaminhamentos possíveis, tomando como
inspiração o algoritmo dos \emph{banded digital waveguides} (Serafin,
Huang e Smith, 2001) é distribuir o esforço de previsão entre mais de
uma rede, por exemplo, ocupando cada uma de prever uma faixa de
frequências. Observou-se, no caso dos samples utilizados até agora, que
aproximadamente 4 faixas seriam necessárias, refletindo-se em 4 redes
por família de peças, em um total de 16 redes. As exatas fronteiras
entre as faixas dependem tanto da arquitetura quanto das características
do instrumento a ser emulado e requerem uma investigação empírica
bastante extensa que deixou de ser desenvolvida neste trabalho, por
limitações de tempo.

Uma outra abordagem seria reservar à uma rede densa a previsão dos
transientes, deixando sob responsabilidade de uma rede recorrente o
período estacionário, mais bem comportado. É ainda possível uma
abordagem experimental em que ao mecanismo de memória da rede sejam
adicionados mecanismos que façam as vezes de filtros digitais
\emph{low-pass}, que tem como resultado amenizar as frequências mais
altas das ondas.

\hypertarget{domuxednio-da-frequuxeancia}{%
\subsection{Domínio da Frequência}\label{domuxednio-da-frequuxeancia}}

Podemos tirar proveito do caráter periódico dos samples e representá-las
diretamente no domínio da frequência. Utilizando a transformada de
Fourier, como visto, temos à disposição uma forma de representar
perfeitamente a onda utilizando um vetor de números complexos com metade
da dimensão da onda original, explorando o fato de que estamos nos
restringindo a representações temporais no domínio dos números reais.

A princípio, tal transformação de domínio não introduziria uma forma
mais compacta de representação das ondas, do ponto de vista
computacional, haja vista que números complexos são representados por
pares de números reais. Levando em conta, contudo, que o ouvido humano
não é capaz de perceber frequências fora da faixa entre 20 Hz e 20 KHz
(Howard e Angus, 2017), identificamos uma das vantagens de trabalharmos
no domínio das frequências: podemos truncar o resultado da FFT à este
intervalo (tomando o cuidado de traduzi-lo em termos das frequências
locais da transformada).

Além disso, trata-se de uma representação independente do tempo de
duração da onda representada, o que no permite trabalhar com uma
arquitetura densa prevendo ondas de tamanhos variados. Mais a frente,
serão ilustradas outras vantagens dessa abordagem, levando em conta
características físicas do instrumento a ser emulado e propriedades da
transformada.

A imagem a seguir compara a distribuição de frequências do som gerado
por um prato de bateria e o som gerado pela Dó central de um piano
(tecla 49 em um piano padrão), e ajuda a entender a diferença entre os
dois instrumentos, e como, mais adiante, poderemos explorar o carater
bem comportado dos sons produzidos por intrumentos harmonicos.

\begin{figure}
\centering
\includegraphics{im/CrashxPiano.png}
\caption{Espectro da intensidade das Frequências Audíveis\label{fig:id}}
\end{figure}

Fonte: Elaboração própria

Para a emulação de um \emph{kit} completo de bateria 4 redes foram
treinadas; as amostras sonoras foram divididas de acordo com as
afinidades de diferentes conjuntos de peças e levando em conta também o
grau de expressividade demandado. As amostras utilizadas no treinamento,
divididas em suas respectivas pastas, podem ser acessadas no repositório
do trabalho (Tesserato, 2018) no caminho
\href{https://github.com/tesserato/tesserato.github.io/tree/master/resources/samples}{``resources/samples/''}.
De maneira semelhante, a implementação da rede está disponível no
endereço
\href{https://github.com/tesserato/tesserato.github.io/blob/master/resources/02_RNN.py}{``resources/02\_RNN.py''}.

Somente para a caixa e tons foram utilizados \emph{round-robins}, devido
à alta propensão ao efeito \emph{machine-gun} quando as peças são
utilizadas em rápida sucessão, articulação comum em vários estilos
musicais.

A arquitetura utilizada é como apresentada na figura
\xrefname{figura}\cref{fig:arqd} e é uniforme para todas as peças;
variações no número de parâmetros ocorrem em função da diferença de
duração das ondas no caso dos pratos. Para o caso da caixa, contratempo
e dos tambores, incluído o bumbo, cada uma das redes é compostapor 120
008 parâmetros treináveis, incluindo pesos das conexões e biases, com um
tamanho em disco de aproximadamente 1,4 Mb.

Os pratos, por possuírem um som com maior tempo de sustentação,
demandaram uma rede com saídas maiores, totalizando 360 008 parametros
treináveis e quase o triplo do tamanho em disco. Temos assim que todo o
sistema possui em torno de 9 Mb de tamanho em disco. Foram utilizadas
funções de ativação lineares, tanto para evitar \emph{overfitting}
quanto para permitir uma implementação menos intensiva, do ponto de
vista computacional.

\begin{figure}
\centering
\includegraphics[width=4.16667in,height=\textheight]{im/Modelo-Bateria-Caixa.png}
\caption{Arquitetura para a emulação da Caixa de uma
Bateria\label{fig:arqd}}
\end{figure}

Fonte: Elaboração própria

Os resultados são bastante satisfatórios, tanto em termos de acurácia
quanto de generalização. A rede não decora nenhuma das distribuições,
como é esperado, mas aprende a generaliza-las no espaço bidimensional
que tem como eixos o tipo da peça, que reflete a grosso modo uma
variação entre agudo e grave e a intensidade da batida. A
\xrefname{figura}\cref{fig:hihat} ilustra uma onda onde foram
concatenadas as saídas da rede para o contratempo aberto, golpeado com
intensidade crescente. A onda concatenada pode ser ouvida no caminho
\href{https://github.com/tesserato/tesserato.github.io/blob/master/resources/04_freq_domain/Exemplos\%20das\%20Predi\%C3\%A7\%C3\%B5es\%20-\%20Bateria/contratempo_piano_para_forte.wav}{``resources/04\_freq\_domain/Exemplos
das Predições - Bateria/contratempo\_piano\_para\_forte.wav''} do
repositório do trabalho, enquanto as amostras originais encontram-se em
\href{https://github.com/tesserato/tesserato.github.io/tree/master/resources/04_freq_domain/Exemplos\%20das\%20Predi\%C3\%A7\%C3\%B5es\%20-\%20Bateria/Intensidade\%20-\%20Contratempo\%20Aberto}{"resources/04\_freq\_domain/Exemplos
das Predições - Bateria/Intensidade - Contratempo Aberto}

\begin{figure}
\centering
\includegraphics[width=\textwidth,height=3.125in]{resources/04_freq_domain/Exemplos das Predições - Bateria/contratempo_piano_para_forte.png}
\caption{Dinâmicas geradas para Contratempo Aberto - Pianíssimo até
Molto Forte\label{fig:hihat}}
\end{figure}

Fonte: Elaboração própria

No casa das redes treinadas com o uso de \emph{round-robins} pode-se
imaginar um terceiro eixo no qual a rede generaliza variações percebidas
nas amostras utilizados para o treinamento. A rigor, o efeito é
determinístico, totalmente dependente dos inputs da rede. Pode-se pensar
nele, contudo, como um termo que adiciona certa aleatoriedade aos sons
gerados, na medida em que pode-se alimentar a rede com um número
arbitrário entre -1 e 1 para para cada amostra gerada, introduzindo
variações que emulam a aleatoriedade natural do instrumento. Essas
variações são ilustradas na \xrefname{figura}\cref{fig:rrobins}, para o
caso do tom mais agudo. As ondas geradas pela rede podem ser encontradas
em
\href{https://github.com/tesserato/tesserato.github.io/tree/master/resources/04_freq_domain/Exemplos\%20das\%20Predi\%C3\%A7\%C3\%B5es\%20-\%20Bateria/Round-Robins\%20-\%20Tom\%2001}{``resources/04\_freq\_domain/Exemplos
das Predições - Bateria/Round-Robins - Tom 01''} e a versão concatenada,
exibida na \xrefname{figura}\cref{fig:rrobins}, para maior comodidade,
encontra-se em
\href{https://github.com/tesserato/tesserato.github.io/blob/master/resources/04_freq_domain/Exemplos\%20das\%20Predi\%C3\%A7\%C3\%B5es\%20-\%20Bateria/tom_roundrobins.wav}{``resources/04\_freq\_domain/Exemplos
das Predições - Bateria/tom\_roundrobins.wav''}

\begin{figure}
\centering
\includegraphics[width=\textwidth,height=3.125in]{resources/04_freq_domain/Exemplos das Predições - Bateria/tom_roundrobins.png}
\caption{7 \emph{Round-Robins} gerados para o Tom 1\label{fig:rrobins}}
\end{figure}

Fonte: Elaboração própria

Um efeito colateral desse comportamento geral da rede é que ela não
reproduz exatamente nenhuma das amostras utilizados no treinamento, e
não há garantias do grau de semelhança entre os timbres gerados e os
exemplos utilizados. Empiricamente, no entanto, observou-se que as
características do espectro de frequências aprendidas e generalizadas
pela rede imprimem ao resultado um timbre totalmente verossímil; em
outras palavras, os tons soam como tons, não exatamente com o mesmo
timbre que cada uma das peças utilizadas no treinamento, mas como uma
amalgama entre elas. Assim é com os pratos, e sucessivamente para todas
as peças.

Quando o esforço é o de emular duas dinâmicas diferentes de uma mesma
peça, como é o caso da caixa, em que uma das variáveis é, como para
todas as outras peças, a intensidade, e a outra é a centralidade da
batida (mais próxima ou afastada do centro da caixa), observa-se
resultados ainda mais orgânicos, como pode ser visto na
\xrefname{figura}\cref{fig:cent}. As ondas geradas podem ser ouvidas em
\href{https://github.com/tesserato/tesserato.github.io/tree/master/resources/04_freq_domain/Exemplos\%20das\%20Predi\%C3\%A7\%C3\%B5es\%20-\%20Bateria/Centralidade\%20-\%20Caixa}{``resources/04\_freq\_domain/Exemplos
das Predições - Bateria/Centralidade - Caixa/''}, enquanto que a versão
concatenada, como paresentada na \xrefname{figura}\cref{fig:cent} está
disponível em
\href{https://github.com/tesserato/tesserato.github.io/blob/master/resources/04_freq_domain/Exemplos\%20das\%20Predi\%C3\%A7\%C3\%B5es\%20-\%20Bateria/caixa_centro_para_bordas.wav}{``resources/04\_freq\_domain/Exemplos
das Predições - Bateria/caixa\_centro\_para\_bordas.wav''}.

\begin{figure}
\centering
\includegraphics[width=\textwidth,height=3.125in]{resources/04_freq_domain/Exemplos das Predições - Bateria/caixa_centro_para_bordas.png}
\caption{Dinâmica posição do golpe para caixa\label{fig:cent}}
\end{figure}

Fonte: Elaboração própria

Para sons com características harmonicas, apresentando frequências
parciais com caráter aproximadamente periódico, esse método não oferece
bons resultados, na medida em que incorpora as características gerais de
cada um dos sons utilizados no treinamento. É curioso observar que, por
exemplo, treinando a rede para aprender a interjeição `ah' cantada em
diferentes notas, o resultado tem uma qualidade parecida com um coral.

As transformadas são exibidas na \xrefname{figura}\cref{fig:voz}:, tanto
dos sons originais quanto das previsões, para a nota mais baixa e mais
alta. É possível reparar que a previsão encaminha-se para os picos em
ambos os casos, mas incorporando uma grande quantidade de dados
indesejáveis. Os áudios são encontrados no repositório do trabalho, no
caminho
\href{https://github.com/tesserato/tesserato.github.io/tree/master/resources/04_freq_domain/Experimento\%20-\%20Voz}{``resources/04\_freq\_domain/Experimento
- Voz/''}

Amostra

Original

Previsão

1-1-1

1-19-1

\begin{figure}
\centering
\includegraphics{im/voice.png}
\caption{Experimento Emulação da Voz Cantada\label{fig:voz}}
\end{figure}

Fonte: Elaboração própria

\hypertarget{modelo-misto}{%
\subsection{Modelo Misto}\label{modelo-misto}}

Pudemos notar que instrumentos de caráter harmonico possuem um
distribuição de frequências bem comportada, consistindo basicamente de
alguns picos em sua representação no domínio da frequência. Ademais,
para que seu som seja considerado harmonioso, esses picos distribuem-se,
forçosamente, de maneira ordenada; idealmente, são múltiplos inteiros da
frequência fundamental da nota reproduzida.

É importante notar que a largura da `base' desses padrões de picos no
domínio da frequência é proporcional ao decaimento de uma senoide pura:
no caso extremo, uma senoide perfeitamente periódica possui uma
distribuição de frequências zero em todos os pontos, com excessão do
ponto que representa sua frequência. A figura a
\xrefname{figura}\cref{fig:decay} ilustra essa relação.

\begin{figure}
\centering
\includegraphics{im/fourierDecay.png}
\caption{Efeito do Decaimento sobre a Tranformada de Fourier: Senoide
440hz\label{fig:decay}}
\end{figure}

Fonte: Elaboração própria

Podemos observar que o decaimento introduz novas frequências, ao redor
da frequência principal, e mudanças de fase na representação no domínio
da frequência; observa-se empiricamente que o `objetivo' primordial
dessas novas frequencias e fases é reproduzir o efeito de decaimento
(ou, de forma mais ampla, o envelope) da onda.

Dessa observação, duas intuições importantes podem ser retiradas: A
primeira é que, com um grau razoavel de aproximação, podemos descrever
um som hamonico gerado por uma excitação impulsiva em função da
localização de algumas de suas frequências, suas respectivas
intensidades e os seus decaimentos.

É natural supor também que a influência das fases dessas ondas pode ser
ignorada, o que comprova-se empiricamente a partir da resconstrução de
ondas com as fases originais e fases nulas ou aleatórias. No caminho
\href{https://github.com/tesserato/tesserato.github.io/tree/master/resources/05\%20Modelo\%20Final/03_waves_from_01_info/piano}{``resources/05
Modelo Final/03\_waves\_from\_01\_info/piano/''} do repositório podem
ser encontradas ondas reconstruídas, descartando-se informações sobre as
fases de cada uma das 100 frequências parciais, a partir das amostras
utilizadas para o treinamento da rede, disponíveis em
\href{https://github.com/tesserato/tesserato.github.io/tree/master/resources/05\%20Modelo\%20Final/00_samples/piano}{resources/05
Modelo Final/00\_samples/piano/}. Comparando as amostras, pode-se
perceber que a reprodução é bastante verossímil. A maior parte da
diferença perceptual tem origem na quantidade de parciais consideradas,
que não contempla as frequências presentes na fase transiente da onda,
nas teclas mais graves.

Amostra

Original

Reconstrução

Tecla 01

Tecla 35

Tecla 49

\begin{figure}
\centering
\includegraphics{im/fourierDecay-comp.png}
\caption{Onda composta por senoides a 110, 220 e 330hz\label{fig:id}}
\end{figure}

Fonte: Elaboração própria

Podemos, então, descartar as informações sobre a fase específica de cada
parcial da onda, restando desenvolver um método para extração das
frequências principais de uma onda sonora, e seus respectivos envelopes,
que aqui assumiremos, simplificadamente, da forma exponencial, sendo
caracterizados por uma taxa de decaimento e uma amplitude inicial.

Se lembrarmos que a transformada de Fourier oferece a intensidade de
cada uma das frequências que compõe a onda (ou mesmo a amplitude, se
normalizada de forma adequada) pode-se supor que o levantamento das duas
primeiras informações sobre a onda torna-se trivial.

Algumas dificuldades, contudo, logo apresentam-se: elencando as
frequencias de maior intensidade de uma onda arbitrária, rapidamente
notamos que muitas delas fazem parte de um mesmo pico. Além disso, as
amplitudes referidas na transformada são as amplitudes médias da
frequência em toda a onda que, além de não serem diretamente
utilizáveis, não oferecem informação sobre o comportamento delas no
tempo.

\hypertarget{frequuxeancias}{%
\subsubsection{Frequências}\label{frequuxeancias}}

Atacando a questão das frequências, é conveniente lembrar de que, em uma
onda pefeitamente harmonica, temos uma frequência fundamental \(f_0\),
da qual todas as outras frequências relevantes são múltiplos inteiros,
da forma \(f_n = n f_0, n \in \mathbb{Z}^+\). Além disso, a frequência
fundamental é, geralmente, a de maior intensidade (mas nem sempre, como
ilustrado na \xrefname{figura}\cref{fig:freq}). Podemos fazer uso desta
informação para encontrar os picos em uma onda.

No caso das teclas de um piano, a relação entre as 88 teclas e sua
frequência fundamental pode ser encontrada a priori a partir da fórmula
\(f_0 = (tecla^2 - 49) / 12 * 440\). Podemos, então, buscar os máximos
nos intervalos apropriados. A figura a seguir compara esse algoritmo com
o simples levantamento dos maiores valores de intensidade, para uma onda
correspondete ao som emitido pela tecla 35 de um piano, da qual busca-se
os 30 primeiros picos, correspondendo às 30 primeiras frequências
parciais.

\begin{figure}
\centering
\includegraphics{im/piano_peaks_tecla_35.png}
\caption{Picos e Valores Máximos - Piano, tecla 35\label{fig:freq}}
\end{figure}

Fonte: Elaboração própria

Repara-se que simplesmente elencando os maiores valores, todas as
frequências identificadas gravitam em torno as 3 primeiras parciais. Com
o algoritmo de busca de picos, as verdadeiras parciais da onda são
capturadas, em intervalos bastante próximos à multiplos inteiros de
\(f_0\). É interessante observar que a onda ilustrada em
\xrefname{figura}\cref{fig:freq} representa um caso onde a frequência
fundamental não corresponde a elevação mais alta.

\hypertarget{decaimentos-e-amplitudes}{%
\subsubsection{Decaimentos e
Amplitudes}\label{decaimentos-e-amplitudes}}

Uma abordagem possível para a investigação das frequências parciais
identificadas seria a aplicação de um método iterativo, seja a partir de
um algoritmo de \emph{curve fitting}, seja com a utilização de redes
neurais. Ambos os métodos foram investigados e limitações críticas foram
identificadas: O algoritmo iterativo usado para \emph{curve fitting} não
apresenta garantias de convergência.

Além disso, o processo iterativo deve ser utilizado para cada uma das
ondas, e para cada uma das frequências parciais que deseja-se extrair de
cada uma delas. Mesmo tratando-se de uma etapa prévia, que não fará
parte do modelo em tempo real, os custos computacionais mostraram-se
proibitivos (sobretudo para ondas com vários segundos de duração).

De maneira semelhante, o método baseado em redes neurais, embora não
deixe de convergir, não apresenta garantias sobre a qualidade dos
resultados. Ademais, os dois algoritmos são míopes, no sentido de que,
na busca por minimizar suas funções objetivo, geram ,por vezes, senoides
com aplitudes e decaimentos aberrantes.

Considerando, de forma aproximada, que todos os decaimentos são
exponenciais, é possível estimá-los observando a diferença de
intensidade da frequência de interesse entre intervalos arbitrários,
como por exemplo a primeira e a segunda metades, da onda.

Tomando o cuidado de converter as frequências absolutas para as
frequências locais, o somatório
\(z(f) = \sum_{t-l/2}^{t+l/2}s[t] ~ e^{-2 \pi t f i } dt\) nos fornece
essa informação: trata-se do caso discreto da integral da multiplicação
ponto a ponto do sinal de interesse \(s(t)\) com o \emph{kernel} da
transformada de Fourier em função unicamente do tempo \(t\), com \(f\)
fixo na frequência local para a qual deseja-se a amplitude. \(z(f)\)
fornece também a fase média no intervalo, que não nos é útil, como foi
visto. Poderíamos utilizar um número arbitrário de intervalos; dois, no
entanto, provaram-se suficientes empiricamente.

O decaimento pode ser estimado, então, a partir da relação
\(d = 2 \ln(a_1 / a_2) / l\). No caso específico de dois intervalos,
temos que a amplitude \(a\) é
\(a_1 / e^{-d l/ 4} = a_2 / e^{-d 3l/ 4}\). Na
\xrefname{figura}\cref{fig:ddcay} são apresentadas as estimativas para
as frequências mais presentes nos sons emitidos por algumas teclas de
piano, além de um exemplo demonstrativo com um senoide puro.

\begin{figure}
\centering
\includegraphics{im/frequencydecay.png}
\caption{Decaimentos Estimados\label{fig:ddcay}}
\end{figure}

Fonte: Elaboração própria

Cabe observar que nos casos dos sons do piano ilustrados nos 5 primeiros
retangulos da imagem o decaimento estimado é para uma das muitas
frequências que compõe a onda - especificamente a de maior intensidade.
No último retangulo, o decaimento coincide com o envelope da onda, já
que esta é composta por apenas uma frequência, e podemos ter uma ideia
melhor de sua acurácia. Estamos, desta forma, equipados para introduzir
as redes neurais no modelo, que será usada para aprender a estimar e
generalizar os parametros extraídos das amostras.

\hypertarget{formulauxe7uxe3o-final}{%
\subsubsection{Formulação Final}\label{formulauxe7uxe3o-final}}

Podemos escrever as frequências, tanto a fundamental quanto as parciais,
em função das teclas de um piano, da seguinte forma:
\(f(k,n)=2^{(k-49)/12}440(n+1), k \in {1, 2,...,88}, n \in {0,1,2,...}\).
Esse \emph{framework} é conveniente para uma vasta gama de instrumentos,
já que essas 88 teclas vão desde a nota A0 até C8, cobrindo com folga o
espectro de frequências da maioria dos instrumentos de interesse.

Para treinar a rede a partir de qualquer instrumento, basta rotular as
amostras dos sons pertinentes com o número da tecla equivalente de um
piano. Como pode ser ouvido nos exemplos apresentados ao fim deste
trabalho, esse procedimento permitiu o treinamento de um instrumento
híbrido, utilizando amostras de um baixo acústico, para o registro mais
grave, um violoncello para o meio do registro, e um violino para o mais
agudo; todos esses instrumentos, ao mesmo tempo, cobriram até a tecla 75
do piano, somente.

Como foi visto, essa equação desconsidera as inarmonicidades presentes
nos instrumentos, responsáveis por características importantes de seus
timbres. Contudo, apresenta uma aproximação inicial bastante razoável
que serve tanto para reforçar as características harmonicas básicas no
modelo final quanto para aliviar o esforço de predição da rede, na
medida em que podemos adicionar um termo de inarmonicidade na equação
acima, a ser aprendido pela rede.

Temos, assim:
\(f(k,n)=2^{(k-49)/12}440(n+1)i(k,n), k \in {0,1,...88}, n \in {0,1,2,...}\).
Considerando tanto frequências fundamentais quanto parciais no espectro
audível, teríamos um intervalo de 27 hz, a frequência fundamental de A0,
até um pouco mais de 20Khz, correspondendo, por exemplo, à quinta
parcial de C8, caso tentássemos estimar diretamente as frequências
parciais de um piano.

Escolhendo trabalhar com as inarmonicidades, por outro lado, reduzimos o
intervalo para o limite entre 0 e 3. Além disso, o comportamento das
inarmonicidades é razoavelmente previsível, com caráter ligeiramente
exponencial, como ilustrado a na \xrefname{figura}\cref{fig:inh}, para a
primeira tecla de um piano, e a tecla 34, a última a ter todas as
parciais no espectro audível.

\begin{figure}
\centering
\includegraphics{im/inharmonicities.png}
\caption{inarmonicidades\label{fig:inh}}
\end{figure}

Fonte: Elaboração própria

Este mesmo rationale é aplicado aos decaimentos e amplitudes, que são
previstos como uma fração dos valores máximos encontrados em cada uma
das teclas (ou notas, mais genericamente, no caso de estarmos treinando
um instrumento arbitrário), e encontram-se no intervalo fechado entre 0
e 1.

O comportamento das amplitudes e decaimentos máximos, por tecla, é
apresentado na \xrefname{figura}\cref{fig:dkey}, e poderia ser, também,
modelado, tanto a partir de redes neurais quanto por meio de modelos
matemáticos. Empiricamente, contudo, nota-se que os valores médios são
suficientes para um bom resultado.

Isso é oportuno, pois reduz a complexidade do modelo final, além de
conferir mais generalidade: na \xrefname{figura}\cref{fig:dkey} abaixo
notamos, por exemplo, que há uma descontinuidade nos valores dos
decaimentos ocorrendo a partir da tecla 55, que deve-se provavelmente à
mudança de cordas duplas para cordas simples.

\begin{figure}
\centering
\includegraphics{im/piano_max_decays.png}
\caption{Decaimentos Máximos por Notas\label{fig:dkey}}
\end{figure}

Fonte: Elaboração própria

Esse fenomeno é partircular dos pianos, e, além disso, essa
descontinuidade é arbitrária, e varia de piano para piano. Teríamos,
então, caso não fizessemos uso da média, que formular uma aproximação
para cada instrumento treinado (ou treinar uma rede adicional). As
amplitudes, embora pareçam possuir uma tendêncial geral de diminuição na
medida em que avançamos na escala do piano, possuem comportamento
bastante arbritrário.

\begin{figure}
\centering
\includegraphics{im/piano_max_amps.png}
\caption{Amplitudes Máximas por Notas\label{fig:id}}
\end{figure}

Fonte: Elaboração própria

Tendo sido observado que, de posse de uma estimativa do decaimento das
frequências parciais, as fases de cada uma delas tem pouco impacto na
reconstrução da onda, optou-se por randomizar as fases: Essa decisão
apresenta a vantagem de caonferir um resultado mais orgânico e variado à
sintese executada pelo modelo final, na medida em que nenhuma onda
gerada será exatamente igual á qualquer uma outra. Outra opção, menos
rica do ponto de vista percpetual, seria ainda dar às fases uma valor
arbitrário (como zero, por exemplo).

Para a previsão de cada uma das 3 quantidades explicitadas, uma rede
densa, com 3 camadas foi utilizada. A arquitetura das 3 redes é
identica, com exceção do número de neurons em cada uma das camadas
ocultas.

Optou-se por utilizar como função de ativação uma versão modificada da
tangente hiperbólica, na forma \(a(x) = \tanh(6 * x - 3) / 2 + 1 / 2\),
de forma a melhor cobrir o intervalo \([0,1]\). Além disso, o método de
inicialização de pesos proposto por LeCun \emph{et al.} (2012) Klambauer
\emph{et al.} (2017) ofereceu uma melhora considerável no tempo de
convergência das redes, e foi a inicialização utilizado na versão final
do modelo, tanto para os \emph{neurons} quanto para os \emph{biases}.

As redes responsáveis pelas amplitudes e decaimentos possuem 50 neurons
em cada uma das camadas ocultas, e um total de 5 301 parametros
treináveis, com aproximadamente 100kb de espaço em disco. Sua
arquitetura é apresentada abaixo. Para as frequências, o número de
neurons foi reduzido para 10, para evitar \emph{overfitting}, o que gera
um total de 261 parametros treináveis e aproximadamente metade do
tamanho em disco. A \xrefname{figura}\cref{fig:fin} a seguir compara os
dados originais com os produzidos pelo modelo.

\begin{figure}
\centering
\includegraphics{im/finalnet.png}
\caption{Previsões - Modelo Misto\label{fig:fin}}
\end{figure}

Fonte: Elaboração própria

Na pasta
\href{https://github.com/tesserato/tesserato.github.io/tree/master/resources/05\%20Modelo\%20Final/02_predictions/piano}{resources/05
Modelo Final/02\_predictions/piano/} do repositório preparado para este
trabalho podem ser observadas imagens semelhantes, para outros números
de neurons em todas as camadas ocultas, incluindo alguns exemplos em que
ocorre \emph{overfitting}.

Podemos, assim, definir uma nova metodologia voltada à modelagem sonora
de instrumentos harmonicos, batizada de \textbf{Neuro-spectral
Synthesis}, resumida de maneira esquemática na
\xrefname{figura}\cref{fig:neurospectral}.

\begin{figure}
\centering
\includegraphics{im/neurospectral.png}
\caption{Modelo final\label{fig:neurospectral}}
\end{figure}

\hypertarget{resultados}{%
\section{Resultados}\label{resultados}}

O modelo final apresenta resultados mais realistas do que as
implementações baseadas no modelo de \emph{Digital Waveguides} e no
método das diferenças finitas, e é ao menos uma ordem de grandeza mais
eficiente, nas implementações aqui apresentadas.

Para efeito de comparação, levando em conta a geração de 5 ondas de 1
segundo, a 44100 FPS, o método proposto foi, em média, 17 vezes mais
rápido que o método baseado em Digital Waveguides e 26 vezes mais rápido
do que o método das diferenças finitas. A
\xrefname{tabela}\cref{tbl:compar} abaixo elenca os resultados.

\begin{longtable}[]{@{}lccc@{}}
\caption{Comparação Tempo de Geração(s)
\label{tbl:compar}}\tabularnewline
\toprule
& Finite Differences & Digital Waveguides & Neuro-Spectral
Synthesis\tabularnewline
\midrule
\endfirsthead
\toprule
& Finite Differences & Digital Waveguides & Neuro-Spectral
Synthesis\tabularnewline
\midrule
\endhead
1 & 16.95213819 & 13.81114888 & 0.5896253586\tabularnewline
2 & 19.62042928 & 13.02065825 & 0.8354649544\tabularnewline
3 & 24.25646234 & 11.33273816 & 0.7555158138\tabularnewline
4 & 17.89053726 & 12.3500874 & 0.7475204468\tabularnewline
5 & 18.00246739 & 12.12223339 & 0.7735056877\tabularnewline
\textbf{Média} & \textbf{19.34440689} & \textbf{12.52737322} &
\textbf{0.7403264523}\tabularnewline
\textbf{Desvio Padrão} & \textbf{2.908692366} & \textbf{0.937300851} &
\textbf{0.09102950916}\tabularnewline
\bottomrule
\end{longtable}

Fonte: Elaboração própria

Treinadas as redes, uma implementação final do modelo é bastante
simples, e é apresentada no quadro 4.

\begin{Shaded}
\begin{Highlighting}[numbers=left,,]
\NormalTok{key }\OperatorTok{=} \DecValTok{66}
\NormalTok{partials }\OperatorTok{=} \DecValTok{100}
\NormalTok{n }\OperatorTok{=} \DecValTok{44100}
\NormalTok{fps }\OperatorTok{=} \DecValTok{44100}

\NormalTok{P }\OperatorTok{=}\NormalTok{ np.arange(}\DecValTok{1}\NormalTok{, partials }\OperatorTok{+} \DecValTok{1}\NormalTok{, }\DecValTok{1}\NormalTok{)}

\NormalTok{[ma, Ma] }\OperatorTok{=}\NormalTok{ np.genfromtxt(}\StringTok{'a_m_M.csv'}\NormalTok{, delimiter}\OperatorTok{=}\StringTok{','}\NormalTok{)}
\NormalTok{[md, Md] }\OperatorTok{=}\NormalTok{ np.genfromtxt(}\StringTok{'d_m_M.csv'}\NormalTok{, delimiter}\OperatorTok{=}\StringTok{','}\NormalTok{)}
\NormalTok{[mf, Mf] }\OperatorTok{=}\NormalTok{ np.genfromtxt(}\StringTok{'f_m_M.csv'}\NormalTok{, delimiter}\OperatorTok{=}\StringTok{','}\NormalTok{)}

\NormalTok{a_net }\OperatorTok{=}\NormalTok{ K.models.load_model(}\StringTok{'a_model.H5'}\NormalTok{, \{}\StringTok{'scaled_tanh'}\NormalTok{: scaled_tanh\})}
\NormalTok{d_net }\OperatorTok{=}\NormalTok{ K.models.load_model(}\StringTok{'d_model.H5'}\NormalTok{, \{}\StringTok{'scaled_tanh'}\NormalTok{: scaled_tanh\})}
\NormalTok{f_net }\OperatorTok{=}\NormalTok{ K.models.load_model(}\StringTok{'f_model.H5'}\NormalTok{, \{}\StringTok{'scaled_tanh'}\NormalTok{: scaled_tanh\})}

\NormalTok{f0 }\OperatorTok{=}\NormalTok{ np.power(}\DecValTok{2}\NormalTok{, (key }\OperatorTok{-} \DecValTok{49}\NormalTok{) }\OperatorTok{/} \DecValTok{12}\NormalTok{) }\OperatorTok{*} \DecValTok{440}
\NormalTok{norm_key }\OperatorTok{=}\NormalTok{ (key }\OperatorTok{-} \DecValTok{1}\NormalTok{) }\OperatorTok{/} \DecValTok{87}
\NormalTok{ipt }\OperatorTok{=}\NormalTok{ np.array([[norm_key, p] }\ControlFlowTok{for}\NormalTok{ p }\KeywordTok{in}\NormalTok{ np.linspace(}\DecValTok{0}\NormalTok{, }\DecValTok{1}\NormalTok{, partials)])}

\NormalTok{A }\OperatorTok{=}\NormalTok{ np.ndarray.flatten(a_net.predict(ipt)) }\OperatorTok{*}\NormalTok{ ((Ma }\OperatorTok{-}\NormalTok{ ma) }\OperatorTok{+}\NormalTok{ ma) }\OperatorTok{*} \DecValTok{5000}
\NormalTok{D }\OperatorTok{=}\NormalTok{ np.ndarray.flatten(d_net.predict(ipt)) }\OperatorTok{*}\NormalTok{ ((Md }\OperatorTok{-}\NormalTok{ md) }\OperatorTok{+}\NormalTok{ md) }\OperatorTok{*} \FloatTok{0.000113026536390889} \CommentTok{#average decay}
\NormalTok{F }\OperatorTok{=}\NormalTok{ (np.ndarray.flatten(f_net.predict(ipt)) }\OperatorTok{*}\NormalTok{ (Mf }\OperatorTok{-}\NormalTok{ mf) }\OperatorTok{+}\NormalTok{ mf) }\OperatorTok{*}\NormalTok{ P }\OperatorTok{*}\NormalTok{ f0}

\NormalTok{w }\OperatorTok{=}\NormalTok{ np.zeros(n)}
\ControlFlowTok{for}\NormalTok{ i }\KeywordTok{in} \BuiltInTok{range}\NormalTok{(partials):}
  \BuiltInTok{print}\NormalTok{(}\StringTok{'partial:'}\NormalTok{, i)}
\NormalTok{  ph }\OperatorTok{=}\NormalTok{ np.random.uniform(}\DecValTok{0}\NormalTok{, }\DecValTok{2} \OperatorTok{*}\NormalTok{ np.pi)}
\NormalTok{  w }\OperatorTok{+=}\NormalTok{ create_signal(n, fps, F[i], ph, A[i], D[i])}
\end{Highlighting}
\end{Shaded}

Implementação na linguagem Python do modelo final

Fonte: Elaboração própria

Uma implementação como essa, funcional, mais os arquivos de definição da
arquitetura das redes e seus respectivos pesos estão disponíveis no
repositório preparado para este trabalho, em
\href{https://github.com/tesserato/tesserato.github.io/tree/master/resources/06\%20Tempo\%20Real}{resources/06
Tempo Real/}, e ocupam aproximadamente 250kb de armazenamento em disco.
No repositório podem ser encontradas alguns áudios gerados pelo modelo.

A principal limitação da metodologia proposta consiste no fato de que
todos os parametros a serem manipulados no modelo final devem, antes,
ter sido incorporados ao processo de treinamento. Os dois paradigmas de
modelagem física utilizados para comparação, tanto o método das
diferenças finitas quanto os \emph{digital waveguides}, permitem alguma
manipulação em tempo real, a posteriori, de seus parametros: nos
exemplos apresentados, o ponto de excitação pode ser mudado de onda para
onda e, a qualquer tempo, o ponto de captação pode ser alterado, até
mesmo durante a simulação, refletindo no timbre do som gerado.

Além disso, esses modelos prestam-se à incorporação razoavelmente
trivial de uma fonte, contínua ou periódica, de excitação, e podem ser
utilizados para a simulação de instrumentos de som continuo, como
violinos acionados pelo arco (ao contrario do acionamento
\emph{pizzicato} aqui apresentado) e metais, por exemplo; esse não é o
caso do modelo aqui proposto, e constitui uma possibilidade de
desenvolvimento futuro interessante.

Uma outra barreira é que o modelo aprende características sonoras a
partir de exemplos e não presta-se, de forma prática, à exploração
sonora direta; esta deve ser efetuada a partir de outra ferramenta, e
então incorporada ao modelo, à época de seu treinamento.

Para ilustração, algumas músicas criadas com os sons gerados pelos
modelos propostos, treinados para emular um kit regular de bateria, um
piano e um híbrido entre baixo acústico, cello e violino podem ser
ouvidos no endereço
\href{https://soundcloud.com/carlos-tarjano/sets/spectral-neural-synthesis}{soundcloud.com/carlos-tarjano/sets/spectral-neural-synthesis}

Do ponto de vista conceitual, o trabalho apresenta alguns
\emph{insights} sobre a transformada discreta de Fourier que, até onde
alcança o conhecimento do autor, não foram anteriormente explorados na
literatura: A Interpretação geométrica da simetria da transformada
discreta de Fourier aplicada à sinais no domínio real, embora não se
preste à uma aplicação direta, ajuda no entendimento deste algoritmo e
pode vir a ser explorada em desenvolvimentos futuros.

A relação entre o envelope de uma onda, no domínio do tempo, e o formato
de seus picos no domínio da frequência, por outro lado, possibilitou a
formulação do modelo fisicamente informado, e pode ser explorada
futuramente de maneira mais profunda, potencialmente baseando um método
mais acurado de identificação de envelopes em ondas sonoras.

\hypertarget{conclusuxe3o}{%
\section{Conclusão}\label{conclusuxe3o}}

O presente trabalho, ao introduzir uma nova técnica de modelagem sonora,
demonstra o potencial do uso de redes neurais à síntese de áudio,
provando sobretudo a possibilidade desse uso em tempo real. O segundo
modelo introduzido apresenta resultados mais verossímeis do que o
algoritmo de modelagem acústica em tempo real mais utilizado e eficiente
encontrado na literatura, à um custo computacional bastante inferior.

Uma outra vantagem frente aos modelos convencionais baseados em
simulação física é que ele pode aprender características importantes do
som que tem origem em partes do sistema difíceis de serem modeladas
fisicamente, como a influência de ressonadores de geometria complexa,
por exemplo.

O trabalho evidencia que arquiteturas densas e recorrentes são capazes
de aprender representações internas que possibilitem a reprodução e
generalização de amostras sonoras de maneira direta; a partir da
introdução de uma representação compacta, fisicamente informada, de
ondas sonoras harmonicas, o trabalho evidencia o potencial sinergético
entre os desenvolvimentos na pesquisa sobre a acústica de intrumentos
musicais e a utilização de redes neurais para basear modelos destinados
à emulação desses instrumentos, ou famílias de instrumentos,
específicos.

Além disso, através da utilização de funções de ativação especialmente
elaboradas para acomodar os parametros da representação pertinentes e
métodos de inicialização de pesos e viezes apropriados, essas
aquiteturas pode ser simplificadas e o número necessário de parametros
treináveis sensivelmente diminuído, tornando-as mais efetivas do ponto
de vista da síntese sonora em tempo real.

Com a escolha apropriada da representação no domínio da frequência, os
modelos prestam-se, com maior ou menor grau de eficiência, à emulação de
uma vasta gama de intrumentos, respectivamente percussivos e harmonicos,
dentro do escopo proposto para o trabalho; um kit de bateria, como o
treinado durante o trabalho, representa por si só grande parte da
família de instrumentos percussivos, em suas diferentes peças.

De maneira semelhante, o treinamento de um piano e um instrumento
híbrido, a partir do modelo fisicamente informado apresentado
representam de maneira bastante completa as aplicações em relação à
instrumentos acústicos harmonicos de caráter impulsivo.

\hypertarget{desenvolvimentos-futuros}{%
\subsection{Desenvolvimentos Futuros}\label{desenvolvimentos-futuros}}

As possibilidades de desenvolvimentos futuros nesta área de intersecção
entre redes neurais e acústica são inúmeras, haja vista, inclusive, a
escassez de investigações semelhantes: Seria interessante, por exemplo,
utilizar as saídas de um modelo elaborado a partir do método das
diferenças finitas, que pode ser formulado de forma a simular
características mais sofisticadas de um instrumento como, como rigidez,
ressonância e vários termos de perda de um dado sistema acústico, ao
custo de uma alta demanda de recursos computacionais, para treinar um
modelo baseado em \emph{digital waveguides} que tenha uma rede neural no
ponto onde as perdas e demais cálculos são concatenados.

Devido ao alto grau de recursividade do algoritmo \emph{digital
waveguides}, o treinamento à partir do resultado final esperado para o
modelo é bastante complexo de ser implementado; os vetores de saída de
uma simulação baseada em diferenças finitas, no entanto, são plenamente
compatíveis com este treinamento, e a inserção de uma rede neural
poderia levar a um modelo que retenha ao menos parte da acurácia da
simulação pelo método de diferenças finitas, com eficiência
computacional próxima, ou até superior, à apresentada pelo algoritmo de
digital waveguides. Em Gully \emph{et al.} (2017), por exemplo,
encontramos um exemplo de trabalho nesse sentido, que explora o uso de
redes neurais para a identificação de parametros relevantes à uma
simulação por digital waveguides do trato vocal humano.

Relaxar a simplificação adotada durante o trabalho em relação à
decaimentos exponenciais é um outro desenvolvimento futuro com potencial
interessante: para algumas categorias de som, como a voz humana por
exemplo, o envelope do som apresenta importância maior do que as
próprias frequências contidas em relação à características como
inteligibilidade.

Estimar os envelopes com a técnica utilizada no segundo modelo, a partir
de uma número maior de pontos, e utilizar uma rede, possivelmente
recorrente, para aprender as características desse envelope para um
conjunto de sons de determinado instrumento (inclusive de excitação
continuada), ou mesmo a voz humana, constitui outra interessante direção
a ser investigada.

Além disso, um algoritmo eficiente de extração de envelope permitiria
investigar a abordagem baseada em \emph{autoencoders} que, alimentados
com um senoide pura modelada pelo envelope da onda, teriam maior
possibilidade de aprender a reconstruir a onda original de maneira
verossímil, em uma abordagem inspirada nas técnicas de transferência de
estilo utilizadas no campo da visão de computadores.

Um outro desenvolvimento seria produzir uma implementação em uma
linguagem de programação mais eficiente, como C ou C++, somada à uma
interface visual e compatibilidade com controladores MIDI pode gerar uma
linha de produtos comercialmente viável, a ser comercializado em formato
standalone e/ou em formato de \emph{plugin} VST, em que um módulo
inicial pode ser oferecido e bancos adicionais, consistindo de novas
arquiteturas e pesos, podem ser criados e comercializdos mediante a
demanda dos usuários ou o desenvolvimento da técnica.

\hypertarget{referuxeancias}{%
\section*{Referências}\label{referuxeancias}}
\addcontentsline{toc}{section}{Referências}

\hypertarget{refs}{}
\leavevmode\hypertarget{ref-aljumah2016novel}{}%
ALJUMAH, A.; AHAMAD, T. A novel approach for detecting DDoS using
artificial neural networks. \textbf{International Journal of Computer
Science and Network Security}, v. 16, n. 12, p. 132--138, 2016.

\leavevmode\hypertarget{ref-bahrampour2015comparative}{}%
BAHRAMPOUR, S. \emph{et al.} Comparative study of deep learning software
frameworks. \textbf{arXiv preprint arXiv:1511.06435}, 2015.

\leavevmode\hypertarget{ref-bahrampour2016comparative}{}%
\_\_\_.. Comparative study of caffe, neon, theano, and torch for deep
learning. 2016.

\leavevmode\hypertarget{ref-bensa2005computational}{}%
BENSA, J. \emph{et al.} Computational modeling of stiff piano strings
using digital waveguides and finite differences. \textbf{Acta Acustica
united with Acustica}, v. 91, n. 2, p. 289--298, 2005.

\leavevmode\hypertarget{ref-bensa2003simulation}{}%
\_\_\_.. The simulation of piano string vibration: From physical models
to finite difference schemes and digital waveguides. \textbf{The Journal
of the Acoustical Society of America}, v. 114, n. 2, p. 1095--1107,
2003.

\leavevmode\hypertarget{ref-bair}{}%
\textbf{Berkeley Artificial Intelligence Research Lab}. \textbf{The
Berkeley Artificial Intelligence Research Blog}, {[}s.d.{]}. Disponível
em: \textless{}\url{http://bair.berkeley.edu/}\textgreater{}

\leavevmode\hypertarget{ref-bilbao2009numerical}{}%
BILBAO, S. \textbf{Numerical sound synthesis: finite difference schemes
and simulation in musical acoustics}. Traducao. {[}s.l.{]} John Wiley \&
Sons, 2009.

\leavevmode\hypertarget{ref-bishop06}{}%
BISHOP, C. M. \textbf{Pattern recognition and machine learning}.
Traducao. {[}s.l.{]} springer, 2006.

\leavevmode\hypertarget{ref-boulanger2013high}{}%
BOULANGER-LEWANDOWSKI, N.; BENGIO, Y.; VINCENT, P.
\textbf{High-dimensional sequence transduction}Acoustics, Speech and
Signal Processing (ICASSP), 2013 IEEE International Conference on.
\textbf{Anais}\ldots{}IEEE, 2013

\leavevmode\hypertarget{ref-boulanger2014phone}{}%
BOULANGER-LEWANDOWSKI, N. \emph{et al.} \textbf{Phone sequence modeling
with recurrent neural networks}Acoustics, Speech and Signal Processing
(ICASSP), 2014 IEEE International Conference on.
\textbf{Anais}\ldots{}IEEE, 2014

\leavevmode\hypertarget{ref-bovermann2016musical}{}%
BOVERMANN, T. \emph{et al.} \textbf{Musical Instruments in the 21st
Century}. Traducao. {[}s.l.{]} Springer, 2016.

\leavevmode\hypertarget{ref-bock2012polyphonic}{}%
BÖCK, S.; SCHEDL, M. \textbf{Polyphonic piano note transcription with
recurrent neural networks}Acoustics, speech and signal processing
(ICASSP), 2012 ieee international conference on.
\textbf{Anais}\ldots{}IEEE, 2012

\leavevmode\hypertarget{ref-brunette2009review}{}%
BRUNETTE, E. S.; FLEMMER, R. C.; FLEMMER, C. L. \textbf{A review of
artificial intelligence}Autonomous Robots and Agents, 2009. ICARA 2009.
4th International Conference on. \textbf{Anais}\ldots{}IEEE, 2009

\leavevmode\hypertarget{ref-caffe}{}%
\textbf{Caffe}. \textbf{Caffe \textbar{} Model Zoo}, {[}s.d.{]}.
Disponível em:
\textless{}\url{http://caffe.berkeleyvision.org/}\textgreater{}

\leavevmode\hypertarget{ref-chaigne2016acoustics}{}%
CHAIGNE, A.; KERGOMARD, J. \textbf{Acoustics of musical instruments}.
Traducao. {[}s.l.{]} Springer, 2016.

\leavevmode\hypertarget{ref-choi2016automatic}{}%
CHOI, K.; FAZEKAS, G.; SANDLER, M. Automatic tagging using deep
convolutional neural networks. \textbf{arXiv preprint arXiv:1606.00298},
2016.

\leavevmode\hypertarget{ref-choi2017convolutional}{}%
CHOI, K. \emph{et al.} \textbf{Convolutional recurrent neural networks
for music classification}Acoustics, Speech and Signal Processing
(ICASSP), 2017 IEEE International Conference on.
\textbf{Anais}\ldots{}IEEE, 2017

\leavevmode\hypertarget{ref-coppin2004artificial}{}%
COPPIN, B. \textbf{Artificial intelligence illuminated}. Traducao.
{[}s.l.{]} Jones; Bartlett Publishers, 2004.

\leavevmode\hypertarget{ref-costa2017evaluation}{}%
COSTA, Y. M.; OLIVEIRA, L. S.; SILLA, C. N. An evaluation of
Convolutional Neural Networks for music classification using
spectrograms. \textbf{Applied Soft Computing}, v. 52, p. 28--38, 2017.

\leavevmode\hypertarget{ref-crabacus_2015}{}%
CRABACUS. \textbf{crabacus/the-open-source-drumkit}\textbf{GitHub}, ago.
2015. Disponível em:
\textless{}\url{https://github.com/crabacus/the-open-source-drumkit}\textgreater{}

\leavevmode\hypertarget{ref-dalgleish2014blurring}{}%
DALGLEISH, M.; SPENCER, S.; FOSTER, C. \textbf{BLURRING THE LINES: AN
INTEGRATED COMPOSITIONAL MODEL FOR DIGITAL MUSICAL INSTRUMENT
DESIGN}\textbf{Anais}\ldots{}Proceedings of the 9th Conference on
Interdisciplinary Musicology. CIM14. Berlin, Germany, 2014

\leavevmode\hypertarget{ref-donahue2018synthesizing}{}%
DONAHUE, C.; MCAULEY, J.; PUCKETTE, M. Synthesizing Audio with
Generative Adversarial Networks. \textbf{arXiv preprint
arXiv:1802.04208}, 2018.

\leavevmode\hypertarget{ref-dozat2016incorporating}{}%
DOZAT, T. Incorporating nesterov momentum into adam. 2016.

\leavevmode\hypertarget{ref-drumgizmo}{}%
\textbf{DrumGizmo Wiki}. \textbf{DrumGizmo Wiki}, {[}s.d.{]}. Disponível
em:
\textless{}\url{https://www.drumgizmo.org/wiki/doku.php}\textgreater{}

\leavevmode\hypertarget{ref-duchi2011adaptive}{}%
DUCHI, J.; HAZAN, E.; SINGER, Y. Adaptive subgradient methods for online
learning and stochastic optimization. \textbf{Journal of Machine
Learning Research}, v. 12, n. Jul, p. 2121--2159, 2011.

\leavevmode\hypertarget{ref-dumoulin2016guide}{}%
DUMOULIN, V.; VISIN, F. A guide to convolution arithmetic for deep
learning. \textbf{ArXiv e-prints}, mar. 2016.

\leavevmode\hypertarget{ref-eck_2018}{}%
ECK, D. \textbf{Making music using new sounds generated with machine
learning}\textbf{Google}Google,, mar. 2018. Disponível em:
\textless{}\url{https://www.blog.google/technology/ai/making-music-using-new-sounds-generated-machine-learning/}\textgreater{}

\leavevmode\hypertarget{ref-elman1990finding}{}%
ELMAN, J. L. Finding structure in time. \textbf{Cognitive science}, v.
14, n. 2, p. 179--211, 1990.

\leavevmode\hypertarget{ref-engel2017neural}{}%
ENGEL, J. \emph{et al.} Neural audio synthesis of musical notes with
wavenet autoencoders. \textbf{arXiv preprint arXiv:1704.01279}, 2017.

\leavevmode\hypertarget{ref-erickson2017toolkits}{}%
ERICKSON, B. J. \emph{et al.} Toolkits and libraries for deep learning.
\textbf{Journal of digital imaging}, v. 30, n. 4, p. 400--405, 2017.

\leavevmode\hypertarget{ref-esteva2017dermatologist}{}%
ESTEVA, A. \emph{et al.} Dermatologist-level classification of skin
cancer with deep neural networks. \textbf{Nature}, v. 542, n. 7639, p.
115--118, 2017.

\leavevmode\hypertarget{ref-fontana2000using}{}%
FONTANA, F.; ROCCHESSO, D.; APOLLONIO, E. \textbf{Using the waveguide
mesh in modelling 3D resonators}Proc. Conf. on Digital Audio Effects
(DAFX-00). \textbf{Anais}\ldots{}2000

\leavevmode\hypertarget{ref-forgeard2008practicing}{}%
FORGEARD, M. \emph{et al.} Practicing a musical instrument in childhood
is associated with enhanced verbal ability and nonverbal reasoning.
\textbf{PloS one}, v. 3, n. 10, p. e3566, 2008.

\leavevmode\hypertarget{ref-frans2017outline}{}%
FRANS, K. Outline Colorization through Tandem Adversarial Networks.
\textbf{arXiv preprint arXiv:1704.08834}, 2017.

\leavevmode\hypertarget{ref-gatys2016image}{}%
GATYS, L. A.; ECKER, A. S.; BETHGE, M. \textbf{Image style transfer
using convolutional neural networks}Proceedings of the IEEE Conference
on Computer Vision and Pattern Recognition. \textbf{Anais}\ldots{}2016

\leavevmode\hypertarget{ref-gazi2018understanding}{}%
GAZI, O. \textbf{Understanding digital signal processing}. Traducao.
{[}s.l.{]} Springer, 2018.

\leavevmode\hypertarget{ref-goldberg2016primer}{}%
GOLDBERG, Y. A Primer on Neural Network Models for Natural Language
Processing. \textbf{J. Artif. Intell. Res.(JAIR)}, v. 57, p. 345--420,
2016.

\leavevmode\hypertarget{ref-goodfellow2016deep}{}%
GOODFELLOW, I. \emph{et al.} \textbf{Deep learning}. Traducao.
{[}s.l.{]} MIT press Cambridge, 2016. v. 1

\leavevmode\hypertarget{ref-gracia2016wave}{}%
GRACIA, X.; SANZ-PERELA, T. The wave equation for stiff strings and
piano tuning. \textbf{arXiv preprint arXiv:1603.05516}, 2016.

\leavevmode\hypertarget{ref-graves2013speech}{}%
GRAVES, A.; MOHAMED, A.-R.; HINTON, G. \textbf{Speech recognition with
deep recurrent neural networks}Acoustics, speech and signal processing
(icassp), 2013 ieee international conference on.
\textbf{Anais}\ldots{}IEEE, 2013

\leavevmode\hypertarget{ref-grinstein2017audio}{}%
GRINSTEIN, E. \emph{et al.} Audio style transfer. \textbf{arXiv preprint
arXiv:1710.11385}, 2017.

\leavevmode\hypertarget{ref-gully2017articulatory}{}%
GULLY, A. J. \emph{et al.} Articulatory Text-to-Speech Synthesis using
the Digital Waveguide Mesh driven by a Deep Neural Network.
\textbf{Proc. Interspeech 2017}, p. 234--238, 2017.

\leavevmode\hypertarget{ref-hagan1996neural}{}%
HAGAN, M. T. \emph{et al.} \textbf{Neural network design}. Traducao.
{[}s.l.{]} Pws Pub. Boston, 1996. v. 20

\leavevmode\hypertarget{ref-he2015multimodal}{}%
HE, L. \emph{et al.} \textbf{Multimodal affective dimension prediction
using deep bidirectional long short-term memory recurrent neural
networks}Proceedings of the 5th International Workshop on Audio/Visual
Emotion Challenge. \textbf{Anais}\ldots{}ACM, 2015

\leavevmode\hypertarget{ref-hinton2012deep}{}%
HINTON, G. \emph{et al.} Deep neural networks for acoustic modeling in
speech recognition: The shared views of four research groups.
\textbf{IEEE Signal Processing Magazine}, v. 29, n. 6, p. 82--97, 2012.

\leavevmode\hypertarget{ref-hinton2018matrix}{}%
HINTON, G. E.; SABOUR, S.; FROSST, N. Matrix capsules with EM routing.
2018.

\leavevmode\hypertarget{ref-hornik91}{}%
HORNIK, K. Approximation capabilities of multilayer feedforward
networks. \textbf{Neural networks}, v. 4, n. 2, p. 251--257, 1991.

\leavevmode\hypertarget{ref-hornik89}{}%
HORNIK, K.; STINCHCOMBE, M.; WHITE, H. Multilayer feedforward networks
are universal approximators. \textbf{Neural networks}, v. 2, n. 5, p.
359--366, 1989.

\leavevmode\hypertarget{ref-howard2017acoustics}{}%
HOWARD, D. M.; ANGUS, J. \textbf{Acoustics and psychoacoustics}.
Traducao. {[}s.l.{]} Focal press, 2017.

\leavevmode\hypertarget{ref-hutchings2017talking}{}%
HUTCHINGS, P. Talking Drums: Generating drum grooves with neural
networks. \textbf{arXiv preprint arXiv:1706.09558}, 2017.

\leavevmode\hypertarget{ref-hwangimage}{}%
HWANG, J.; ZHOU, Y. Image Colorization with Deep Convolutional Neural
Networks. 2016.

\leavevmode\hypertarget{ref-iizuka2016let}{}%
IIZUKA, S.; SIMO-SERRA, E.; ISHIKAWA, H. Let there be color!: joint
end-to-end learning of global and local image priors for automatic image
colorization with simultaneous classification. \textbf{ACM Transactions
on Graphics (TOG)}, v. 35, n. 4, p. 110, 2016.

\leavevmode\hypertarget{ref-isola2016image}{}%
ISOLA, P. \emph{et al.} Image-to-image translation with conditional
adversarial networks. \textbf{arXiv preprint arXiv:1611.07004}, 2016.

\leavevmode\hypertarget{ref-ivy_audio}{}%
\textbf{Ivy Audio}. \textbf{Ivy Audio}, {[}s.d.{]}. Disponível em:
\textless{}\url{http://www.ivyaudio.com/}\textgreater{}

\leavevmode\hypertarget{ref-karjalainen2004digital}{}%
KARJALAINEN, M.; ERKUT, C. Digital waveguides versus finite difference
structures: Equivalence and mixed modeling. \textbf{EURASIP Journal on
Applied Signal Processing}, v. 2004, p. 978--989, 2004.

\leavevmode\hypertarget{ref-karpathy2014large}{}%
KARPATHY, A. \emph{et al.} \textbf{Large-scale video classification with
convolutional neural networks}Proceedings of the IEEE conference on
Computer Vision and Pattern Recognition. \textbf{Anais}\ldots{}2014

\leavevmode\hypertarget{ref-keras}{}%
\textbf{Keras: The Python Deep Learning library}. \textbf{Keras
Documentation}, {[}s.d.{]}. Disponível em:
\textless{}\url{https://keras.io/}\textgreater{}

\leavevmode\hypertarget{ref-khorrami2016deep}{}%
KHORRAMI, P. \emph{et al.} \textbf{How deep neural networks can improve
emotion recognition on video data}Image Processing (ICIP), 2016 IEEE
International Conference on. \textbf{Anais}\ldots{}IEEE, 2016

\leavevmode\hypertarget{ref-kingma2014adam}{}%
KINGMA, D. P.; BA, J. Adam: A method for stochastic optimization.
\textbf{arXiv preprint arXiv:1412.6980}, 2014.

\leavevmode\hypertarget{ref-klambauer2017self}{}%
KLAMBAUER, G. \emph{et al.} \textbf{Self-normalizing neural
networks}Advances in Neural Information Processing Systems.
\textbf{Anais}\ldots{}2017

\leavevmode\hypertarget{ref-kulkarni2015deep}{}%
KULKARNI, T. D. \emph{et al.} \textbf{Deep convolutional inverse
graphics network}Advances in Neural Information Processing Systems.
\textbf{Anais}\ldots{}2015

\leavevmode\hypertarget{ref-laird2001physical}{}%
LAIRD, J. A. \textbf{The physical modelling of drums using digital
waveguides}. {[}s.l.{]} University of Bristol, 2001.

\leavevmode\hypertarget{ref-larsson2016learning}{}%
LARSSON, G.; MAIRE, M.; SHAKHNAROVICH, G. \textbf{Learning
representations for automatic colorization}European Conference on
Computer Vision. \textbf{Anais}\ldots{}Springer, 2016

\leavevmode\hypertarget{ref-lecun2012efficient}{}%
LECUN, Y. A. \emph{et al.} Efficient backprop. \emph{In}: \textbf{Neural
networks: Tricks of the trade}. Traducao. {[}s.l.{]} Springer, 2012. p.
9--48.

\leavevmode\hypertarget{ref-lecun1998gradient}{}%
LECUN, Y. \emph{et al.} Gradient-based learning applied to document
recognition. \textbf{Proceedings of the IEEE}, v. 86, n. 11, p.
2278--2324, 1998.

\leavevmode\hypertarget{ref-leshno93}{}%
LESHNO, M. \emph{et al.} Multilayer feedforward networks with a
nonpolynomial activation function can approximate any function.
\textbf{Neural networks}, v. 6, n. 6, p. 861--867, 1993.

\leavevmode\hypertarget{ref-lyons2011understanding}{}%
LYONS, R. G. \textbf{Understanding Digital Signal Processing, 3/E}.
Traducao. {[}s.l.{]} Pearson Education India, 2011.

\leavevmode\hypertarget{ref-maas2013rectifier}{}%
MAAS, A. L.; HANNUN, A. Y.; NG, A. Y. \textbf{Rectifier nonlinearities
improve neural network acoustic models}Proc. ICML.
\textbf{Anais}\ldots{}2013

\leavevmode\hypertarget{ref-magenta}{}%
\textbf{Magenta}. \textbf{Magenta}, {[}s.d.{]}. Disponível em:
\textless{}\url{https://magenta.tensorflow.org/}\textgreater{}

\leavevmode\hypertarget{ref-mcculloch1943logical}{}%
MCCULLOCH, W. S.; PITTS, W. A logical calculus of the ideas immanent in
nervous activity. \textbf{The bulletin of mathematical biophysics}, v.
5, n. 4, p. 115--133, 1943.

\leavevmode\hypertarget{ref-minski1969perceptrons}{}%
MINSKI, M. L.; PAPERT, S. A. Perceptrons: an introduction to
computational geometry. \textbf{MA: MIT Press, Cambridge}, 1969.

\leavevmode\hypertarget{ref-mital2017time}{}%
MITAL, P. K. Time Domain Neural Audio Style Transfer. \textbf{arXiv
preprint arXiv:1711.11160}, 2017.

\leavevmode\hypertarget{ref-mizutani2000derivation}{}%
MIZUTANI, E.; DREYFUS, S. E.; NISHIO, K. \textbf{On derivation of MLP
backpropagation from the Kelley-Bryson optimal-control gradient formula
and its application}Neural Networks, 2000. IJCNN 2000, Proceedings of
the IEEE-INNS-ENNS International Joint Conference on.
\textbf{Anais}\ldots{}IEEE, 2000

\leavevmode\hypertarget{ref-mullen2006physical}{}%
MULLEN, J. \textbf{Physical modelling of the vocal tract with the 2D
digital waveguide mesh}. {[}s.l.{]} University of York, 2006.

\leavevmode\hypertarget{ref-iowa}{}%
\textbf{Musical Instrument Samples}. \textbf{University of Iowa
Electronic Music Studios}, {[}s.d.{]}. Disponível em:
\textless{}\url{http://theremin.music.uiowa.edu/MIS.html}\textgreater{}

\leavevmode\hypertarget{ref-neon}{}%
\textbf{Neon}. \textbf{Intel AI}, {[}s.d.{]}. Disponível em:
\textless{}\url{https://ai.intel.com/neon/}\textgreater{}

\leavevmode\hypertarget{ref-google}{}%
\textbf{NSynthSuper}. \textbf{Google}Google,, {[}s.d.{]}. Disponível em:
\textless{}\url{https://nsynthsuper.withgoogle.com/}\textgreater{}

\leavevmode\hypertarget{ref-olson2017applied}{}%
OLSON, T.; LEVITT. \textbf{Applied Fourier Analysis}. Traducao.
{[}s.l.{]} Springer, 2017.

\leavevmode\hypertarget{ref-omar2017predicting}{}%
OMAR, N.; JOHARI, Z. `.; SMITH, M. Predicting fraudulent financial
reporting using artificial neural network. \textbf{Journal of Financial
Crime}, v. 24, n. 2, p. 362--387, 2017.

\leavevmode\hypertarget{ref-oord2016pixel}{}%
OORD, A. VAN DEN; KALCHBRENNER, N.; KAVUKCUOGLU, K. Pixel recurrent
neural networks. \textbf{arXiv preprint arXiv:1601.06759}, 2016.

\leavevmode\hypertarget{ref-van2016conditional}{}%
OORD, A. VAN DEN \emph{et al.} \textbf{Conditional image generation with
pixelcnn decoders}Advances in Neural Information Processing Systems.
\textbf{Anais}\ldots{}2016

\leavevmode\hypertarget{ref-pang2017convolution}{}%
PANG, Y. \emph{et al.} Convolution in convolution for network in
network. \textbf{IEEE Transactions on Neural Networks and Learning
Systems}, 2017.

\leavevmode\hypertarget{ref-parker_2018}{}%
PARKER, C. \textbf{Artificial intelligence could be our saviour,
according to the CEO of Google}\textbf{World Economic Forum}, jan. 2018.
Disponível em:
\textless{}\url{https://www.weforum.org/agenda/2018/01/google-ceo-ai-will-be-bigger-than-electricity-or-fire/}\textgreater{}

\leavevmode\hypertarget{ref-parker1985learning}{}%
PARKER, D. B. Learning logic. 1985.

\leavevmode\hypertarget{ref-parvat2017survey}{}%
PARVAT, A. \emph{et al.} \textbf{A survey of deep-learning
frameworks}Inventive Systems and Control (ICISC), 2017 International
Conference on. \textbf{Anais}\ldots{}IEEE, 2017

\leavevmode\hypertarget{ref-pereyra2012harmonic}{}%
PEREYRA, M. C.; WARD, L. A. \textbf{Harmonic analysis: from Fourier to
wavelets}. Traducao. {[}s.l.{]} American Mathematical Soc., 2012. v. 63

\leavevmode\hypertarget{ref-pfalz2018generating}{}%
PFALZ, A. Generating Audio Using Recurrent Neural Networks. 2018.

\leavevmode\hypertarget{ref-raschka2015python}{}%
RASCHKA, S. \textbf{Python machine learning}. Traducao. {[}s.l.{]} Packt
Publishing Ltd, 2015.

\leavevmode\hypertarget{ref-reddi2018convergence}{}%
REDDI, S. J.; KALE, S.; KUMAR, S. On the convergence of adam and beyond.
2018.

\leavevmode\hypertarget{ref-reiffenstein2006codification}{}%
REIFFENSTEIN, T. Codification, patents and the geography of knowledge
transfer in the electronic musical instrument industry. \textbf{The
Canadian Geographer/Le Géographe canadien}, v. 50, n. 3, p. 298--318,
2006.

\leavevmode\hypertarget{ref-rigaud2013parametric}{}%
RIGAUD, F.; DAVID, B.; DAUDET, L. A parametric model and estimation
techniques for the inharmonicity and tuning of the piano. \textbf{The
Journal of the Acoustical Society of America}, v. 133, n. 5, p.
3107--3118, 2013.

\leavevmode\hypertarget{ref-risi2017neuroevolution}{}%
RISI, S.; TOGELIUS, J. Neuroevolution in games: State of the art and
open challenges. \textbf{IEEE Transactions on Computational Intelligence
and AI in Games}, v. 9, n. 1, p. 25--41, 2017.

\leavevmode\hypertarget{ref-roberts2018learning}{}%
ROBERTS, A. \emph{et al.} Learning Latent Representations of Music to
Generate Interactive Musical Palettes. 2018.

\leavevmode\hypertarget{ref-roberts2018hierarchical}{}%
\_\_\_.. A Hierarchical Latent Vector Model for Learning Long-Term
Structure in Music. \textbf{arXiv preprint arXiv:1803.05428}, 2018.

\leavevmode\hypertarget{ref-Rosenblatt57}{}%
ROSENBLATT, F. \textbf{The perceptron, a perceiving and recognizing
automaton Project Para}. Traducao. {[}s.l.{]} Cornell Aeronautical
Laboratory, 1957.

\leavevmode\hypertarget{ref-rosenblatt1958perceptron}{}%
\_\_\_. The perceptron: a probabilistic model for information storage
and organization in the brain. \textbf{Psychological review}, v. 65, n.
6, p. 386, 1958.

\leavevmode\hypertarget{ref-ruder2016overview}{}%
RUDER, S. An overview of gradient descent optimization algorithms.
\textbf{arXiv preprint arXiv:1609.04747}, 2016.

\leavevmode\hypertarget{ref-rumelhart1985learning}{}%
RUMELHART, D. E.; HINTON, G. E.; WILLIAMS, R. J. \textbf{Learning
internal representations by error propagation}. {[}s.l.{]} California
Univ San Diego La Jolla Inst for Cognitive Science, 1985.

\leavevmode\hypertarget{ref-russell2016artificial}{}%
RUSSELL, S. J.; NORVIG, P. \textbf{Artificial intelligence: a modern
approach}. Traducao. {[}s.l.{]} Malaysia; Pearson Education Limited,
2016.

\leavevmode\hypertarget{ref-sabour2017dynamic}{}%
SABOUR, S.; FROSST, N.; HINTON, G. E. \textbf{Dynamic routing between
capsules}Advances in Neural Information Processing Systems.
\textbf{Anais}\ldots{}2017

\leavevmode\hypertarget{ref-sainath2015deep}{}%
SAINATH, T. N. \emph{et al.} Deep convolutional neural networks for
large-scale speech tasks. \textbf{Neural Networks}, v. 64, p. 39--48,
2015.

\leavevmode\hypertarget{ref-sak2015fast}{}%
SAK, H. \emph{et al.} Fast and accurate recurrent neural network
acoustic models for speech recognition. \textbf{arXiv preprint
arXiv:1507.06947}, 2015.

\leavevmode\hypertarget{ref-salsa2016partial}{}%
SALSA, S. \textbf{Partial differential equations in action: from
modelling to theory}. Traducao. {[}s.l.{]} Springer, 2016. v. 99

\leavevmode\hypertarget{ref-sangkloy2016scribbler}{}%
SANGKLOY, P. \emph{et al.} Scribbler: Controlling deep image synthesis
with sketch and color. \textbf{arXiv preprint arXiv:1612.00835}, 2016.

\leavevmode\hypertarget{ref-sarroff2014musical}{}%
SARROFF, A. M.; CASEY, M. A. \textbf{Musical audio synthesis using
autoencoding neural nets}ICMC. \textbf{Anais}\ldots{}2014

\leavevmode\hypertarget{ref-schmidhuber2015deep}{}%
SCHMIDHUBER, J. Deep learning in neural networks: An overview.
\textbf{Neural networks}, v. 61, p. 85--117, 2015.

\leavevmode\hypertarget{ref-torch}{}%
\textbf{Scientific computing for LuaJIT.} \textbf{Torch}, {[}s.d.{]}.
Disponível em: \textless{}\url{http://torch.ch/}\textgreater{}

\leavevmode\hypertarget{ref-serafin2001banded}{}%
SERAFIN, S.; HUANG, P.; SMITH, J. \textbf{The banded digital waveguide
mesh}Proc. Workshop on Future Directions of Computer Music (Mosart-01),
Barcelona, Spain. \textbf{Anais}\ldots{}Citeseer, 2001

\leavevmode\hypertarget{ref-serra2007state}{}%
SERRA, X. \textbf{State of the art and future directions in musical
sound synthesis}Multimedia Signal Processing, 2007. MMSP 2007. IEEE 9th
Workshop on. \textbf{Anais}\ldots{}IEEE, 2007

\leavevmode\hypertarget{ref-shead_2016}{}%
SHEAD, S. \textbf{Microsoft exec: 'AI is the most important technology
that anybody on the planet is working on today'}\textbf{Business Insider
Deutschland}BI,, maio 2016. Disponível em:
\textless{}\url{https://www.businessinsider.de/microsoft-exec-ai-is-the-most-important-technology-that-anybody-on-the-planet-is-working-on-today-2016-5?r=UK\&IR=T}\textgreater{}

\leavevmode\hypertarget{ref-shi2016benchmarking}{}%
SHI, S. \emph{et al.} \textbf{Benchmarking state-of-the-art deep
learning software tools}Cloud Computing and Big Data (CCBD), 2016 7th
International Conference on. \textbf{Anais}\ldots{}IEEE, 2016

\leavevmode\hypertarget{ref-sigtia2016end}{}%
SIGTIA, S.; BENETOS, E.; DIXON, S. An end-to-end neural network for
polyphonic piano music transcription. \textbf{IEEE/ACM Transactions on
Audio, Speech and Language Processing (TASLP)}, v. 24, n. 5, p.
927--939, 2016.

\leavevmode\hypertarget{ref-smith2006basic}{}%
SMITH, J. A basic introduction to digital waveguide synthesis (for the
technically inclined). \textbf{Center for Computer Research in Music and
Acoustics (CCRMA), Stanford University. http://ccrma. stanford.
edu/\textasciitilde{} jos/swgt}, 2006.

\leavevmode\hypertarget{ref-smith1992physical}{}%
SMITH, J. O. Physical modeling using digital waveguides.
\textbf{Computer music journal}, v. 16, n. 4, p. 74--91, 1992.

\leavevmode\hypertarget{ref-smith2008digital}{}%
\_\_\_. Digital waveguide architectures for virtual musical instruments.
\emph{In}: \textbf{Handbook of Signal Processing in Acoustics}.
Traducao. {[}s.l.{]} Springer, 2008. p. 399--417.

\leavevmode\hypertarget{ref-smith1997scientist}{}%
SMITH, S. W.; OTHERS. The scientist and engineer's guide to digital
signal processing. 1997.

\leavevmode\hypertarget{ref-socher2014recursive}{}%
SOCHER, R. Recursive deep learning for natural language processing and
computer vision. 2014.

\leavevmode\hypertarget{ref-sorensen1987real}{}%
SORENSEN, H. V. \emph{et al.} Real-valued fast Fourier transform
algorithms. \textbf{IEEE Transactions on Acoustics, Speech, and Signal
Processing}, v. 35, n. 6, p. 849--863, 1987.

\leavevmode\hypertarget{ref-philharmoniaorchestra}{}%
\textbf{Sound Samples}. \textbf{Philharmonia Orchestra}, {[}s.d.{]}.
Disponível em:
\textless{}\url{http://www.philharmonia.co.uk/explore/sound_samples}\textgreater{}

\leavevmode\hypertarget{ref-southall2016automatic}{}%
SOUTHALL, C.; STABLES, R.; HOCKMAN, J. \textbf{Automatic Drum
Transcription Using Bi-Directional Recurrent Neural Networks.}ISMIR.
\textbf{Anais}\ldots{}2016

\leavevmode\hypertarget{ref-stanley2007compositional}{}%
STANLEY, K. O. Compositional pattern producing networks: A novel
abstraction of development. \textbf{Genetic programming and evolvable
machines}, v. 8, n. 2, p. 131--162, 2007.

\leavevmode\hypertarget{ref-staudt2016development}{}%
STAUDT, P. Development of a Digital Musical Instrument with Embedded
Sound Synthesis. 2016.

\leavevmode\hypertarget{ref-stein2003princeton}{}%
STEIN, E. M.; SHAKARCHI, R. \textbf{Princeton lectures in analysis}.
Traducao. {[}s.l.{]} Princeton University Press, 2003.

\leavevmode\hypertarget{ref-sutskever2013importance}{}%
SUTSKEVER, I. \emph{et al.} \textbf{On the importance of initialization
and momentum in deep learning}International conference on machine
learning. \textbf{Anais}\ldots{}2013

\leavevmode\hypertarget{ref-tensorflow}{}%
\textbf{TensorFlow}. \textbf{TensorFlow}, {[}s.d.{]}. Disponível em:
\textless{}\url{https://www.tensorflow.org/}\textgreater{}

\leavevmode\hypertarget{ref-tesserato_2018}{}%
TESSERATO. \textbf{tesserato/tesserato.github.io}\textbf{GitHub}Carlos
Tarjano,, jul. 2018. Disponível em:
\textless{}\url{https://github.com/tesserato/tesserato.github.io}\textgreater{}

\leavevmode\hypertarget{ref-theano}{}%
\textbf{Theano}. \textbf{Multilayer Perceptron - DeepLearning 0.1
documentation}, {[}s.d.{]}. Disponível em:
\textless{}\url{http://deeplearning.net/software/theano/}\textgreater{}

\leavevmode\hypertarget{ref-theis2015generative}{}%
THEIS, L.; BETHGE, M. \textbf{Generative image modeling using spatial
LSTMs}Advances in Neural Information Processing Systems.
\textbf{Anais}\ldots{}2015

\leavevmode\hypertarget{ref-tieleman2012lecture}{}%
TIELEMAN, T.; HINTON, G. Lecture 6.5-rmsprop: Divide the gradient by a
running average of its recent magnitude. \textbf{COURSERA: Neural
networks for machine learning}, v. 4, n. 2, p. 26--31, 2012.

\leavevmode\hypertarget{ref-tuohy2006evolved}{}%
TUOHY, D. R.; POTTER, W. D.; CENTER, A. I. \textbf{An Evolved Neural
Network/HC Hybrid for Tablature Creation in GA-based Guitar
Arranging.}ICMC. \textbf{Anais}\ldots{}2006

\leavevmode\hypertarget{ref-van2016wavenet}{}%
VAN DEN OORD, A. \emph{et al.} Wavenet: A generative model for raw
audio. \textbf{arXiv preprint arXiv:1609.03499}, 2016.

\leavevmode\hypertarget{ref-van1993physical}{}%
VAN DUYNE, S. A.; SMITH, J. O. \textbf{Physical modeling with the 2-D
digital waveguide mesh}Proceedings of the International Computer Music
Conference. \textbf{Anais}\ldots{}INTERNATIONAL COMPUTER MUSIC
ACCOCIATION, 1993

\leavevmode\hypertarget{ref-van1995tetrahedral}{}%
\_\_\_. \textbf{The tetrahedral digital waveguide mesh}Applications of
Signal Processing to Audio and Acoustics, 1995., IEEE ASSP Workshop on.
\textbf{Anais}\ldots{}IEEE, 1995

\leavevmode\hypertarget{ref-vaughn2000music}{}%
VAUGHN, K. Music and mathematics: Modest support for the oft-claimed
relationship. \textbf{Journal of aesthetic education}, v. 34, n. 3/4, p.
149--166, 2000.

\leavevmode\hypertarget{ref-veen_2018}{}%
VEEN, F. VAN. \textbf{The Neural Network Zoo}\textbf{The Asimov
Institute}The Asimov Institute,, maio 2018. Disponível em:
\textless{}\url{http://www.asimovinstitute.org/neural-network-zoo/}\textgreater{}

\leavevmode\hypertarget{ref-veit2016residual}{}%
VEIT, A.; WILBER, M. J.; BELONGIE, S. \textbf{Residual networks behave
like ensembles of relatively shallow networks}Advances in Neural
Information Processing Systems. \textbf{Anais}\ldots{}2016

\leavevmode\hypertarget{ref-deepmind}{}%
\textbf{WaveNet: A Generative Model for Raw Audio}. \textbf{DeepMind},
{[}s.d.{]}. Disponível em:
\textless{}\url{https://deepmind.com/blog/wavenet-generative-model-raw-audio/}\textgreater{}

\leavevmode\hypertarget{ref-werbos1974beyond}{}%
WERBOS, P. Beyond Regression:" New Tools for Prediction and Analysis in
the Behavioral Sciences. \textbf{Ph. D. dissertation, Harvard
University}, 1974.

\leavevmode\hypertarget{ref-amazon}{}%
\textbf{What is Artificial Intelligence? - Amazon Web Services}.
\textbf{Amazon}Amazon,, {[}s.d.{]}. Disponível em:
\textless{}\url{https://aws.amazon.com/machine-learning/what-is-ai/}\textgreater{}

\leavevmode\hypertarget{ref-widrow1960adaptive}{}%
WIDROW, B.; HOFF, M. E. \textbf{Adaptive switching circuits}. {[}s.l.{]}
STANFORD UNIV CA STANFORD ELECTRONICS LABS, 1960.

\leavevmode\hypertarget{ref-widrow199030}{}%
WIDROW, B.; LEHR, M. A. 30 years of adaptive neural networks:
perceptron, madaline, and backpropagation. \textbf{Proceedings of the
IEEE}, v. 78, n. 9, p. 1415--1442, 1990.

\leavevmode\hypertarget{ref-xu2015ccg}{}%
XU, W.; AULI, M.; CLARK, S. \textbf{CCG Supertagging with a Recurrent
Neural Network.}ACL (2). \textbf{Anais}\ldots{}2015

\leavevmode\hypertarget{ref-yadav2015introduction}{}%
YADAV, N.; YADAV, A.; KUMAR, M. \textbf{An introduction to neural
network methods for differential equations}. Traducao. {[}s.l.{]}
Springer, 2015.

\leavevmode\hypertarget{ref-Yang2017}{}%
YANG, Y.; KROMPASS, D.; TRESP, V. Tensor-Train Recurrent Neural Networks
for Video Classification. 2017.

\leavevmode\hypertarget{ref-zeiler2012adadelta}{}%
ZEILER, M. D. ADADELTA: an adaptive learning rate method. \textbf{arXiv
preprint arXiv:1212.5701}, 2012.

\leavevmode\hypertarget{ref-zhang2016colorful}{}%
ZHANG, R.; ISOLA, P.; EFROS, A. A. \textbf{Colorful image
colorization}European Conference on Computer Vision.
\textbf{Anais}\ldots{}Springer, 2016

\leavevmode\hypertarget{ref-zhang2017very}{}%
ZHANG, Y.; CHAN, W.; JAITLY, N. \textbf{Very deep convolutional networks
for end-to-end speech recognition}Acoustics, Speech and Signal
Processing (ICASSP), 2017 IEEE International Conference on.
\textbf{Anais}\ldots{}IEEE, 2017

\leavevmode\hypertarget{ref-zhang2017towards}{}%
ZHANG, Y. \emph{et al.} Towards end-to-end speech recognition with deep
convolutional neural networks. \textbf{arXiv preprint arXiv:1701.02720},
2017.

\leavevmode\hypertarget{ref-zhu2016generative}{}%
ZHU, J.-Y. \emph{et al.} \textbf{Generative visual manipulation on the
natural image manifold}European Conference on Computer Vision.
\textbf{Anais}\ldots{}Springer, 2016

\leavevmode\hypertarget{ref-zweig2017advances}{}%
ZWEIG, G. \emph{et al.} \textbf{Advances in all-neural speech
recognition}Acoustics, Speech and Signal Processing (ICASSP), 2017 IEEE
International Conference on. \textbf{Anais}\ldots{}IEEE, 2017

\leavevmode\hypertarget{ref-zwicker2013psychoacoustics}{}%
ZWICKER, E.; FASTL, H. \textbf{Psychoacoustics: Facts and models}.
Traducao. {[}s.l.{]} Springer Science \& Business Media, 2013. v. 22

\end{document}
